% =============================================================================
% APPENDIX MP: CORE-MASTER: The Complete Optimization Problem
% =============================================================================
% Module: BCM2_MASTER (Framework Integration)
% Category: CORE
% Version: 1.0 (January 2026)
% Status: Draft
% Language: English
%
% PURPOSE: This appendix defines the COMPLETE optimization problem that
% underlies the entire EBF framework. It integrates contributions from
% ALL 24 chapters and ALL 10 CORE appendices into a unified formal structure.
%
% This is the "Appendix of Appendices" - the mathematical backbone that
% shows how every component fits together.
%
% =============================================================================

% -----------------------------------------------------------------------------
% HEADER BLOCK (Required)
% -----------------------------------------------------------------------------

\begin{tcolorbox}[colback=gray!5!white, colframe=gray!75!black,
    title=Appendix Header: MP CORE-MASTER]

\begin{tabular}{@{}ll@{}}
\textbf{Code:} & MP \\
\textbf{Category:} & CORE (Fundamental Theory) \\
\textbf{Full Name:} & CORE-MASTER: The Complete Optimization Problem \\
\textbf{Version:} & 1.0 (January 2026) \\
\textbf{Status:} & Draft \\
\textbf{Language:} & English \\
\end{tabular}

\vspace{0.5em}
\textbf{Purpose:} Defines the \textbf{complete optimization problem} that integrates
all 24 chapters and 10 CORE appendices. This is the formal mathematical backbone
of the EBF framework---the ``Appendix of Appendices.''

\textbf{Central Question:}
\begin{quote}
\textit{``How do we maximize the probability of sustainable behavioral change
over time, given the 10C state, context dynamics, and intervention constraints?''}
\end{quote}

\end{tcolorbox}

% -----------------------------------------------------------------------------
% ABSTRACT
% -----------------------------------------------------------------------------

\begin{abstract}
This appendix defines the \textbf{complete optimization problem} underlying the
Evidence-Based Framework for Economic and Social Behavior (EBF). It integrates
all 24 chapters and 10 CORE appendices into a unified formal structure, answering
the fundamental question: \textit{How do we maximize the probability of sustainable
behavioral change over time?} The problem is formulated as a dynamic optimization
with objective function $P_{eff}$ (probability of behavioral change), 10C state space,
continuous intervention instruments $I \in [0,1]^9$, endogenous context dynamics,
and constraints C1-C7. Key structural features include non-concavity (multiple equilibria),
the K vs Q distinction, path dependence, and multi-level coherence requirements.
\end{abstract}

% -----------------------------------------------------------------------------
% QUICK REFERENCE BOX
% -----------------------------------------------------------------------------

\begin{tcolorbox}[colback=blue!5!white, colframe=blue!75!black,
    title=\textbf{Quick Reference: MP CORE-MASTER}]
\small
\begin{tabular}{@{}p{3.5cm}p{9cm}@{}}
\textbf{Objective Function} & $P_{eff} = \sigma(WEC \cdot \alpha_{BCJ} \cdot \beta_{BCS})$ (Ch16) \\[3pt]
\textbf{State Space} & $S = (L, d, \gamma, \Psi, \Theta, A, W, \varphi, N_{L2})$ (10C) \\[3pt]
\textbf{Instruments} & $\vec{I} \in [0,1]^9$ continuous intervention vector (IE) \\[3pt]
\textbf{Constraints} & C1-C7 (Ch20) + Non-concavity (Ch7) \\[3pt]
\textbf{Dynamics} & $\Psi_{t+1} = f(\Psi_t, I_t, a_t)$ endogenous context (Ch9) \\[3pt]
\textbf{Key Insight} & Landscape has multiple equilibria; K $\neq$ Q \\
\end{tabular}
\end{tcolorbox}

% -----------------------------------------------------------------------------
% SCOPE BOX
% -----------------------------------------------------------------------------

\begin{tcolorbox}[colback=orange!5!white, colframe=orange!75!black,
    title=\textbf{Appendix Scope: Ziel / In-Scope / Out-of-Scope / Lieferobjekte}]

\textbf{Ziel (Objective):}
\begin{quote}
\textit{Bereitstellung des vollständigen mathematischen Optimierungsproblems, das alle 24 Kapitel
und 10 CORE Appendices unter einer einheitlichen Formulierung integriert---das ``Cockpit'' des EBF Frameworks.}
\end{quote}

\textbf{Zentrale Frage:} Wie maximieren wir $P_{eff}$ (Wahrscheinlichkeit nachhaltiger Verhaltensänderung)
über Zeit, gegeben 10C Status Quo, Kontext-Dynamik und Interventions-Constraints?

\vspace{0.5em}
\textbf{In-Scope:}
\begin{itemize}[nosep]
    \item \textbf{Master-Gleichung:} $\max \sum_t \delta^t P_{eff}(S_t, I_t, \Psi_t)$ mit allen Komponenten
    \item \textbf{10C $\to$ State Mapping:} Wie jede CORE-Dimension zum Zustandsvektor $S$ beiträgt
    \item \textbf{Constraint-System:} C1--C7 (Budget, Crowding, Timing, Capacity, Ethics, Segment, K*)
    \item \textbf{Dynamik:} Bellman-Formulierung, Kontext-Evolution $\Psi_{t+1} = f(\Psi_t, I_t, a_t)$
    \item \textbf{Strukturelle Eigenschaften:} Non-Konkavität, Multiple Equilibria, K vs Q Distinction
    \item \textbf{Kapitel-Integration:} 140 Axiome (MP-A1 bis MP-A140), 62 Theoreme (MP-T1 bis MP-T62)
    \item \textbf{Diagnostik-Referenz:} Schnelle Fehlersuche (Phase-Affinity, Segment-Multiplier, Crowding)
\end{itemize}

\vspace{0.5em}
\textbf{Out-of-Scope:}
\begin{itemize}[nosep]
    \item \textbf{Lösungsalgorithmen} $\rightarrow$ Appendix UNMAPPED_LLM (METHOD-LLMMC), Appendix EEE (METHOD-DESIGN)
    \item \textbf{Empirische Kalibrierung} $\rightarrow$ Appendix UNMAPPED_BBB (CORE-WHERE), METHOD appendices
    \item \textbf{Domain-spezifische Anwendungen} $\rightarrow$ DOMAIN appendices (AA--AG, AJ, AK, W--Z)
    \item \textbf{Interventions-Design Workflow} $\rightarrow$ Appendix TKT (METHOD-TOOLKIT), /design-intervention
    \item \textbf{Modell-Design Workflow} $\rightarrow$ Appendix EEE (METHOD-DESIGN), /design-model
    \item \textbf{Case Studies} $\rightarrow$ data/case-registry.yaml, /case Skill
\end{itemize}

\vspace{0.5em}
\textbf{Lieferobjekte (Deliverables):}
\begin{enumerate}[nosep]
    \item \textbf{Master-Problem (UNMAPPED_MP-1):} Vollständige Optimierungsgleichung mit Objective, Constraints, Dynamics
    \item \textbf{10C-State-Mapping-Tabelle:} Wie WHO/WHAT/HOW/.../HIERARCHY zu $S$ werden
    \item \textbf{Chapter-Contribution-Matrix:} Welches Kapitel welche Komponente liefert (Sec.~\ref{sec:mp-chapters})
    \item \textbf{Axiom-Katalog:} 140 nummerierte Axiome mit wissenschaftlichen Zitationen
    \item \textbf{Theorem-Katalog:} 62 nummerierte Theoreme mit Beweisverweisen
    \item \textbf{4-Stage Pipeline:} $U_{pot} \to U_{eff} \to WAX \geq \theta \to S(t)$
    \item \textbf{Bellman-Formulierung:} Dynamische Programmierung für EBF
    \item \textbf{Umfassendes Glossar:} 100+ Begriffe in 14 Kategorien (Sec.~\ref{sec:mp-glossary})
    \item \textbf{K*-Formel:} Emergente Portfolio-Größe mit allen Faktoren
    \item \textbf{Diagnostik-Checkliste:} Schnelle Fehlersuche bei Interventions-Versagen
\end{enumerate}

\vspace{0.5em}
\textbf{Voraussetzungen:}
\begin{itemize}[nosep]
    \item Grundverständnis der 10 CORE Appendices (AAA, B, C, V, BBB, AU, AV, AW, HI, IE)
    \item Kenntnis der Kapitel 1--24 (oder selektives Nachschlagen via Chapter-Matrix)
\end{itemize}

\end{tcolorbox}

% -----------------------------------------------------------------------------
% CROSS-REFERENCE STRUCTURE
% -----------------------------------------------------------------------------

\begin{tcolorbox}[colback=purple!5!white, colframe=purple!75!black,
    title=Chapter \& Appendix Integration]

\textbf{This appendix integrates ALL 24 chapters:}
\begin{itemize}[nosep]
\item \textbf{Foundations (Ch1-8):} Motivation, rationality, evidence, complementarity, reference, non-concavity, mathematics
\item \textbf{10C CORE (Ch9-15):} WHEN, WHAT/WHO, AWARE, READY, STAGE, SEGMENTS, HIERARCHY
\item \textbf{Output (Ch16):} The objective function $P_{eff}$
\item \textbf{Interventions (Ch17-21):} Instruments, phase-affinity, segments, portfolios, dynamics
\item \textbf{Scaling (Ch22-24):} Policy, multi-level, lifespan emergence
\end{itemize}

\textbf{And ALL 10 CORE appendices:}
\begin{center}
\small
\begin{tabular}{llll}
AAA (WHO) & C (WHAT) & B (HOW) & V (WHEN) \\
BBB (WHERE) & AU (AWARE) & AV (READY) & AW (STAGE) \\
HI (HIERARCHY) & IE (EIT) & & \\
\end{tabular}
\end{center}

\end{tcolorbox}

% -----------------------------------------------------------------------------
% METHODOLOGICAL FOUNDATION: EXCLUSION PRINCIPLE
% -----------------------------------------------------------------------------

\begin{tcolorbox}[colback=orange!5!white, colframe=orange!75!black,
    title=\textbf{Methodological Foundation: Exclusion Principle (EXC-1 to EXC-4)}]

\textbf{Reference:} Appendix UNMAPPED_FRM, Section \ref{sec:frm-exclusion-principle}

\textbf{Core Principle:} In EBF, \textbf{ADDITIVE aggregation is the default} ($H_0$). Multiplicative (complementary) formulations require explicit justification.

\begin{equation}
\boxed{
\underbrace{Y = Y_0 + \sum_j \delta_j \cdot f_j}_{\text{DEFAULT (EXC-1)}} \quad \text{vs.} \quad
\underbrace{Y = Y_0 \cdot \prod_j f_j}_{\text{requires Veto-Logic (EXC-2)}}
}
\end{equation}

\textbf{Veto-Logic Test (EXC-2):} Multiplicative is justified for factor $f_j$ iff: $f_j = 0 \Rightarrow Y = 0$.

\textbf{Documentation Requirement (EXC-4):} Every multiplicative formula in this appendix includes:
\begin{enumerate}[nosep]
    \item Which factors satisfy veto-logic?
    \item Empirical evidence for interaction
    \item Theoretical mechanism
\end{enumerate}

\textit{Validation: \texttt{python scripts/check\_formula\_compliance.py}}

\end{tcolorbox}

% =============================================================================
\section{The Central Question}
\label{sec:mp-central}
% =============================================================================

\begin{tcolorbox}[colback=red!5!white, colframe=red!75!black,
    title=\textbf{The Fundamental Question of EBF}]

\Large
\centering
\textit{How do we maximize the probability of sustainable behavioral change\\
over time, given the 10C state, context dynamics, and intervention constraints?}

\vspace{0.5em}
\normalsize
Formally: $\displaystyle\max_{\{I_t\}_{t=0}^T} \sum_{t=0}^{T} \delta^t \cdot P_{eff}(S_t, I_t, \Psi_t)$

\end{tcolorbox}

This question integrates:
\begin{itemize}
\item \textbf{What to maximize:} $P_{eff}$ --- the probability of behavioral change (Chapter 16)
\item \textbf{Over what horizon:} $t = 0, \ldots, T$ --- potentially infinite (Chapter 24)
\item \textbf{With what instruments:} $I_t \in [0,1]^9$ --- the intervention vector (Chapter 17, IE)
\item \textbf{Subject to what state:} $S_t$ --- the 10C status quo (Chapters 9-15)
\item \textbf{In what context:} $\Psi_t$ --- endogenously evolving (Chapter 9)
\end{itemize}

% =============================================================================
\section{The Complete Optimization Problem}
\label{sec:mp-problem}
% =============================================================================

\begin{tcolorbox}[colback=blue!5!white, colframe=blue!75!black,
    title=\textbf{The EBF Master Problem (UNMAPPED_MP-1)}]

\begin{equation}
\boxed{
\max_{\{I_t, P_t\}_{t=0}^T} \sum_{t=0}^{T} \delta^t \cdot P_{eff}(S_t, I_t, \Psi_t)
}
\label{eq:mp-objective}
\end{equation}

\textbf{Subject to:}

\vspace{0.3em}
\textit{State Dynamics:}
\begin{align}
S_{t+1} &= S_t + \Delta S(I_t) & \text{[State Transition, EIT-2]} \label{eq:mp-state} \\
\Psi_{t+1} &= f(\Psi_t, I_t, a_t) & \text{[Context Evolution, Ch9]} \label{eq:mp-context}
\end{align}

\textit{Feasibility Constraints:}
\begin{align}
I_t &\in \mathcal{S}(\Psi_t, S_t) & \text{[Solution Space, SS]} \label{eq:mp-feasible} \\
|P_t| &= K^*(B_t, c_t, \varphi_t, \gamma_t, L_t) & \text{[K-Optimal Portfolio, Ch20]} \label{eq:mp-kstar}
\end{align}

\textit{Hard Constraints (C1-C7 from Chapter 20):}
\begin{align}
\text{C1:} \quad & \sum_{i \in P_t} \text{cost}(I_i) \leq B_t & \text{[Budget]} \\
\text{C2:} \quad & \sum_{i < j} \gamma_{ij}(\Psi_t) \geq -0.3 & \text{[Crowding-Out]} \\
\text{C3:} \quad & A_{\min} \leq A_t \leq A_{\max} & \text{[Awareness Bounds]} \\
\text{C4:} \quad & WAX_t \geq \theta_t & \text{[Willingness Threshold]} \\
\text{C5:} \quad & \sigma_s(I) \geq 0 \text{ for } \geq 85\% \text{ pop.} & \text{[Segment Safety]} \\
\text{C6:} \quad & \alpha(\varphi_t, I_i) \geq \alpha_{\min} & \text{[Phase Affinity]} \\
\text{C7:} \quad & |P_t| = K^*_t & \text{[Portfolio Size]}
\end{align}

\textit{Structural Constraints (from Chapter 7):}
\begin{align}
\text{C8:} \quad & \text{Landscape is non-concave (multiple equilibria)} & \text{[Non-Concavity]} \\
\text{C9:} \quad & \text{Valley-crossing requires simultaneous change} & \text{[Coordination]} \\
\text{C10:} \quad & K \approx 1 \not\Rightarrow \max Q & \text{[K-Q Distinction]}
\end{align}

\end{tcolorbox}

% =============================================================================
\section{The Objective Function: $P_{eff}$}
\label{sec:mp-objective}
% =============================================================================

The objective function comes from \textbf{Chapter 16}:

\begin{equation}
\boxed{
P_{eff} = \sigma\left(
    \underbrace{
        \left[
            \underbrace{\tau \cdot \min\{AWX \cdot U\} + (1-\tau) \cdot \sum AWX \cdot U}_{WAX}
            - \underbrace{\theta_{base} + \sum \delta_j \Psi_j}_{\theta}
        \right]
    }_{WEC = WAX - \theta}
    \cdot \alpha_{BCJ} \cdot \beta_{BCS}
\right)
}
\label{eq:mp-peff}
\end{equation}

where $\sigma(x) = \frac{1}{1 + e^{-kx}}$ is the sigmoid function.

\subsection{Components of $P_{eff}$}

\begin{center}
\renewcommand{\arraystretch}{1.4}
\begin{tabular}{llll}
\toprule
\textbf{Component} & \textbf{Symbol} & \textbf{Source} & \textbf{Interpretation} \\
\midrule
Willingness & $WAX$ & Ch12, AV & Readiness to act \\
Threshold & $\theta$ & Ch12, AV & Action barrier (context-dependent) \\
Thinking Style & $\tau$ & Ch12 & Cognitive integration style \\
Awareness & $AWX$ & Ch11, AU & Salience of utility \\
Utility & $U$ & Ch10, C & FEPSDE components \\
Journey Factor & $\alpha_{BCJ}$ & Ch13, AW & Stage-dependent multiplier \\
Segment Factor & $\beta_{BCS}$ & Ch14, BA & Heterogeneity multiplier \\
\bottomrule
\end{tabular}
\end{center}

% =============================================================================
\section{The State Space: 10C Status Quo}
\label{sec:mp-state}
% =============================================================================

The state vector $S$ comprises all 10C dimensions (\textbf{Chapter 1b}):

\begin{equation}
\boxed{
S = (L, d, \gamma, \Psi, \Theta, A, W, \varphi, N_{L2})
}
\label{eq:mp-state-vector}
\end{equation}

\subsection{10C $\to$ State Mapping}

\begin{center}
\renewcommand{\arraystretch}{1.3}
\begin{tabular}{cllll}
\toprule
\textbf{10C} & \textbf{Question} & \textbf{Symbol} & \textbf{Chapter} & \textbf{CORE Appendix} \\
\midrule
WHO & Who has utility? & $L \in \{L_0, \ldots, L_\infty\}$ & Ch10 & AAA \\
WHAT & What is utility? & $d \in$ FEPSDE (144 comp.) & Ch10 & C \\
HOW & How interact? & $\gamma_{ij} \in [-1, +1]$ (15 params) & Ch5 & B \\
WHEN & Context matters? & $\Psi \in [0,1]^8$ & Ch9 & V \\
WHERE & Parameters from? & $\Theta \in \mathbb{R}^m$ & Ch4 & BBB \\
AWARE & How salient? & $A(\cdot) \in [0,1]$ & Ch11 & AU \\
READY & How willing? & $WAX, \theta$ & Ch12 & AV \\
STAGE & Journey phase? & $\varphi \in [0,1]$ (6 phases) & Ch13 & AW \\
HIERARCHY & Decision level? & $N_{L2} \in \{L_0, \ldots, L_3\}$ & Ch15 & HI \\
\bottomrule
\end{tabular}
\end{center}

% =============================================================================
\section{The Instrument Space: Interventions}
\label{sec:mp-instruments}
% =============================================================================

Interventions are vectors in continuous 9-dimensional space (\textbf{Chapter 17, Appendix UNMAPPED_IE}):

\begin{equation}
\boxed{
\vec{I} \in [0,1]^9
}
\label{eq:mp-intervention}
\end{equation}

\subsection{Intervention Vector Components}

\begin{equation}
\vec{I} = (I_{WHO}, I_{WHAT}, I_{HOW}, I_{WHEN}, I_{WHERE}, I_{AWARE}, I_{READY}, I_{STAGE}, I_{HIER})
\end{equation}

Each component $I_c \in [0,1]$ indicates the intensity of intervention on that 10C dimension.

\textbf{Key Insight (EIT-3):} Intervention ``types'' (A1--A8) are \textit{emergent clusters}
in this continuous space, not ontological primitives.

\subsection{Intervention Effectiveness}

From \textbf{Chapter 11/20} (EMERGE Algorithm, EIT-EMP-2), revised per \textbf{EXC-3 (Hybrid Formulation)}:

\begin{tcolorbox}[colback=green!5!white, colframe=green!75!black,
    title=\textbf{EXC-3 Compliance: K* Hybrid Formula}]
\textbf{Veto-Logic Analysis:} Only $f_B$ (Budget) satisfies veto logic: $B=0 \Rightarrow K^*=0$ (no budget $\to$ no interventions possible). Other factors reduce/scale but don't prevent interventions.

\textbf{Therefore:} $f_B$ is \textit{multiplicative}; $\delta_C, \delta_S, \delta_\gamma, \delta_L$ are \textit{additive} (EXC-1 default).
\end{tcolorbox}

\begin{equation}
\boxed{
K^* = K_{base} \cdot f_B + \delta_C + \delta_S + \delta_\gamma + \delta_L
}
\label{eq:mp-kstar-formula}
\end{equation}

where:
\begin{align}
f_B &= \sqrt{B/B_{ref}} & \text{[Budget factor, MULTIPLICATIVE---veto logic]} \\
\delta_C &= -0.2 \cdot K_{base} \cdot (c - c_{ref}) & \text{[Complexity adjustment, ADDITIVE]} \\
\delta_S &= K_{base} \cdot (\alpha_{BCJ}(\varphi) - 1) & \text{[Stage adjustment, ADDITIVE]} \\
\delta_\gamma &= 0.3 \cdot K_{base} \cdot \bar{\gamma} & \text{[Complementarity adjustment, ADDITIVE]} \\
\delta_L &= 0.1 \cdot K_{base} \cdot (L - 1) & \text{[Level adjustment, ADDITIVE]}
\end{align}

\textit{Reference:} Appendix UNMAPPED_FRM (Axiom EXC-3), \texttt{data/formula-registry.yaml} (ID: K*-Ch20)

% =============================================================================
\section{The Solution Space}
\label{sec:mp-solution-space}
% =============================================================================

The solution space defines what interventions are \textit{feasible} given context and state:

\begin{equation}
\boxed{
\mathcal{S}(\Psi_t, S_t) = \{I \in [0,1]^9 : \text{feasible}(I, \Psi_t, S_t)\}
}
\label{eq:mp-solution-space}
\end{equation}

\subsection{Solution Space per 10C Dimension}

\begin{center}
\renewcommand{\arraystretch}{1.3}
\begin{tabular}{lll}
\toprule
\textbf{10C} & \textbf{Solution Space} & \textbf{Constraints} \\
\midrule
WHO & $\mathcal{S}_{WHO}(\Psi) \subseteq \{L_0, \ldots, L_4\}$ & Access, authority, scalability \\
WHAT & $\mathcal{S}_{WHAT}(\Psi) \subseteq$ FEPSDE & Values, culture, preferences \\
HOW & $\mathcal{S}_{HOW}(\Psi) \subseteq [-1, +1]^{15}$ & Existing behavioral patterns \\
WHEN & $\mathcal{S}_{WHEN}(\Psi) \subseteq \Psi$ & Context controllability \\
WHERE & $\mathcal{S}_{WHERE}(\Psi) \subseteq \mathbb{R}^m$ & Data availability, measurement \\
AWARE & $\mathcal{S}_{AWARE}(\Psi) \subseteq [0,1]$ & Attention limits, cognitive load \\
READY & $\mathcal{S}_{READY}(\Psi) \subseteq [0,1]$ & Autonomy, reactance risk \\
STAGE & $\mathcal{S}_{STAGE}(\Psi) \subseteq \{S_1, \ldots, S_6\}$ & Phase receptivity ($\alpha$) \\
HIER & $\mathcal{S}_{HIER}(\Psi) \subseteq \{L_0, \ldots, L_3\}$ & Cognitive capacity \\
\bottomrule
\end{tabular}
\end{center}

\textbf{Full Solution Space:}
\begin{equation}
\mathcal{S}(\Psi, S) = \mathcal{S}_{WHO} \times \mathcal{S}_{WHAT} \times \cdots \times \mathcal{S}_{HIER} \cap \{\text{Budget, Ethics, Physics}\}
\end{equation}

\textit{See Appendix UNMAPPED_SS (CORE-SOLSPACE) for complete formalization.}

% =============================================================================
\section{The Dynamics}
\label{sec:mp-dynamics}
% =============================================================================

The optimization problem is \textit{dynamic}---state and context evolve over time.

\subsection{State Transition (EIT-2)}

From \textbf{Chapter 17, Appendix UNMAPPED_IE}:

\begin{equation}
\boxed{
S_{t+1} = S_t + \Delta S(I_t)
}
\label{eq:mp-state-transition}
\end{equation}

where $\Delta S(I)$ is the state change induced by intervention $I$.

\subsection{Context Evolution (Chapter 9)}

Context is \textit{endogenous}---it evolves based on behavior:

\begin{equation}
\boxed{
\Psi_{t+1} = f(\Psi_t, I_t, a_t)
}
\label{eq:mp-context-evolution}
\end{equation}

\textbf{Key Insight:} This creates a \textit{feedback loop}:
\begin{center}
$S_0, \Psi_0 \to \mathcal{S}(\Psi_0) \to I^*_0 \to S_1, \Psi_1 \to \mathcal{S}(\Psi_1) \to I^*_1 \to \cdots$
\end{center}

The solution space \textit{emerges} from context, and context \textit{evolves} through intervention.

\subsection{Portfolio Dynamics (Chapter 21)}

The optimal portfolio size evolves:

\begin{equation}
\boxed{
K^*(t) = K^*_0 \cdot g(\varphi_t) \cdot h(\Psi_t) \cdot \ell(\Theta_t)
}
\label{eq:mp-kstar-dynamics}
\end{equation}

\textbf{Redesign Triggers:}
\begin{itemize}
\item $|\Delta\varphi| > 0.2$ (Phase jump) $\to$ Recalculate $K^*(t)$
\item $||\Delta\Psi|| > 0.15$ (Context shock) $\to$ Expand portfolio, then refocus
\item $||\Delta\Theta|| < \epsilon$ for $T$ periods (Convergence) $\to$ Focus portfolio
\end{itemize}

\subsection{Lifespan Dynamics (Chapter 24)}

For long-term horizons (30+ years), domain-specific ODEs govern trajectory:

\begin{equation}
\frac{d\varphi_d}{dt} = \alpha_d (1 - \varphi_d) - \beta_d \varphi_d + \lambda_d E(I_t)
\label{eq:mp-lifespan}
\end{equation}

where $d \in \{\text{Career, Health, Money, Relationships, Self, Living, Meaning}\}$.

% =============================================================================
\section{The Context Structure (Chapter 1)}
\label{sec:mp-context-structure}
% =============================================================================

The context vector $\Psi$ determines \textit{when} different theoretical predictions apply.
Chapter~1 operationalizes $\Psi$ through eight primitive dimensions.

\subsection{The Eight Context Dimensions}

\begin{center}
\renewcommand{\arraystretch}{1.3}
\begin{tabular}{@{}clll@{}}
\toprule
\textbf{Symbol} & \textbf{Dimension} & \textbf{What It Captures} & \textbf{Key Measures} \\
\midrule
$\Psi_I$ & Institutional & Formal rules and enforcement & Rule of Law, Polity IV \\
$\Psi_S$ & Social & Trust and social capital & WVS Trust, Civic Index \\
$\Psi_C$ & Cognitive & Information processing capacity & PISA, Financial Literacy \\
$\Psi_K$ & Informational & Transparency and signal quality & Press Freedom, CPI \\
$\Psi_E$ & Economic & Development and resources & GDP/capita, Credit Access \\
$\Psi_T$ & Temporal & Stability and time horizons & EPU Index, Inflation Vol. \\
$\Psi_M$ & Market Scope & Geographic reach and openness & Trade/GDP, LPI \\
$\Psi_F$ & Factor Flexibility & Adjustment capacity & EPL Index, Mobility \\
\bottomrule
\end{tabular}
\end{center}

Each dimension is normalized to $[0,1]$:
\begin{equation}
\Psi = (\Psi_I, \Psi_S, \Psi_C, \Psi_K, \Psi_E, \Psi_T, \Psi_M, \Psi_F) \in [0,1]^8
\end{equation}

\subsection{Context Determines Theory Applicability}

The framework functions as a \textbf{metatheory}: it specifies \textit{when} each
component theory applies:

\begin{center}
\renewcommand{\arraystretch}{1.3}
\begin{tabular}{lll}
\toprule
\textbf{Theory} & \textbf{Applies When} & \textbf{Context Condition} \\
\midrule
Arrow-Debreu & No interactions & $\gamma_{ij} \approx 0$, $\dot{\Psi} \approx 0$ \\
Behavioral Econ & Individual biases dominate & $\Psi_C$ varies, others fixed \\
Institutional Econ & Rules dominate & $\Psi_I$ varies, behavior aggregates \\
Social Preferences & Social context matters & $\Psi_S$ high, $\gamma_{social} > 0$ \\
Full EBF & Multiple $\Psi$ vary & $\gamma \neq 0$, $\dot{\Psi} \neq 0$ \\
\bottomrule
\end{tabular}
\end{center}

\textbf{Key insight:} Existing theories tell you \emph{that} effects occur.
EBF tells you \emph{when}---which contexts amplify, dampen, or reverse each effect.

\subsection{The Level Salience Function}

Which aggregation level $L$ is salient depends on context (Chapter~10, Appendix~AAA):

\begin{equation}
\boxed{
\text{Total Utility} = \sum_{L=0}^{\infty} \alpha^L(\Psi) \cdot U^L
}
\label{eq:mp-level-salience}
\end{equation}

where $\alpha^L(\Psi)$ is the \textbf{level salience function} with $\sum_L \alpha^L(\Psi) = 1$.

\textbf{Context effects on level salience:}
\begin{itemize}
\item High $\Psi_S$ (trust) $\to$ Higher $\alpha^{L \geq 2}$ (community-level thinking)
\item Low $\Psi_E$ (scarcity) $\to$ Higher $\alpha^{L=0}$ (individual survival focus)
\item High $\Psi_I$ (institutions) $\to$ Higher $\alpha^{L \geq 4}$ (organizational thinking)
\end{itemize}

\subsection{Context-Dependent Treatment Effects}

The same intervention produces different effects in different contexts:

\begin{equation}
\tau(I, \Psi) = P_{eff}(S + \Delta S(I) \mid \Psi) - P_{eff}(S \mid \Psi)
\end{equation}

\textbf{Example (from Chapter~1):} A matching grant for charity:
\begin{itemize}
\item Switzerland ($\Psi_S = 0.9$, $\Psi_I = 0.9$, $\Psi_E = 0.9$): $\tau = +45\%$
\item Rural Peru ($\Psi_S = 0.5$, $\Psi_I = 0.4$, $\Psi_E = 0.3$): $\tau = +8\%$
\end{itemize}

The effect is not ``45\%'' or ``8\%''---it is $\tau(\Psi)$, a function of context.

% =============================================================================
\section{Rationality as Stability (Chapter 2)}
\label{sec:mp-rationality}
% =============================================================================

Chapter~2 establishes that classical rationality is \textit{necessary but insufficient}
for the EBF framework. It motivates the extension to $(C, \Psi)$ by identifying
the stability gap in classical economics.

\subsection{The Separability Assumption}

The Arrow-Debreu framework achieves equilibrium existence by assuming:
\begin{equation}
\boxed{
C_{ij} = \frac{\partial^2 U_i}{\partial x_i \partial x_j} = 0 \quad \text{for all } i \neq j
}
\label{eq:mp-separability}
\end{equation}

\textbf{Meaning:} Your consumption doesn't affect my utility. Your production doesn't
affect my costs. Agents are connected only through prices.

\textbf{EBF Relaxation:} We allow $C_{ij} \neq 0$ (complementarity), which:
\begin{itemize}
\item Creates utility interdependence (Chapter 5, Appendix UNMAPPED_B)
\item Generates non-concave landscapes (Chapter 7)
\item Enables social preferences and spillovers (Chapter 10, Appendix UNMAPPED_AAA)
\end{itemize}

\subsection{The Stability Gap (Sonnenschein-Mantel-Debreu)}

The Sonnenschein-Mantel-Debreu theorem shows that aggregate excess demand functions
can have almost any shape consistent with Walras' law. This implies:

\begin{itemize}
\item \textbf{Multiple equilibria} are possible
\item \textbf{Stability is not guaranteed}---the economy may not return to equilibrium after a shock
\item \textbf{Equilibrium selection is indeterminate}---theory cannot predict which equilibrium emerges
\end{itemize}

\textbf{Example: Coordination Failure}
\begin{center}
\renewcommand{\arraystretch}{1.2}
\begin{tabular}{l|cc}
 & Worker 2: Specialize & Worker 2: Autarky \\
\hline
Worker 1: Specialize & (3, 3) & (0, 2) \\
Worker 1: Autarky & (2, 0) & (2, 2) \\
\end{tabular}
\end{center}

Both (Specialize, Specialize) and (Autarky, Autarky) are stable equilibria.
The first is Pareto-superior, but Arrow-Debreu cannot predict which emerges.

\textbf{In EBF terms:} Both have $K \approx 1$ (coherent), but different $Q$ (welfare).
The economy can be \textit{stably trapped} in the inferior equilibrium.

\subsection{Evolutionary Foundation: Selection, Not Calculation}

Evolutionary game theory provides an alternative foundation for rationality
\citep{weibull1995, maynardsmith1982}:

\begin{equation}
\boxed{
\dot{x}_i = x_i \left[ \pi_i(x) - \bar{\pi}(x) \right]
}
\label{eq:mp-replicator}
\end{equation}

where $x_i$ is the population share using strategy $i$, $\pi_i$ is its payoff,
and $\bar{\pi}$ is population average payoff.

\textbf{Key insight:} Rationality \textit{emerges} as a result of selection, not as
an assumption about calculation. This has profound implications:

\begin{itemize}
\item \textbf{Bounded rationality is not a bug:} Heuristics evolved because they worked
      in ancestral environments. $\Psi_C$ reflects evolutionary calibration.
\item \textbf{Social preferences can survive:} Altruism is sustained through kin selection,
      reciprocity, group selection \citep{fehr2004nature}. $\Psi_S$ is evolutionarily stable.
\item \textbf{Multiple equilibria are natural:} Selection can stabilize different configurations.
      History determines which one emerges.
\end{itemize}

\textbf{Framework integration:}
\begin{itemize}
\item $C^*$ is not calculated but \textit{emerges} from evolutionary dynamics
\item Multiple $C^*$ may exist; selection determines which is reached
\item Bounded rationality ($\Psi_C < \infty$) is evolutionarily grounded
\end{itemize}

\subsection{Quantal Response Equilibrium (QRE)}

\citet{mckelvey1995} bridged rationality and bounded rationality:

\begin{equation}
\boxed{
P(s_i) = \frac{\exp(\lambda \cdot \pi_i)}{\sum_j \exp(\lambda \cdot \pi_j)}
}
\label{eq:mp-qre}
\end{equation}

where $\lambda$ measures rationality: $\lambda = 0$ is random choice, $\lambda \to \infty$
is fully rational.

\textbf{Connection to EBF:} QRE operationalizes $\Psi_C$ (cognitive context):
\begin{itemize}
\item Lower cognitive capacity $\to$ lower $\lambda$ $\to$ noisier decisions
\item Higher cognitive load $\to$ behavior closer to random
\item This connects to AWARE (Chapter 11): attention limits affect $\lambda$
\end{itemize}

\subsection{The Coherence Index $K$}

Chapter~2 introduces the fundamental distinction between coherence and welfare:

\begin{tcolorbox}[colback=yellow!5!white, colframe=yellow!75!black,
    title=\textbf{K vs Q: The Central Distinction}]

\begin{itemize}
\item $K \approx 1$ means you are \textbf{AT a peak} (internally consistent, coherent)
\item $Q$ measures the \textbf{HEIGHT of the peak} (welfare, quality)
\item High $K$ does NOT imply high $Q$---you may be at the \textit{wrong} peak
\item Context $\Psi$ determines which peak has highest $Q$
\end{itemize}

\end{tcolorbox}

\textbf{Historical examples:}
\begin{center}
\renewcommand{\arraystretch}{1.3}
\begin{tabular}{llll}
\toprule
\textbf{Period} & \textbf{K} & \textbf{Q} & \textbf{Description} \\
\midrule
Gold Standard (1870--1914) & $\approx 1$ & Moderate & Coherent rules, constrained growth \\
Weimar (1921--1923) & $\to 0$ & Collapsed & Complete institutional breakdown \\
Post-War (1945--1970) & $\approx 1$ & High & Bretton Woods + welfare state \\
2008 Financial Crisis & $0.95 \to 0.3$ & Sharp drop & Correlation breakdown, trust collapse \\
\bottomrule
\end{tabular}
\end{center}

\subsection{Implications for the Optimization Problem}

Chapter~2's rationality analysis has direct implications for MP:

\begin{enumerate}
\item \textbf{Multiple local optima} (from stability gap): The Bellman equation
      (\ref{eq:mp-bellman}) may have multiple solutions. Global search required.

\item \textbf{Selection dynamics} (from evolutionary theory): Which optimum is
      reached depends on initial conditions $(S_0, \Psi_0)$ and path.

\item \textbf{Bounded rationality} (from QRE): Agents don't solve MP exactly.
      Actual behavior is noisy, with $\lambda$ parameterizing rationality.

\item \textbf{K-Q trade-off}: Interventions may increase $K$ (stability) while
      decreasing $Q$ (welfare), or vice versa. Both must be tracked.
\end{enumerate}

% =============================================================================
\section{Limits of Classical Utility (Chapter 3)}
\label{sec:mp-limits}
% =============================================================================

Chapter~3 documents the empirical evidence that \textit{necessitates} $C_{ij} \neq 0$.
Six behavioral domains systematically violate the classical assumptions,
providing calibration data for the complementarity structure.

\subsection{The Six Anomaly Domains}

\begin{center}
\renewcommand{\arraystretch}{1.3}
\begin{tabular}{@{}clll@{}}
\toprule
\textbf{\#} & \textbf{Domain} & \textbf{Classical Violation} & \textbf{Key Parameter} \\
\midrule
1 & Social Preferences & $U_i(x_i, x_j) \neq U_i(x_i)$ & $\alpha, \beta$ (Fehr-Schmidt) \\
2 & Identity \& Self-Image & $U = U^{mat} + U^{IDN}$ & Identity premium \\
3 & Reference Dependence & $U(x) \to U(x - r)$ & $\lambda \approx 2$ (loss aversion) \\
4 & Framing Effects & $\text{Choice}(A,B;\Psi_1) \neq \text{Choice}(A,B;\Psi_2)$ & Frame sensitivity \\
5 & Preference Construction & $U(x;\Psi) = f(\text{process}(\Psi))$ & Construction weight \\
6 & Temporal/Present Bias & $\beta < 1$ in $\beta$-$\delta$ model & $\beta \approx 0.7$ \\
\bottomrule
\end{tabular}
\end{center}

\subsection{Four Types of Utility Interdependence}

The six domains generate four structural interdependencies:

\begin{tcolorbox}[colback=blue!5!white, colframe=blue!75!black,
    title=\textbf{The Four Interdependencies (Chapter 3)}]

\begin{enumerate}
\item \textbf{Cross-agent} (Social Preferences):
$\displaystyle \frac{\partial U_i}{\partial x_j} \neq 0$

\item \textbf{Temporal} (Habits, Identity):
$\displaystyle \frac{\partial U_t}{\partial a_{t-1}} \neq 0$

\item \textbf{Intertemporal Inconsistency} (Present Bias):
$\displaystyle U_t(c_t, c_{t+1}) \neq \delta \cdot U_{t+1}(c_{t+1}, c_{t+2})$

\item \textbf{Context Dependence} (Framing, Norms):
$\displaystyle \frac{\partial U}{\partial \Psi} \neq 0$
\end{enumerate}

\textbf{Key insight:} All four generate non-zero cross-derivatives $C_{ij} \neq 0$,
creating the complementarity structure that classical theory assumes away.

\end{tcolorbox}

\subsection{Distribution of Social Preference Types}

Empirical evidence shows that purely selfish agents are a \textit{minority}:

\begin{center}
\renewcommand{\arraystretch}{1.2}
\begin{tabular}{lcc}
\toprule
\textbf{Type} & \textbf{Students} & \textbf{General Population} \\
\midrule
Selfish & 40--60\% & 20--30\% \\
Inequality averse & 10--15\% & 15--20\% \\
Altruistic & 25--45\% & 45--55\% \\
\bottomrule
\end{tabular}
\end{center}

\textbf{Implication:} Population heterogeneity in $C_{ij}$ is substantial.
The BCS segmentation (Chapter 14, Appendix BA) must account for this distribution.

\subsection{The $\beta$-$\delta$ Model of Present Bias}

\citet{laibson1997} formalized quasi-hyperbolic discounting:
\begin{equation}
\boxed{
U_t = u(c_t) + \beta \sum_{\tau=1}^{T} \delta^\tau u(c_{t+\tau}) \quad \text{where } \beta < 1
}
\label{eq:mp-beta-delta}
\end{equation}

\textbf{Empirical estimates:} $\beta \approx 0.7$, $\delta \approx 0.99$

\textbf{Implications for MP:}
\begin{itemize}
\item \textbf{Time inconsistency:} Plans made today are not followed tomorrow
\item \textbf{Commitment demand:} Agents value devices that restrict future choices
\item \textbf{Sophisticates vs. naïfs:} $\hat{\beta} = \beta$ vs. $\hat{\beta} = 1$
\item \textbf{Interaction:} $\Psi_T \times \Psi_C$ (cognitive load worsens self-control)
\end{itemize}

\subsection{Calibration Data Mapping}

Each domain provides calibration targets for $C^*(\Psi)$:

\begin{center}
\small
\renewcommand{\arraystretch}{1.2}
\begin{tabular}{@{}lll@{}}
\toprule
\textbf{Domain} & \textbf{What Varies} & \textbf{Calibration Target} \\
\midrule
Social preferences & Cultures, stakes, merit & $C_{ij}(\Psi_S, \Psi_I)$ \\
Identity & Groups, priming, salience & $U^{IDN}(\Psi_S, \text{group})$ \\
Reference dependence & Domains, expectations & $r(\Psi)$, $\lambda(\Psi)$ \\
Framing & Presentations, defaults & Frame sensitivity $= f(\Psi_C)$ \\
Preference construction & Familiarity, complexity & Construction weight $= f(\Psi_K)$ \\
Temporal preferences & Cultures, cognitive load & $\beta(\Psi_C, \Psi_T)$ \\
\bottomrule
\end{tabular}
\end{center}

\textbf{Methodological shift:} From \textit{deriving} behavior from axioms to
\textit{calibrating} behavior on experimental variation.

\subsection{Implications for the Optimization Problem}

Chapter~3's anomaly evidence has direct implications for MP:

\begin{enumerate}
\item \textbf{Heterogeneous $C_{ij}$}: Different population segments have different
      complementarity structures. The solution space $\mathcal{S}(\Psi, S)$ varies by segment.

\item \textbf{Time-inconsistent agents}: The Bellman equation (\ref{eq:mp-bellman})
      must account for $\beta < 1$. Commitment devices become intervention instruments.

\item \textbf{Reference-dependent utility}: The objective $P_{eff}$ depends on
      reference points $r(\Psi)$, not just final outcomes.

\item \textbf{Context as intervention target}: Since $\partial U/\partial \Psi \neq 0$,
      changing context (framing, defaults, norms) is a legitimate intervention.

\item \textbf{Propositional $\neq$ functional knowledge}: LLMs know ``that'' effects
      exist but not ``when'' or ``how strongly''---calibration on data is essential.
\end{enumerate}

% =============================================================================
\section{Empirical Foundations (Chapter 4)}
\label{sec:mp-empirical}
% =============================================================================

Chapter~4 establishes the empirical basis for the framework. The parameter set $\Theta$
(CORE-WHERE, Appendix UNMAPPED_BBB) derives from 1,900+ experimental papers, field studies,
and cross-cultural research.

\subsection{Experimental Evidence for $C_{ij} \neq 0$}

\begin{center}
\renewcommand{\arraystretch}{1.3}
\begin{tabular}{@{}llll@{}}
\toprule
\textbf{Experiment} & \textbf{Classical Prediction} & \textbf{Actual Result} & \textbf{Implication} \\
\midrule
Ultimatum Game & Offer \$0.01, accept & Modal offer 40-50\% & Fairness matters \\
Dictator Game & Keep 100\% & Mean transfer 20-30\% & Altruism exists \\
Public Goods & Contribute 0\% & Initial 40-60\% & Conditional cooperation \\
Trust Game & Send \$0, return \$0 & Send 50\%, return 95\% & Reciprocity \\
\bottomrule
\end{tabular}
\end{center}

\textbf{Key insight:} All four games violate $C_{ij} = 0$. The cross-derivative
$\partial^2 U_i / \partial x_i \partial x_j \neq 0$ is empirically established.

\subsection{Calibrated Social Preference Parameters}

The Fehr-Schmidt model parameters estimated from experiments:

\begin{equation}
U_i(x) = x_i - \alpha_i \cdot \frac{1}{n-1} \sum_{j \neq i} \max(x_j - x_i, 0)
              - \beta_i \cdot \frac{1}{n-1} \sum_{j \neq i} \max(x_i - x_j, 0)
\label{eq:mp-fehr-schmidt}
\end{equation}

\begin{center}
\renewcommand{\arraystretch}{1.2}
\begin{tabular}{lll}
\toprule
\textbf{Parameter} & \textbf{Meaning} & \textbf{Estimate} \\
\midrule
$\alpha$ & Envy (disadvantageous inequity) & Mean $\approx 0.85$, range $[0, 4+]$ \\
$\beta$ & Guilt (advantageous inequity) & Mean $\approx 0.315$, range $[0, 0.6]$ \\
\bottomrule
\end{tabular}
\end{center}

\textbf{Constraint:} $\beta_i \leq \alpha_i$ (guilt $\leq$ envy)

\textbf{Heterogeneity:} High variance across individuals---motivates BCS segmentation (Chapter 14).

\subsection{Cross-Cultural Variation}

The structure $C_{ij} \neq 0$ is universal; the parameters are culturally specific:

\begin{center}
\renewcommand{\arraystretch}{1.2}
\begin{tabular}{lcc}
\toprule
\textbf{Society} & \textbf{Ultimatum Offer} & \textbf{Context ($\Psi$)} \\
\midrule
Machiguenga (Peru) & 26\% & Low market integration \\
WEIRD Samples & 40-45\% & High market integration \\
Lamelara (Indonesia) & 58\% & Strong cooperation norms \\
\bottomrule
\end{tabular}
\end{center}

\textbf{Implication:} $C_{ij} = C_{ij}(\Psi_{culture})$---complementarity is context-dependent.

\subsection{The Five Empirical Facts}

Chapter~4 establishes five facts that MP takes as axiomatic:

\begin{tcolorbox}[colback=green!5!white, colframe=green!75!black,
    title=\textbf{Five Empirical Facts (Chapter 4)}]

\begin{enumerate}
\item \textbf{Complementarities are real:} $C_{ij} \neq 0$ confirmed across 100s of experiments
\item \textbf{Complementarities are measurable:} Structural estimation recovers $(\alpha, \beta)$
\item \textbf{Complementarities are context-dependent:} Same individuals, different $C_{ij}(\Psi)$
\item \textbf{Complementarities have neural substrates:} Social and material rewards share circuits
\item \textbf{Complementarities vary cross-culturally:} Universal structure, culture-specific parameters
\end{enumerate}

\end{tcolorbox}

\subsection{The Diagnostic Protocol}

When $C$ deviates from $C^*(\Psi)$, deviations are classified into five types:

\begin{center}
\small
\renewcommand{\arraystretch}{1.2}
\begin{tabular}{@{}llp{5cm}@{}}
\toprule
\textbf{Type} & \textbf{Signature} & \textbf{Action} \\
\midrule
I: Random & Uncorrelated & Better measurement instruments \\
II: $\Psi$-Systematic & Correlates with $\Psi_k$ & Enrich $C^*(\Psi)$ with interactions \\
III: Extra-$\Psi$ & Correlates with $Z \notin \Psi$ & Test candidate 9th dimension \\
IV: Temporal & Autocorrelated & Model transition dynamics \\
V: Localized & Concentrated in units & Case study \\
\bottomrule
\end{tabular}
\end{center}

\textbf{Self-improving architecture:} Deviations are not noise but information.
The framework is systematically improvable (Lakatos progressive research program).

\subsection{Implications for the Optimization Problem}

Chapter~4's empirical foundations have direct implications for MP:

\begin{enumerate}
\item \textbf{Parameter source:} $\Theta$ in equation (\ref{eq:mp-objective}) comes from
      1,900+ papers, not assumption. See Appendix UNMAPPED_BBB (CORE-WHERE).

\item \textbf{Calibration, not simulation:} LLMs know ``that'' behavioral effects exist
      but not ``when'' or ``how strongly.'' Prediction requires calibration on experimental
      variation, not training on text.

\item \textbf{Heterogeneity in $C$:} Population variance in $(\alpha, \beta)$ means
      different segments face different optimization landscapes. The solution space
      $\mathcal{S}(\Psi, S)$ is segment-specific.

\item \textbf{Cultural variation:} Cross-cultural differences in $C_{ij}(\Psi_{culture})$
      mean interventions must be context-adapted. Swiss parameters $\neq$ Peruvian parameters.

\item \textbf{Continuous learning:} The diagnostic protocol enables self-reinforcement.
      Deviations update $\Theta$; $\Theta$ updates predictions; predictions generate new deviations.
\end{enumerate}

\subsection{Calibration vs. Simulation (Chapter 4x)}

Chapter~4x establishes the fundamental distinction: \textbf{calibration on behavioral data}
outperforms \textbf{LLM-based simulation}.

\begin{tcolorbox}[colback=red!5!white, colframe=red!75!black,
    title=\textbf{The Peng et al. Finding (2025)}]

Digital twins (LLMs conditioned on individual-level data) achieve:
\begin{itemize}
\item Individual-level accuracy: 75\% (no better than demographic baseline)
\item Correlation with human responses: $r \approx 0.2$
\item Variance: Under-dispersed (heterogeneity lost)
\end{itemize}

\textbf{Pattern of failure:} Digital twins show theoretically expected behavior that
humans \textit{don't} show, and fail to show actual patterns humans \textit{do} exhibit.

\end{tcolorbox}

\textbf{The Knowledge Distinction:}
\begin{itemize}
\item \textbf{Propositional knowledge} (LLMs have): ``Loss aversion exists and $\lambda \approx 2$''
\item \textbf{Functional knowledge} (prediction requires): $\lambda(\Psi) = f(\Psi_C, \Psi_T, \Psi_E, \text{stake}, \text{domain}, \ldots)$
\end{itemize}

LLMs know \textit{that} effects exist but not \textit{when} or \textit{how strongly}.

\textbf{Architecture Comparison:}
\begin{center}
\small
\renewcommand{\arraystretch}{1.2}
\begin{tabular}{@{}lll@{}}
\toprule
\textbf{Aspect} & \textbf{LLM / Digital Twin} & \textbf{EBF / $C^*(\Psi)$} \\
\midrule
Data source & Text corpora & Experimental behavioral data \\
Knowledge form & ``Effect is...'' & $E(\Psi) = \alpha + \beta_C\Psi_C + \ldots$ \\
Replication & $\sim$50\% (Peng) & $>$85\% (by construction) \\
Heterogeneity & Averaged away & Explicitly modeled as $\sigma^2(\Psi)$ \\
\bottomrule
\end{tabular}
\end{center}

\textbf{Scaling Logic:}
\begin{center}
\small
\renewcommand{\arraystretch}{1.2}
\begin{tabular}{rccc}
\toprule
\textbf{Experiments} & \textbf{Parameters} & \textbf{What's Identifiable} & \textbf{Expected $R^2_{OOS}$} \\
\midrule
50 & $\sim$50-100 & Main effects & 0.3-0.5 \\
500 & $\sim$500-1,000 & + Two-way interactions & 0.5-0.7 \\
1,500 & $\sim$2,000-5,000 & + Nonlinearities, thresholds & 0.7-0.85 \\
5,000+ & $\sim$10,000+ & + Domain transfer, dynamics & 0.85+ \\
\bottomrule
\end{tabular}
\end{center}

\textbf{The Testable Claim:} EBF calibrated on $N$ experiments outperforms (a) naive baselines,
(b) LLM-based predictions, and (c) expert predictions on held-out experiments.

% =============================================================================
\section{Complementarity as Structural Principle (Chapter 5)}
\label{sec:mp-complementarity}
% =============================================================================

Chapter~5 answers the 3rd CORE question: \textbf{HOW do utility components interact?}
This is the mathematical foundation of the $\gamma$-structure in the optimization problem.

\subsection{The Formal Definition}

\begin{equation}
\boxed{
\gamma_{ij} = \frac{\partial^2 U}{\partial x_i \partial x_j}
\quad \begin{cases}
> 0 & \text{complements (doing more of } i \text{ raises return to } j\text{)} \\
< 0 & \text{substitutes (doing more of } i \text{ lowers return to } j\text{)} \\
= 0 & \text{independent (classical assumption)}
\end{cases}
}
\label{eq:mp-gamma-def}
\end{equation}

\textbf{Supermodularity:} A function $f$ is supermodular if $\partial^2 f / \partial x_i \partial x_j \geq 0$ for all $i \neq j$, equivalently: $f(x \vee y) + f(x \wedge y) \geq f(x) + f(y)$.

\subsection{Four Types of Complementarity}

\begin{center}
\renewcommand{\arraystretch}{1.3}
\begin{tabular}{@{}llll@{}}
\toprule
\textbf{Type} & \textbf{Definition} & \textbf{Example} & \textbf{Consequence} \\
\midrule
Strategic & $\frac{\partial^2 \Pi_i}{\partial a_i \partial a_j} > 0$ & Bank runs, adoption & Multiple equilibria \\
Temporal & $\frac{\partial^2 U}{\partial a_t \partial a_{t+1}} > 0$ & Habits, identity & Path dependence \\
Cross-Domain & $\frac{\partial^2 U}{\partial a^A \partial a^B} > 0$ & Health-wealth & Spillovers \\
Cross-Agent & $\frac{\partial^2 U_i}{\partial x_i \partial x_j} > 0$ & Networks, norms & Coordination problems \\
\bottomrule
\end{tabular}
\end{center}

\subsection{The $\Gamma$-Matrix Structure}

The complete complementarity structure has 163 parameters:

\begin{center}
\renewcommand{\arraystretch}{1.2}
\begin{tabular}{lll}
\toprule
\textbf{Matrix} & \textbf{Parameters} & \textbf{Meaning} \\
\midrule
$\Gamma = [\gamma_{dd'}]$ & 15 & FEPSDE dimension interactions \\
$\Lambda = [\lambda_{kk'}]$ & 28 & Context ($\Psi$) interactions \\
$\delta = [\delta_{dd',k}]$ & 120 & Modulation tensor \\
\midrule
\textbf{Total} & \textbf{163} & Full parameter space \\
\bottomrule
\end{tabular}
\end{center}

\subsection{Context-Dependent Complementarity}

Complementarity is not constant---it depends on context $\Psi$:
\begin{equation}
\boxed{
\gamma_{dd'}(\Psi) = \bar{\gamma}_{dd'} + \sum_k \delta_{dd',k} \cdot (\Psi_k - 0.5)
}
\label{eq:mp-gamma-context}
\end{equation}

\textbf{Example:} Trust complementarity $\gamma_{trust}$ is higher in high-$\Psi_I$ (institutional quality) contexts.

\subsection{The Coherence Index Revisited}

Chapter~5 formalizes the coherence index from Chapter~2:
\begin{equation}
K = 1 - \frac{\|C - C^*\|}{\|C^*\|}
\label{eq:mp-coherence-ch5}
\end{equation}

\begin{itemize}
\item $K \approx 1$: Perfect alignment with reference structure (stable)
\item $K \ll 1$: Misalignment (instability, crisis)
\item $K$ vs $Q$: Coherence differs from quality (stable trap possible)
\end{itemize}

\subsection{Implications for the Optimization Problem}

Chapter~5's complementarity structure has direct implications for MP:

\begin{enumerate}
\item \textbf{Non-separable objective:} The payoff function in (\ref{eq:mp-objective}) includes
      $\gamma$-terms: $\Pi = \sum_d w_d u_d + \sum_{d<d'} \gamma_{dd'} u_d u_{d'}$.

\item \textbf{Multiple equilibria:} Strong complementarity ($\gamma > \gamma_{crit}$) creates
      multiple local optima. Global search required.

\item \textbf{Context-dependent landscape:} Since $\gamma = \gamma(\Psi)$, the optimization
      landscape changes with context. Same intervention, different $\Psi$ → different effect.

\item \textbf{163 parameters to estimate:} The solution space $\mathcal{S}(\Psi, S)$ depends
      on calibrated $\Gamma$, $\Lambda$, and $\delta$ values from Appendix UNMAPPED_BBB.

\item \textbf{Classical as special case:} When $\gamma_{ij} = 0$ for all $i,j$ (Arrow-Debreu),
      the optimization problem reduces to independent maximization.
\end{enumerate}

% =============================================================================
\section{The Reference Structure $C^*$ (Chapter 6)}
\label{sec:mp-reference}
% =============================================================================

Chapter~6 defines the \textbf{reference structure} $C^*$---the self-consistent configuration
against which coherence is measured. This is foundational for the optimization problem:
$K = 1 - \|C - C^*\|/\|C^*\|$ requires knowing $C^*$.

\subsection{The Fixed-Point Definition}

\begin{equation}
\boxed{
C^* = \Phi(C^*, \Psi)
}
\label{eq:mp-cstar-fixedpoint}
\end{equation}

$C^*$ is the unique structure where beliefs reproduce themselves through experience.
Under regularity (D-REG), reciprocity (D-REC), and contraction (D-CON) axioms,
Banach fixed-point theorem guarantees existence and uniqueness.

\subsection{Three Convergent Derivations}

\begin{center}
\renewcommand{\arraystretch}{1.3}
\begin{tabular}{@{}llll@{}}
\toprule
\textbf{Derivation} & \textbf{Method} & \textbf{Condition} & \textbf{Result} \\
\midrule
Axiomatic & Fixed-point iteration & $\|\Phi(C_1) - \Phi(C_2)\| \leq k\|C_1 - C_2\|$, $k<1$ & Uniqueness \\
Evolutionary & ESS analysis & $\pi(C^*, C^*) > \pi(C', C^*)$ for all $C' \neq C^*$ & Stability \\
Empirical & Ergodic average & $C^*_{emp} = \lim_{T\to\infty} \frac{1}{T}\sum_t C_t$ & Observability \\
\bottomrule
\end{tabular}
\end{center}

\textbf{Key Result:} All three derivations yield the same $C^*$:
\begin{equation}
C^*_{axiomatic} = C^*_{evolutionary} = C^*_{empirical}
\label{eq:mp-cstar-convergence}
\end{equation}

\subsection{Context Dependence}

$C^*$ is not a constant---it depends on context:
\begin{equation}
C^*(\Psi_1) \neq C^*(\Psi_2) \quad \text{for different contexts}
\label{eq:mp-cstar-context}
\end{equation}

\textbf{Implication for optimization:} The target coherence structure $C^*$ shifts with $\Psi$.
Interventions must track this moving target.

\subsection{Convergence Rate}

The approach to $C^*$ is exponentially fast:
\begin{equation}
\|C_n - C^*\| \leq k^n \|C_0 - C^*\|
\label{eq:mp-cstar-rate}
\end{equation}

where $k < 1$ is the contraction constant. Typical values: $k \approx 0.7$, so 8 iterations achieve 90\% convergence.

\subsection{The Stigler-Becker Insight}

Chapter~6 preserves the Stigler-Becker stability insight at a higher level:
not stable preferences $U$, but stable reference structure $C^*(\Psi)$.
The advance: $C^*$ is \textit{derived} from axioms, not exogenously assumed.

\subsection{Implications for the Optimization Problem}

\begin{enumerate}
\item \textbf{Coherence target:} The optimization seeks $K \to 1$, meaning $C \to C^*(\Psi)$.
\item \textbf{Moving target:} Since $C^* = C^*(\Psi)$ and $\Psi$ evolves, the target shifts.
\item \textbf{Calibration requirement:} Parameters $\alpha, \beta, \gamma$ in $\Phi$ must be calibrated from experimental data (Appendix UNMAPPED_BBB).
\item \textbf{Multiple $C^*$:} For fixed $\Psi$, $C^*$ is unique; across contexts, different $C^*(\Psi)$ exist.
\end{enumerate}

% =============================================================================
\section{The Landscape Structure (Chapter 7)}
\label{sec:mp-landscape}
% =============================================================================

\begin{tcolorbox}[colback=orange!5!white, colframe=orange!75!black,
    title=\textbf{Critical: Non-Concavity and Multiple Equilibria}]

The optimization landscape is \textbf{NOT concave}. Complementarity ($\gamma > 0$) creates:

\begin{itemize}
\item \textbf{Multiple local maxima} (peaks): Different coherent configurations
\item \textbf{Valleys}: Incoherent configurations between peaks
\item \textbf{Path dependence}: Which peak you reach depends on your path
\end{itemize}

\textbf{The K vs Q Distinction:}
\begin{itemize}
\item $K \approx 1$ means you are AT a peak (coherent)
\item $Q$ measures the HEIGHT of the peak (welfare)
\item High $K$ does NOT imply high $Q$---you may be at the wrong peak
\item Context $\Psi$ determines which peak has highest $Q$
\end{itemize}

\end{tcolorbox}

\subsection{Proposition: Non-Concavity from Complementarity}

\begin{proposition}[MP-P1: Non-Concavity]
If $\frac{\partial^2 \Pi}{\partial x_i \partial x_j} > 0$ for all $i \neq j$ (supermodularity),
then the payoff landscape $\Pi$ is not globally concave and has multiple local maxima.
\end{proposition}

\textit{Proof: See Chapter 7, Proposition 7.1, and Appendix UNMAPPED_NC (FORMAL-NC).}

\subsection{Implication: Valley-Crossing Problem}

To transition from Peak A to Peak B:
\begin{enumerate}
\item Leave Peak A (coherence falls: $K \downarrow$)
\item Cross the valley (incoherence, temporary performance loss)
\item Climb to Peak B (coherence restored: $K \uparrow$)
\end{enumerate}

\textbf{Sequential change fails}---you must change all complementary elements \textit{simultaneously}.

\subsection{Roberts' Fit Concept}

\citet{roberts2004} explains why successful firms look different yet internally coherent:
\textbf{fit} = organizational elements are supermodular complements.

\begin{center}
\renewcommand{\arraystretch}{1.2}
\begin{tabular}{@{}lll@{}}
\toprule
\textbf{Roberts (2004)} & \textbf{EBF Framework} & \textbf{Meaning} \\
\midrule
Fit & $K \approx 1$ & At a local maximum \\
Misfit & $K < 1$ & In the valley between peaks \\
Configuration quality & $Q(C^*)$ & Height of the peak \\
Multiple configurations & Multiple $C^*$ & Multiple local maxima \\
\bottomrule
\end{tabular}
\end{center}

\textbf{Examples of coherent configurations:}
\begin{itemize}[nosep]
\item \textbf{Toyota:} Lean + empowerment + continuous improvement + lifetime employment
\item \textbf{Southwest:} Low cost + point-to-point + single aircraft + fun culture
\item \textbf{Goldman:} Partnership + high risk + aggressive incentives + up-or-out
\end{itemize}

Each is $K \approx 1$ but incompatible with each other's elements.

\subsection{Coherent Obsolescence: The K-Q Trap}

\textbf{Kodak (2000):} Highly coherent ($K \approx 1$) in film photography.
But context shifted ($\Psi \to$ digital). Result: High $K$, collapsing $Q$.

\begin{center}
\renewcommand{\arraystretch}{1.2}
\begin{tabular}{@{}lcc@{}}
\toprule
\textbf{Metric} & \textbf{1995} & \textbf{2005} \\
\midrule
Coherence $K$ & $\approx 0.95$ & $\approx 0.90$ \\
Peak Quality $Q$ & $\approx 0.85$ & $\approx 0.25$ \\
Context $\Psi_{digital}$ & 0.1 & 0.7 \\
\bottomrule
\end{tabular}
\end{center}

\textbf{Lesson:} $K$ measures fit with \textit{self}; $Q$ measures fit with \textit{world}.

\subsection{Transition Cost Structure}

\begin{center}
\renewcommand{\arraystretch}{1.2}
\begin{tabular}{@{}lccc@{}}
\toprule
\textbf{Change Strategy} & $\Delta s$ & $\Delta \sigma$ & \textbf{Result} \\
\midrule
Full transformation & $+0.8$ & $+0.8$ & Reach new peak \\
Strategy only & $+0.8$ & $0$ & Stuck in valley \\
Structure only & $0$ & $+0.8$ & Stuck in valley \\
Gradual (sequential) & staged & staged & Prolonged valley \\
\bottomrule
\end{tabular}
\end{center}

\subsection{Multi-Level Fit Application}

The complementarity-fit-non-concavity logic applies at all levels:

\begin{center}
\renewcommand{\arraystretch}{1.2}
\small
\begin{tabular}{@{}llll@{}}
\toprule
\textbf{Level} & \textbf{Elements} & \textbf{Fit Means} & \textbf{Example} \\
\midrule
Individual & Values, goals, habits & Integrity & Live your values \\
Household & Roles, expectations & Harmony & Shared parenting \\
Organization & Strategy, structure, culture & Competitive fit & Southwest \\
Region & Economy, education & Economic coherence & Silicon Valley \\
Nation & Institutions, norms, policies & Institutional fit & Nordic model \\
\bottomrule
\end{tabular}
\end{center}

\subsection{The Calibration Imperative for Navigation}

Strategic questions require \textit{quantitative} answers:
\begin{enumerate}[nosep]
\item How deep is the valley? (transition cost)
\item How high is the target peak? (potential gain)
\item What is $P(success)$? (transformation probability)
\end{enumerate}

Text-based analysis explains \textit{that} complementarity exists, but cannot quantify \textit{how much}.
The landscape must be calibrated from data, not described from principles.

% =============================================================================
\section{Mathematical Foundations (Chapter 8)}
\label{sec:mp-mathematical}
% =============================================================================

Chapter~8 provides the formal mathematical machinery for the optimization problem,
bridging intuition with rigorous structure.

\subsection{The Two-Axis Configuration Space}

Following \citet{roberts2004}, configurations reduce to two fundamental dimensions:
\begin{equation}
\boxed{
(s, \sigma) \in [0,1]^2
\quad \begin{cases}
s: & \text{Strategy (stability } 0 \leftrightarrow \text{ dynamism } 1\text{)} \\
\sigma: & \text{Structure (hierarchy } 0 \leftrightarrow \text{ flexibility } 1\text{)}
\end{cases}
}
\label{eq:mp-config-space}
\end{equation}

\subsection{The Formal Payoff Function}

\begin{equation}
\boxed{
\Pi(s, \sigma; \Psi) = \underbrace{-\gamma (s - \sigma)^2}_{\text{Misfit penalty}}
                      + \underbrace{\alpha(\Psi) \cdot s \cdot \sigma}_{\text{Innovation synergy}}
                      + \underbrace{\beta(\Psi) \cdot (1-s)(1-\sigma)}_{\text{Stability synergy}}
}
\label{eq:mp-payoff-formal}
\end{equation}

\textbf{Equilibrium structure:}
\begin{center}
\renewcommand{\arraystretch}{1.2}
\begin{tabular}{@{}lccc@{}}
\toprule
\textbf{Point} & $(s, \sigma)$ & $\Pi$ & \textbf{Type} \\
\midrule
Peak A & $(0, 0)$ & $\beta$ & Local max (stability) \\
Peak B & $(1, 1)$ & $\alpha$ & Local max (innovation) \\
Saddle M & $(0.5, 0.5)$ & $(\alpha+\beta)/4$ & Saddle point \\
Misfit & $(0, 1), (1, 0)$ & $-\gamma$ & Local minima \\
\bottomrule
\end{tabular}
\end{center}

\subsection{The 144-Component Welfare Structure}

The full welfare function aggregates 144 utility components:
\begin{equation}
W = \sum_{i \in \{INU, KNU, IDN\}} \sum_{d \in FEPSDE} \sum_{t=0}^{3}
    \delta_d(\Psi)^t \left[ G_{i,d,t} - \lambda_d(\Psi) \cdot P_{i,d,t} \right]
\label{eq:mp-144-welfare}
\end{equation}

\textbf{Component structure:} $3 \times 2 \times 6 \times 4 = 144$
\begin{itemize}[nosep]
\item $i \in \{INU, KNU, IDN\}$: Utility category (Individual, Collective, Identity)
\item $v \in \{G, P\}$: Valence (Gain, Pain)
\item $d \in FEPSDE$: Six dimensions (Financial, Emotional, Physical, Social, Development, Ecological)
\item $t \in \{0,1,2,3\}$: Time horizon
\end{itemize}

\subsection{The Complementarity Matrix}

The full matrix $C \in \mathbb{R}^{144 \times 144}$ has block structure:
\begin{equation}
C = \begin{pmatrix}
  C_{INU,INU} & C_{INU,KNU} & C_{INU,IDN} \\
  C_{KNU,INU} & C_{KNU,KNU} & C_{KNU,IDN} \\
  C_{IDN,INU} & C_{IDN,KNU} & C_{IDN,IDN}
\end{pmatrix}
\label{eq:mp-block-matrix}
\end{equation}

\textbf{Tractability:} 10,440 unique entries reduced to $\sim$500 via:
\begin{enumerate}[nosep]
\item Block structure (within-category dense, cross-category sparse)
\item Low-rank approximation: $C \approx U\Sigma V^T$ with $k \approx 10$--$20$ factors
\item Hierarchical Bayesian estimation pooling across similar contexts
\end{enumerate}

\subsection{Formal Coherence and Quality Indices}

\begin{equation}
K(s, \sigma) = 1 - \frac{\min\{d((s,\sigma), A), d((s,\sigma), B)\}}{d_{max}}
\label{eq:mp-k-formal}
\end{equation}

\begin{equation}
Q(C^*) = \Pi(C^*; \Psi) \quad \text{(payoff at equilibrium)}
\label{eq:mp-q-formal}
\end{equation}

\textbf{Realized welfare:} $W = Q(C^*) \cdot g(K)$ where $g(1) = 1$, $g(0) = 0$.

\subsection{What Mathematics Provides vs. Requires}

\begin{center}
\renewcommand{\arraystretch}{1.2}
\begin{tabular}{@{}ll@{}}
\toprule
\textbf{Mathematics Provides} & \textbf{Calibration Required} \\
\midrule
Structure of $\Pi(s,\sigma;\Psi)$ & Values of $\gamma, \alpha(\Psi), \beta(\Psi)$ \\
Form of 144-component aggregation & Discount factors $\delta_d$, loss aversion $\lambda_d$ \\
Fixed-point conditions for $C^*$ & 10,440 entries of $C$ matrix \\
Measures $K$ and $Q$ & Context-specific parameter estimates \\
\bottomrule
\end{tabular}
\end{center}

% =============================================================================
\section{Context as Endogenous Variable (Chapter 9)}
\label{sec:mp-context}
% =============================================================================

Chapter~9 answers the 4th CORE question: \textbf{WHEN and how does context modulate behavior?}
This is foundational because every parameter in the optimization problem is context-dependent.

\subsection{The 8 Ψ-Dimensions}

\begin{equation}
\boxed{
\boldsymbol{\Psi} = (\Psi_I, \Psi_S, \Psi_C, \Psi_K, \Psi_E, \Psi_T, \Psi_M, \Psi_F) \in [0,1]^8
}
\label{eq:mp-8psi}
\end{equation}

\begin{center}
\renewcommand{\arraystretch}{1.2}
\small
\begin{tabular}{@{}llll@{}}
\toprule
\textbf{Dimension} & \textbf{Full Name} & \textbf{Examples} & \textbf{Speed $\eta$} \\
\midrule
$\Psi_I$ & Institutional & Laws, property rights & Very slow \\
$\Psi_S$ & Social & Norms, expectations & Medium \\
$\Psi_C$ & Cognitive & Bounded rationality, attention & Fast \\
$\Psi_K$ & Cultural & Values, worldviews & Very slow \\
$\Psi_E$ & Economic & Markets, prices, income & Fast \\
$\Psi_T$ & Temporal & Time pressure, horizons & Variable \\
$\Psi_M$ & Material & Technology, infrastructure & Medium \\
$\Psi_F$ & Physical & Geography, climate & Very slow \\
\bottomrule
\end{tabular}
\end{center}

\subsection{Endogenous Context Dynamics}

Context is \textbf{not exogenous background}---it evolves through aggregate behavior:

\begin{equation}
\boxed{
\Psi_{t+1} = f(\Psi_t, a_t) = \Psi_t + \eta \cdot (a_t - a^*(\Psi_t))
}
\label{eq:mp-context-dynamics}
\end{equation}

\textbf{The Feedback Loop:}
\[
\Psi_t \to C^*(\Psi_t) \to a_t^*(C^*) \to \Psi_{t+1} \to \text{[cycle]}
\]

If behavior matches expectation ($a_t = a^*$), context persists.
If behavior deviates ($a_t \neq a^*$), context shifts.

\subsection{Multi-Speed Context}

Different context components have different adjustment speeds:
\begin{align}
\eta_{tech} &\gg \eta_{norm} > \eta_{inst} \gg \eta_{culture} \\
\text{(years)} &\quad \text{(decades)} \quad \text{(generations)}
\end{align}

This creates \textbf{cultural lag}: technology moves fast while institutions and culture lag behind.

\subsection{Tipping Points and Regime Change}

\textbf{Self-reinforcing feedback} ($\partial^2\Psi/\partial\Psi \partial a > 0$): Multiple equilibria, path dependence.

\textbf{S-curve dynamics:}
\begin{equation}
\Psi_{t+1} = \Psi_t + \eta \cdot \Psi_t (1 - \Psi_t) \cdot \Delta a_t
\label{eq:mp-s-curve}
\end{equation}

Slow start → rapid middle → saturation. Examples: smartphone adoption (2007-2015), smoking decline (1965-2020).

\subsection{Universal Modulation}

\textbf{Every parameter in EBF is context-dependent:}
\begin{itemize}[nosep]
\item Complementarity: $\gamma_{ij}(\Psi)$
\item Loss aversion: $\lambda(\Psi)$
\item Action threshold: $\theta(\Psi)$
\item Discount factor: $\delta(\Psi)$
\item Reference structure: $C^*(\Psi)$
\end{itemize}

This explains why ``one size fits all'' policies fail and treatment effects vary systematically.

\subsection{Implications for the Optimization Problem}

\begin{enumerate}[nosep]
\item \textbf{Moving target:} Optimal $I^*_t$ depends on $\Psi_t$ which evolves.
\item \textbf{Indirect effects:} Interventions change $\Psi$, which changes future optima.
\item \textbf{Path dependence:} Early interventions shape context for later ones.
\item \textbf{Tipping opportunities:} Small interventions near tipping points have large effects.
\end{enumerate}

% =============================================================================
\section{Welfare and FEPSDE (Chapter 10)}
\label{sec:mp-fepsde}
% =============================================================================

Chapter~10 answers two fundamental CORE questions: \textbf{WHAT constitutes utility?} (FEPSDE dimensions) and \textbf{WHO has utility?} (Level hierarchy). Without these, complementarity (HOW) and context (WHEN) have nothing to operate on.

\subsection{The K vs Q Distinction}

\textbf{Critical insight:} Coherence ($K$) measures whether a system is in equilibrium. It does \emph{not} measure whether the equilibrium is \textbf{good}. A society of mutual exploitation can have $K = 1$ yet be deeply impoverished in welfare.

\begin{center}
\renewcommand{\arraystretch}{1.2}
\begin{tabular}{@{}lll@{}}
\toprule
\textbf{Index} & \textbf{Measures} & \textbf{High Value Means} \\
\midrule
$K$ (Coherence) & At equilibrium? & Beliefs match reality \\
$Q$ (Quality) & Is equilibrium good? & High welfare for participants \\
\bottomrule
\end{tabular}
\end{center}

\textbf{Example:} Two societies with $K \approx 1$:
\begin{itemize}[nosep]
\item \textit{Extractive Coherence:} Elites extract, masses don't invest, institutions enforce extraction. Coherent but low welfare.
\item \textit{Inclusive Coherence:} Property rights protect, people invest, institutions enforce contracts. Coherent and high welfare.
\end{itemize}

\subsection{The Six FEPSDE Dimensions}

\begin{equation}
\boxed{
d \in \text{FEPSDE} = \{F, E, P, S, D, Eco\}
}
\label{eq:mp-fepsde-dims}
\end{equation}

\begin{center}
\renewcommand{\arraystretch}{1.2}
\small
\begin{tabular}{@{}clll@{}}
\toprule
\textbf{Dim} & \textbf{Full Name} & \textbf{Covers} & \textbf{Key Insight} \\
\midrule
$U^F$ & Financial & Income, wealth, security & Necessary not sufficient \\
$U^E$ & Emotional & Life satisfaction, meaning & Distinct from $U^F$ (Nordic paradox) \\
$U^P$ & Physical & Health, longevity, vitality & Precondition for others \\
$U^S$ & Social & Relationships, trust, status & Strongest longevity predictor \\
$U^D$ & Digital & Connectivity, literacy, rights & Post-COVID fundamental \\
$U^{Eco}$ & Ecological & Environment, sustainability & Intergenerational equity \\
\bottomrule
\end{tabular}
\end{center}

\textbf{Selection criteria:} (1) Empirically distinct, (2) Universally relevant, (3) Measurable, (4) Policy-relevant, (5) Theoretically grounded.

\subsection{The Three Utility Categories}

\begin{equation}
\boxed{
i \in \{INU, KNU, IDN\}
}
\label{eq:mp-utility-categories}
\end{equation}

\begin{center}
\renewcommand{\arraystretch}{1.2}
\begin{tabular}{@{}cll@{}}
\toprule
\textbf{Category} & \textbf{Question} & \textbf{Example} \\
\midrule
INU & What do I want for \textbf{myself}? & My income, my health \\
KNU & What do I want for \textbf{us}? & Family welfare, team success \\
IDN & \textbf{Who} am I / want to be? & Being honest, being competent \\
\bottomrule
\end{tabular}
\end{center}

\textbf{KNU is level-indexed:}
\begin{equation}
KNU^L, \quad L \in \{1, 2, 3, ..., \infty\}
\label{eq:mp-knu-levels}
\end{equation}

\begin{itemize}[nosep]
\item $KNU^1$: Household
\item $KNU^2$: Organization/Community
\item $KNU^3$: Region/Nation
\item $KNU^\infty$: Humanity/all sentient beings
\end{itemize}

\textbf{Key insight:} People sacrifice INU for KNU (parents for children, soldiers for country) and both for IDN (refusing to lie even when profitable). These are \emph{distinct} categories, not reducible to each other.

\subsection{Time Structure and Discounting}

\begin{equation}
\boxed{
t \in \{0, 1, 2, 3\}
}
\label{eq:mp-time-horizons}
\end{equation}

\begin{center}
\renewcommand{\arraystretch}{1.2}
\begin{tabular}{@{}cll@{}}
\toprule
$t$ & \textbf{Horizon} & \textbf{Timeframe} \\
\midrule
0 & Instant & Right now \\
1 & Short-term & Days to weeks \\
2 & Medium-term & Months to years \\
3 & Long-term & Years to lifetime \\
\bottomrule
\end{tabular}
\end{center}

\textbf{Dimension-specific discounting:}
\begin{equation}
\text{Present value of } U_{d,t} = \delta_d(\Psi)^t \cdot U_{d,t}
\label{eq:mp-discounting}
\end{equation}

Discount rates $\delta_d$ may vary by dimension: Financial (standard exponential), Health (lower---``health is priceless''), Ecological (intergenerational debates).

\subsection{Loss Aversion Structure}

Following \citet{kahnemantversky1979}, losses loom larger than gains:

\begin{equation}
\boxed{
U_d = G_d - \lambda_d(\Psi) \cdot P_d, \quad \lambda_d > 1
}
\label{eq:mp-loss-aversion}
\end{equation}

\begin{center}
\renewcommand{\arraystretch}{1.2}
\begin{tabular}{@{}llc@{}}
\toprule
\textbf{Dimension} & \textbf{Evidence} & $\lambda_d$ \\
\midrule
Financial & Kahneman \& Tversky (1979) & $\approx 2$ \\
Physical & Health state valuations & $\approx 2$--$3$ \\
Social & Rejection sensitivity research & $\approx 2$--$4$ \\
\bottomrule
\end{tabular}
\end{center}

\textbf{Policy implication:} Preventing harm ($-\Delta P$) creates more welfare than equivalent benefit ($+\Delta G$).

\subsection{The Complete 144-Component Structure}

\begin{equation}
\boxed{
3 \times 2 \times 6 \times 4 = 144 \text{ utility components}
}
\label{eq:mp-144-components}
\end{equation}

\begin{center}
\renewcommand{\arraystretch}{1.2}
\begin{tabular}{@{}rl@{}}
\toprule
\textbf{Factor} & \textbf{Elements} \\
\midrule
3 categories & INU, KNU, IDN \\
2 valences & Gain ($G$), Pain ($P$) \\
6 dimensions & F, E, P, S, D, Eco \\
4 time horizons & $t \in \{0, 1, 2, 3\}$ \\
\bottomrule
\end{tabular}
\end{center}

\textbf{Component notation:} $G_{i,d,t}$ or $P_{i,d,t}$, e.g., $G_{INU,F,0}$ = ``my financial gain now'' (bonus payment), $P_{KNU,Eco,3}$ = ``our collective ecological pain long-term'' (climate change).

\subsection{The Aggregate Quality Function}

\begin{equation}
\boxed{
Q(C^*) = \sum_{i \in \{INU, KNU, IDN\}} \sum_{d \in FEPSDE} \sum_{t=0}^{3}
         \omega_i \cdot \delta_d(\Psi)^t \left[ G_{i,d,t} - \lambda_d(\Psi) \cdot P_{i,d,t} \right]
}
\label{eq:mp-quality-function}
\end{equation}

where:
\begin{itemize}[nosep]
\item $\omega_i$: Weight on utility category $i$
\item $\delta_d(\Psi)$: Context-dependent discount factor for dimension $d$
\item $\lambda_d(\Psi)$: Context-dependent loss aversion for dimension $d$
\end{itemize}

\subsection{Implications for the Optimization Problem}

\begin{enumerate}[nosep]
\item \textbf{Objective enrichment:} $P_{eff}$ now optimizes over 144-dimensional $Q$, not scalar utility.
\item \textbf{Parameter complexity:} Need $\omega_i$ (3), $\delta_d$ (6), $\lambda_d$ (6) = 15 additional parameters.
\item \textbf{Tradeoff analysis:} Interventions may improve some $(i, d, t)$ while worsening others.
\item \textbf{K-Q separation:} High $K$ does not guarantee high $Q$---both must be optimized.
\item \textbf{Level aggregation:} $KNU^L$ requires aggregation weights $\alpha^L(\Psi)$ across levels.
\end{enumerate}

% =============================================================================
\section{Awareness: The Transformation Condition (Chapter 11)}
\label{sec:mp-awareness}
% =============================================================================

Chapter~11 answers the 6th CORE question: \textbf{AWARE---How aware am I of which utility at the Moment of Truth?} Awareness is a \emph{necessary condition} for utility to influence behavior.

\subsection{The Transformation Principle}

Potential utility (from Chapter 10) must pass through an awareness filter:

\begin{equation}
\boxed{
U_{eff,i,d,t} = A_{i,d,t}(t^*) \times U_{pot,i,d,t}
}
\label{eq:mp-transformation}
\end{equation}

where:
\begin{itemize}[nosep]
\item $A_{i,d,t}(t^*) \in [0,1]$: Awareness of component $(i, d, t)$ at Moment of Truth $t^*$
\item $U_{pot}$: Potential utility (144 components from Chapter 10)
\item $U_{eff}$: Effective utility (what actually influences behavior)
\end{itemize}

\textbf{Critical insight:} Without awareness ($A = 0$), even high potential utility has no behavioral effect. Anna choosing a promotion focuses on salary ($A(U^F) = 0.95$) but ignores health impact ($A(U^P) = 0.05$).

\subsection{The Moment of Truth}

\begin{equation}
t^* = \text{decision point where awareness matters}
\label{eq:mp-moment-truth}
\end{equation}

Awareness can vary dramatically before and after $t^*$. Post-decision regret (``I knew this would happen'') indicates $A(t^*) \neq A(t > t^*)$.

\subsection{Three System Levels}

\begin{center}
\renewcommand{\arraystretch}{1.2}
\begin{tabular}{@{}clll@{}}
\toprule
\textbf{Level} & \textbf{Name} & \textbf{Speed} & \textbf{Awareness Type} \\
\midrule
SYS 0 & Instinct & Instantaneous & Hardwired (pain, fear) \\
SYS 1 & Intuition & Fast & Automatic, heuristic \\
SYS 2 & Deliberation & Slow & Effortful, analytical \\
\bottomrule
\end{tabular}
\end{center}

\textbf{Instant Exception Rule:} $A(U^{P}_{t=0}) = 1$ always (immediate pain overrides everything).

\subsection{The K-Optimal Formula (EMERGE)}

The number of active intervention dimensions emerges from context:

\begin{equation}
\boxed{
K^*(\Psi) = K_{base} + \delta_{budget} + \delta_{complexity} + \delta_{stage} + \delta_{crowding}
}
\label{eq:mp-k-optimal}
\end{equation}

\begin{center}
\renewcommand{\arraystretch}{1.2}
\small
\begin{tabular}{@{}llll@{}}
\toprule
\textbf{Factor} & \textbf{Condition} & \textbf{Adjustment} & \textbf{Source} \\
\midrule
$K_{base}$ & --- & 4 & Michie et al. (2013) \\
$\delta_{budget}$ & $B > 2 \times B_{ref}$ & $+1$ per doubling & Resource constraint \\
$\delta_{complexity}$ & High education & $+1$ & Miller (1956) \\
$\delta_{stage}$ & Pre-contemplation & $-2$ & Prochaska (1983) \\
$\delta_{stage}$ & Action & $+1$ & Prochaska (1983) \\
$\delta_{crowding}$ & $P(\text{crowding}) > 0.3$ & $-1$ & Frey (1997) \\
\bottomrule
\end{tabular}
\end{center}

\textbf{Bounds:} $K^* \in [2, 9]$ based on cognitive capacity limits.

\subsection{The EMERGE Algorithm}

Interventions \emph{emerge} from context optimization, not menu selection:

\begin{equation}
\boxed{
\vec{I}^* = \text{EMERGE}(\Psi, \Gamma(\Psi), B, K^*) = \arg\max_{\|\vec{I}\|_0 \leq K^*} E(\vec{I} | \Psi)
}
\label{eq:mp-emerge}
\end{equation}

\textbf{Five-Step Pipeline:}
\begin{enumerate}[nosep]
\item Context Encoding: $\Psi \in [0,1]^8$
\item Complementarity Matrix: $\gamma_{ij}(\Psi) = \gamma_0 + \sum_k \beta_k \Psi_k$
\item Portfolio Size: Compute $K^*(\Psi)$
\item EMERGE: $\vec{I}^* = \arg\max E(\vec{I} | \Psi)$
\item Mechanism Mapping: Dimensions $\to$ concrete actions
\end{enumerate}

\subsection{Axiom and Corollary Counts}

\begin{center}
\renewcommand{\arraystretch}{1.2}
\begin{tabular}{@{}lr@{}}
\toprule
\textbf{Component} & \textbf{Count} \\
\midrule
Axioms (AWX-A1 to AWX-A90) & 104 \\
Bias Corollaries (derived) & 59 \\
Ψ-Modulations (PSI-01 to PSI-08) & 8 \\
Application Cases (CASE-01 to CASE-09) & 9 \\
\bottomrule
\end{tabular}
\end{center}

The 59 cognitive biases are \emph{derived} as corollaries from the axiom system, not added ad-hoc. This provides a unified explanation for diverse phenomena (confirmation bias, availability heuristic, etc.).

\subsection{Implications for the Optimization Problem}

\begin{enumerate}[nosep]
\item \textbf{Utility filtering:} Objective function uses $U_{eff}$, not $U_{pot}$---awareness creates the behavioral signal.
\item \textbf{Intervention target:} Many interventions aim to increase $A(\cdot)$ for underweighted dimensions.
\item \textbf{Portfolio optimization:} $K^*$ constrains the number of simultaneous interventions.
\item \textbf{Emergence principle:} $\vec{I}^*$ computed, not selected---context determines optimal profile.
\item \textbf{Moment sensitivity:} Timing matters---same $A$ profile at different $t$ yields different behavior.
\end{enumerate}

% =============================================================================
\section{Willingness to Act (Chapter 12)}
\label{sec:mp-willingness}
% =============================================================================

Chapter~12 answers the 7th CORE question: \textbf{READY---How ready am I to act on what I'm aware of?} Awareness (AWARE) is necessary but not sufficient; action requires willingness to exceed a threshold.

\subsection{The Action Condition}

The entire EBF framework culminates in one comparison:

\begin{equation}
\boxed{
\text{Action} \Leftrightarrow WAX \geq \theta
}
\label{eq:mp-action-condition}
\end{equation}

\textbf{The knowing-doing gap:} Thomas knows exercise would improve his health, feels his doctor recommends it, yet hasn't exercised in months. He is \emph{aware} but not \emph{willing}---his $WAX < \theta$.

\subsection{The Thinking Style Parameter ($\tau$)}

\begin{equation}
\boxed{
\tau \in [0, 1], \quad \tau = \tau_0 + \Delta\tau(\Psi)
}
\label{eq:mp-tau}
\end{equation}

\begin{center}
\renewcommand{\arraystretch}{1.2}
\begin{tabular}{@{}clp{6cm}@{}}
\toprule
$\tau$ & \textbf{Thinking Style} & \textbf{Logic} \\
\midrule
$\to 1$ & Connected (Complementary) & ``Only act if EVERYTHING is right'' \\
$\to 0$ & Separated (Additive) & ``Look at total benefit'' \\
\bottomrule
\end{tabular}
\end{center}

\textbf{Context shifts $\tau$:} High uncertainty ($\Psi_{UNS}$) $\to \tau \downarrow$; strong group identity ($\Psi_{SEG}$) $\to \tau \uparrow$; salient identity ($\Psi_{IDN}$) $\to \tau \uparrow$.

\subsection{The WAX Master Equation}

\begin{equation}
\boxed{
WAX = \tau \cdot \min_i\{U^*_i\} + (1-\tau) \cdot \sum_i U^*_i
}
\label{eq:mp-wax}
\end{equation}

where $U^*_i = A_i \cdot U_i$ is effective utility (from AWARE transformation).

\textbf{Interpretation:}
\begin{itemize}[nosep]
\item $\min\{U^*\}$ = weakest dimension determines value (chain logic)
\item $\sum U^*$ = all dimensions contribute (sum logic)
\item $\tau$ weights between these two logics
\end{itemize}

\subsection{Threshold Decomposition}

\begin{equation}
\boxed{
\theta = \theta_{base}(d) + \sum_j \delta_j \cdot \Psi_j
}
\label{eq:mp-theta}
\end{equation}

\begin{center}
\renewcommand{\arraystretch}{1.2}
\small
\begin{tabular}{@{}lc@{}}
\toprule
\textbf{Decision Type} & $\theta_{base}$ \\
\midrule
Continue habit & 0.2 \\
Small adjustment & 0.5 \\
Moderate change & 0.8 \\
Major investment & 1.2 \\
Life-changing & 1.5+ \\
\bottomrule
\end{tabular}
\end{center}

\textbf{Context adjustments:} High uncertainty $\to \theta \uparrow$; good subsidies $\to \theta \downarrow$; social proof $\to \theta \downarrow$.

\subsection{The Gap and Three Regimes}

\begin{equation}
\boxed{
\Delta = WAX - \theta
}
\label{eq:mp-gap}
\end{equation}

\begin{center}
\renewcommand{\arraystretch}{1.2}
\begin{tabular}{@{}cccc@{}}
\toprule
\textbf{Regime} & \textbf{Condition} & \textbf{P(Action)} & \textbf{Meaning} \\
\midrule
Below & $WAX < \theta$ & $< 50\%$ & Not ready \\
Critical & $WAX \approx \theta$ & $\approx 50\%$ & Tipping point \\
Above & $WAX > \theta$ & $> 50\%$ & Ready to act \\
\bottomrule
\end{tabular}
\end{center}

\subsection{Two Paths to Close the Gap}

\begin{center}
\renewcommand{\arraystretch}{1.2}
\begin{tabular}{@{}lll@{}}
\toprule
\textbf{Path} & \textbf{Strategy} & \textbf{Levers} \\
\midrule
1: $\uparrow WAX$ & ``Make it more attractive'' & $\uparrow \tau$, $\uparrow A$, $\uparrow U$ \\
2: $\downarrow \theta$ & ``Make it easier'' & $\downarrow$ friction, change defaults \\
\bottomrule
\end{tabular}
\end{center}

\textbf{Best practice:} Combine both paths for maximum effect.

\subsection{Implications for the Optimization Problem}

\begin{enumerate}[nosep]
\item \textbf{Final gate:} Action requires $WAX > \theta$---the last condition in the behavioral chain.
\item \textbf{Dual intervention:} Can target $WAX$ (motivation) or $\theta$ (barriers) or both.
\item \textbf{$\tau$ determines portfolio:} High $\tau$ requires addressing all dimensions; low $\tau$ allows compensation.
\item \textbf{Context sensitivity:} Both $\tau$ and $\theta$ are context-dependent---interventions can shift the context.
\item \textbf{99+ axioms:} Complete specification in Appendix UNMAPPED_AV (CORE-READY).
\end{enumerate}

% =============================================================================
\section{Behavioral Change Journey (Chapter 13)}
\label{sec:mp-bcj}
% =============================================================================

Chapter~13 answers the 8th CORE question: \textbf{STAGE---Where in the Behavioral Change Journey does the decision-maker find themselves?} Stage determines which interventions are effective.

\subsection{The Six Stages}

\begin{equation}
\boxed{
S \in \{S_1, S_2, S_3, S_4, S_5, S_6\}
}
\label{eq:mp-bcj-stages}
\end{equation}

\begin{center}
\renewcommand{\arraystretch}{1.2}
\small
\begin{tabular}{@{}clccc@{}}
\toprule
\textbf{Stage} & \textbf{Name} & $AWX$ & $\theta$ & \textbf{Intervention Focus} \\
\midrule
$S_1$ & Unaware & $\approx 0$ & $\to \infty$ & Create awareness \\
$S_2$ & Aware & low & high & Establish relevance \\
$S_3$ & Considering & medium & medium & Concretize benefits \\
$S_4$ & Intending & high & medium & Lower threshold \\
$S_5$ & Acting & high & low & Eliminate friction \\
$S_6$ & Maintaining & -- & $\to 0$ & Reinforce habit \\
\bottomrule
\end{tabular}
\end{center}

\textbf{Key insight:} Maria (``I know I should exercise'') and Thomas (``I've planned to start Monday'') are at different stages. An intervention that works for Thomas may fail for Maria.

\subsection{Stage-Dependent Parameter Modulation}

\begin{equation}
\boxed{
\theta^{eff}(S) = \theta_0 \cdot m(S), \quad m(S_1) \to \infty, \quad m(S_6) \to 0
}
\label{eq:mp-theta-modulation}
\end{equation}

\subsection{Stage Factor in EMERGE}

The stage factor $f_S$ modulates optimal portfolio size $K^*$:

\begin{center}
\renewcommand{\arraystretch}{1.2}
\begin{tabular}{@{}lcl@{}}
\toprule
\textbf{Stage} & $f_S$ & \textbf{Rationale} \\
\midrule
$S_1$ (Unaware) & 0.6 & Focus on awareness only \\
$S_2$ (Aware) & 0.7 & Build personal relevance \\
$S_3$ (Considering) & 1.0 & Address barriers \\
$S_4$ (Intending) & \textbf{1.2} & Maximum---close intention-behavior gap \\
$S_5$ (Acting) & 1.0 & Minimize friction \\
$S_6$ (Maintaining) & 0.8 & Habit lock-in requires less \\
\bottomrule
\end{tabular}
\end{center}

\subsection{Transition Dynamics}

\begin{equation}
P(S_i \to S_{i+1}) = f(AWX, WAX, \theta, \Psi, \text{Intervention})
\label{eq:mp-transition}
\end{equation}

\textbf{Forward transitions:}
\begin{itemize}[nosep]
\item $S_1 \to S_2$: $AWX$ crosses awareness threshold
\item $S_3 \to S_4$: Intention forms ($WAX \to \theta$)
\item $S_4 \to S_5$: $WAX > \theta$ triggers action
\item $S_5 \to S_6$: Repeated action until automaticity
\end{itemize}

\textbf{Backward transitions (relapse):} $P(S_i \to S_{i-1}) = g(\Delta\Psi, \text{time}, \text{alternatives})$

\subsection{Habit Lock-In Time ($T_{lock}$)}

\textbf{Novel prediction:} When $\theta \to 0$ at $S_6$, behavior becomes ``locked in''---minimal psychological cost, behavioral inertia resisting change.

\subsection{Implications for the Optimization Problem}

\begin{enumerate}[nosep]
\item \textbf{Stage-dependent effectiveness:} Same intervention has different effects at different stages.
\item \textbf{$f_S$ in $K^*$:} $K^* = K_{base} \cdot f_B + \delta_C + \delta_S + \delta_\gamma + \delta_L$ (EXC-3 hybrid) uses stage factor.
\item \textbf{Transition optimization:} Can optimize for stage transitions, not just final behavior.
\item \textbf{Relapse prevention:} Must consider backward transition risk in maintenance.
\item \textbf{85 axioms, 152 boundary conditions:} Complete specification in Appendix UNMAPPED_AW (CORE-STAGE).
\end{enumerate}

% =============================================================================
\section{Behavioral Change Segments: Population Heterogeneity (Chapter 14)}
\label{sec:mp-bcs}
% =============================================================================

Chapter~14 addresses population heterogeneity: \textbf{BCS---Which segment does the decision-maker belong to, and how does this modulate intervention effectiveness?} Segments emerge naturally from preference distributions.

\subsection{The Heterogeneity Space}

Population heterogeneity is defined over utility weights, awareness, and willingness:

\begin{equation}
\boxed{
\mathcal{H} = \{(\omega_i, A_i, W_i) : i \in \text{Population}\}
}
\label{eq:mp-heterogeneity-space}
\end{equation}

where:
\begin{itemize}[nosep]
\item $\omega_i = (\omega^F_i, \omega^E_i, \omega^P_i, \omega^S_i, \omega^D_i, \omega^{Eco}_i)$: FEPSDE utility weights
\item $A_i$: Individual awareness profile
\item $W_i$: Individual willingness profile
\end{itemize}

\subsection{Natural Clustering (NOT Fixed at 5)}

\textbf{Critical insight:} The number of segments is context-dependent, not universal:

\begin{equation}
\boxed{
k^* = \arg\min_k \{ BIC(k) \} = f(\mathcal{H}, \Psi)
}
\label{eq:mp-optimal-k}
\end{equation}

\begin{center}
\renewcommand{\arraystretch}{1.2}
\begin{tabular}{@{}lcc@{}}
\toprule
\textbf{Context Type} & \textbf{Typical $k^*$} & \textbf{Example} \\
\midrule
Novel + Voluntary + Social & 5 & Electric vehicles \\
Mandatory + High-Stakes & 2--3 & COVID vaccination \\
Established + Ubiquitous & 2--3 & Smartphone usage \\
\bottomrule
\end{tabular}
\end{center}

\textbf{Rogers 5-segment structure emerges when:}
\begin{itemize}[nosep]
\item Behavior is new/unfamiliar (high awareness heterogeneity)
\item Adoption is voluntary (individual choice)
\item Social influence present (peer effects matter)
\item Population large ($N \geq 1000$)
\end{itemize}

\subsection{Segment-Multiplier: Heterogeneous Treatment Effects}

\begin{equation}
\boxed{
\sigma_s = \frac{CATE(s)}{ATE} = \frac{E[Y(1) - Y(0) | S = s]}{E[Y(1) - Y(0)]}
}
\label{eq:mp-segment-multiplier}
\end{equation}

\begin{center}
\renewcommand{\arraystretch}{1.2}
\begin{tabular}{@{}lcp{6cm}@{}}
\toprule
\textbf{Multiplier} & \textbf{Value} & \textbf{Interpretation} \\
\midrule
$\sigma_s > 1.3$ & High & Segment responds above average \\
$\sigma_s \approx 1$ & Neutral & Segment responds like average \\
$\sigma_s < 0.5$ & Low & Segment underresponds \\
$\sigma_s < 0$ & Backfire & Intervention harms this segment \\
\bottomrule
\end{tabular}
\end{center}

\subsection{BCJ $\times$ BCS Interaction Matrix}

Journey stage (BCJ) and segment type (BCS) jointly determine intervention design:

\begin{equation}
\boxed{
\text{Intervention}^* = f(S_{BCJ}, C_{BCS})
}
\label{eq:mp-bcj-bcs}
\end{equation}

\begin{center}
\small
\renewcommand{\arraystretch}{1.2}
\begin{tabular}{@{}llp{5cm}@{}}
\toprule
\textbf{BCJ Stage} & \textbf{BCS Segment} & \textbf{Tailored Intervention} \\
\midrule
$S_3$ (Considering) & Incentive-driven & Concrete ROI calculation \\
$S_3$ (Considering) & Social-driven & ``Your neighbors are considering it'' \\
$S_4$ (Intending) & Identity-driven & One-click enrollment \\
$S_4$ (Intending) & Risk-averse & Money-back guarantee \\
$S_5$ (Acting) & Effort-averse & Make continuation the default \\
\bottomrule
\end{tabular}
\end{center}

\subsection{Leverage Profile}

Each segment has characteristic leverage points:

\begin{equation}
\boxed{
\ell_k^d = \frac{\text{Var}(\omega_i^d | i \in C_k)}{\sum_{d'} \text{Var}(\omega_i^{d'} | i \in C_k)}
}
\label{eq:mp-leverage}
\end{equation}

High $\ell_k^d$ = targeting dimension $d$ creates maximum impact in segment $k$.

\subsection{Cluster Size Diagnostics}

Cluster size distribution reveals heterogeneity structure:

\begin{equation}
CV = \frac{\sigma_{n_k}}{\bar{n}_k}, \quad R = \frac{|C_{max}|}{|C_{min}|}
\label{eq:mp-cluster-diagnostics}
\end{equation}

\begin{center}
\renewcommand{\arraystretch}{1.2}
\small
\begin{tabular}{@{}llll@{}}
\toprule
\textbf{CV} & \textbf{R} & \textbf{Diagnosis} & \textbf{Distribution} \\
\midrule
$< 0.4$ & $< 10$ & Well-balanced & Normal \\
$0.4$--$0.7$ & $10$--$50$ & Slightly skewed & Near-normal \\
$0.7$--$1.2$ & $50$--$200$ & Moderately skewed & Right-skewed \\
$> 1.2$ & $> 200$ & Highly polarized & Bimodal \\
\bottomrule
\end{tabular}
\end{center}

\subsection{Implications for the Optimization Problem}

\begin{enumerate}[nosep]
\item \textbf{Heterogeneous $V$:} Different $V(S, \Psi)$ for each segment $s$.
\item \textbf{Portfolio allocation:} Optimize budget across segment-specific interventions.
\item \textbf{$\sigma_s$ in $P_{eff}$:} $P_{eff,s} = \sigma_s \cdot P_{eff,avg}$.
\item \textbf{Cascade effects:} Early adopter segments trigger social norm shifts.
\item \textbf{13 axioms (UNMAPPED_SEG-1 to UNMAPPED_SEG-13):} Complete specification in Appendix BA (FORMAL-SEGMENT).
\end{enumerate}

% =============================================================================
\section{Decision Hierarchy and WEC Coherence (Chapter 15)}
\label{sec:mp-wec}
% =============================================================================

Chapter~15 answers: \textbf{How do we organize intervention decisions across multiple levels to ensure coherence?} The WEC framework integrates Willingness (W), Expectation (E), and Compatibility (C).

\subsection{The WEC Framework}

\begin{center}
\renewcommand{\arraystretch}{1.2}
\begin{tabular}{@{}clp{5cm}@{}}
\toprule
\textbf{Component} & \textbf{Question} & \textbf{Source} \\
\midrule
\textbf{W} (Willingness) & Can/will the person act? & Axioms W11-W13, AV \\
\textbf{E} (Expectation) & Will the intervention work? & Equilibrium theory, BB \\
\textbf{C} (Compatibility) & Do decision layers reinforce? & HI CORE-HIERARCHY \\
\bottomrule
\end{tabular}
\end{center}

\subsection{Four-Level Decision Hierarchy (HI)}

\begin{center}
\renewcommand{\arraystretch}{1.2}
\begin{tabular}{@{}clll@{}}
\toprule
\textbf{Level} & \textbf{Name} & \textbf{Nature} & \textbf{Strategic?} \\
\midrule
L0 & Meta-Meta & What counts as a decision? & Always \\
L1 & Strategic Roots & Few, high interaction & Always \\
L2 & Derived Strategic & Relational to L1 & Some \\
L3 & Operative & Low interaction & Rarely \\
\bottomrule
\end{tabular}
\end{center}

\subsection{The Universal L2 Formula (HI Theorem T1)}

\begin{equation}
\boxed{
N_{L2} = \frac{\alpha \cdot \gamma_{avg} \times n \times (1-m)}{\log(n)}
}
\label{eq:mp-n-l2}
\end{equation}

where:
\begin{itemize}[nosep]
\item $\alpha \approx 0.8$: Bounded rationality constant
\item $\gamma_{avg}$: Average complementarity
\item $n$: Total decision space size
\item $m$: Modularity affordance
\end{itemize}

\subsection{Context-Specific L2 Predictions}

\begin{center}
\renewcommand{\arraystretch}{1.2}
\begin{tabular}{@{}lcccc@{}}
\toprule
\textbf{Context} & $\gamma_{avg}$ & $n$ & $m$ & $N_{L2}$ \\
\midrule
Individual & 0.42 & 15 & 0.70 & 3.1 \\
Household & 0.58 & 19 & 0.60 & 3.8 \\
Organization & 0.62 & 35 & 0.50 & 6.2 \\
Nation & 0.68 & 45 & 0.35 & 7.8 \\
Tribal & 0.75 & 40 & 0.15 & 9.2 \\
\bottomrule
\end{tabular}
\end{center}

Empirical error: $< 5\%$ across all contexts.

\subsection{Dynamic Willingness (W11-W13)}

\textbf{W11:} Willingness is dynamically modulated:
\begin{equation}
WAX_i(t) = WAX_i^0 \times f(A(t)) \times g(KON(t)) \times h(JNY(t))
\label{eq:mp-wax-dynamic}
\end{equation}

\textbf{W12:} Boundary conditions exist: below $WAX_{blockade}$, no action regardless of other factors.

\textbf{W13:} Behavior probability requires all three:
\begin{equation}
P(\text{Action}|t) = P(WAX_t \geq WAX_{blockade}) \times P(\text{Opportunity}_t) \times P(\text{Capability}_t)
\label{eq:mp-action-probability}
\end{equation}

\subsection{Critical Foundations (CF1-CF5)}

\begin{enumerate}[nosep]
\item \textbf{Hierarchy Emerges:} Decision levels arise from complementarity, not authority.
\item \textbf{Modularity-Coherence Trade-off:} High autonomy risks drift; high coherence restricts innovation.
\item \textbf{Bounded Rationality Universal:} Humans coherently manage $\sim$4 L2s ($\alpha \approx 0.8$ invariant).
\item \textbf{Context Determines Scale:} Individual has 3-4 L2, tribal has 8-12 L2.
\item \textbf{Incoherence Costlier Than Complexity:} 50-70\% efficiency loss from missing L2s.
\end{enumerate}

\subsection{Implications for the Optimization Problem}

\begin{enumerate}[nosep]
\item \textbf{WEC coherence constraint:} All L2s must satisfy W, E, and C simultaneously.
\item \textbf{$N_{L2}$ is context-dependent:} Cannot assume fixed number of intervention layers.
\item \textbf{Incoherence penalty:} Missing L2 causes 50-70\% efficiency loss.
\item \textbf{Bounded rationality:} Optimal $K^*$ limited by cognitive capacity ($\alpha \approx 0.8$).
\item \textbf{12 axioms, 15 theorems:} Complete specification in Appendix UNMAPPED_HI (CORE-HIERARCHY).
\end{enumerate}

% =============================================================================
\section{Effective Probability: The Objective Function (Chapter 16)}
\label{sec:mp-peff}
% =============================================================================

Chapter~16 answers the ultimate question: \textbf{What is the probability that the person actually changes behavior?} This is the objective function $P_{eff}$ for the entire optimization problem.

\subsection{The Master Equation}

\begin{equation}
\boxed{
P_{eff} = \sigma\left( \underbrace{[WAX - \theta]}_{WEC} \cdot \alpha_{BCJ} \cdot \beta_{BCS} \right)
}
\label{eq:mp-peff-master}
\end{equation}

where:
\begin{itemize}[nosep]
\item $\sigma(x) = 1/(1 + e^{-kx})$: Sigmoid transformation
\item $WAX - \theta$: Willingness gap (from Chapter 12)
\item $\alpha_{BCJ} \in (0,1]$: Journey stage factor (from Chapter 13)
\item $\beta_{BCS} \in (0.5, 1.2)$: Segment type factor (from Chapter 14)
\end{itemize}

\subsection{Sigmoid Properties}

\begin{center}
\renewcommand{\arraystretch}{1.2}
\begin{tabular}{@{}rcl@{}}
\toprule
\textbf{Condition} & & \textbf{Interpretation} \\
\midrule
$WEC < 0$ & $\Rightarrow$ & $P_{eff} < 50\%$ (unlikely to act) \\
$WEC = 0$ & $\Rightarrow$ & $P_{eff} = 50\%$ (indifferent) \\
$WEC > 0$ & $\Rightarrow$ & $P_{eff} > 50\%$ (likely to act) \\
\bottomrule
\end{tabular}
\end{center}

\subsection{Natural Trajectory Model (NTM)}

Without intervention, probability decays:

\begin{equation}
A(t) = A_0 \cdot e^{-\lambda t}, \quad P_{eff}(t) \downarrow \text{ as } t \uparrow
\label{eq:mp-ntm-decay}
\end{equation}

\subsection{UNTCM: Unified NTM-Context Model}

The coupled ODE system from Appendix UNMAPPED_UN:

\begin{align}
\frac{dA}{dt} &= -\lambda(\Psi) \cdot A + \kappa_A \cdot f(B) + I_A(t) \\
\frac{dW}{dt} &= -\mu(\Psi) \cdot W + \kappa_W \cdot g(B) + I_W(t) \\
\frac{dB}{dt} &= \alpha \cdot P(A, W, \theta) - \delta \cdot B
\end{align}

\textbf{Three Regimes:}
\begin{itemize}[nosep]
\item \textbf{Regime I (Decay-Dominant):} $\kappa < \kappa^*$ --- fades without intervention
\item \textbf{Regime II (Feedback-Dominant):} $\kappa > \kappa^*$ --- self-sustaining lock-in
\item \textbf{Regime III (Oscillatory):} Complex $\kappa$ --- cycles between high/low
\end{itemize}

\textbf{Tipping Point:} $\bar{B}^* \approx 10$--$25\%$ determines regime transition.

\subsection{Loss Aversion Lock-in (X5)}

Status quo strengthens over time via loss aversion:

\begin{equation}
\theta_{eff}(t) = \theta_0 + \lambda_{LA} \cdot |S_{target} - r(t)|, \quad \lambda_{LA} \approx 2.25
\label{eq:mp-loss-aversion-lockin}
\end{equation}

Late intervention faces $\sim$225\% higher effective threshold.

\subsection{Implications for the Optimization Problem}

\begin{enumerate}[nosep]
\item \textbf{$P_{eff}$ is the objective:} $\max \sum_t \delta^t P_{eff}(S_t, I_t, \Psi_t)$.
\item \textbf{Sigmoid nonlinearity:} Small changes near $WEC = 0$ have large $\Delta P_{eff}$.
\item \textbf{Decay dynamics:} Without intervention, $P_{eff} \to 0$ (NTM).
\item \textbf{Tipping points:} Near $\bar{B}^*$, small interventions trigger regime change.
\item \textbf{Maintenance cost:} Status quo preservation costs $\sim$3.25$\times$ base (superlinear).
\item \textbf{10 axioms (P1-P10):} Complete specification in Appendix UNMAPPED_PBB (FORMAL-PROBABILITY).
\end{enumerate}

% =============================================================================
\section{Emergent Intervention Theory: The Instruments (Chapter 17)}
\label{sec:mp-eit}
% =============================================================================

Chapter~17 establishes: \textbf{How do we design interventions?} Interventions are not primitive categories but emerge from the continuous 10C space.

\subsection{Axiom A0: The Foundation}

All of EBF---and thus intervention theory---emerges from one axiom:

\begin{equation}
\boxed{
A0: \exists a, G, U, K: \text{Agent } a \text{ pursues goals } G \text{ under uncertainty } U \text{ in context } K
}
\label{eq:mp-axiom-a0}
\end{equation}

From A0, the 10C questions emerge \textit{necessarily}.

\subsection{The 10C Status Quo Vector}

Before intervening, diagnose current state across all dimensions:

\begin{equation}
\boxed{
\vec{S}_0 = (L, d, \gamma_{current}, \Psi, \Theta, A_0, W_0, \varphi_0, \text{Scope})
}
\label{eq:mp-status-quo}
\end{equation}

\subsection{Intervention as Transformation}

An intervention $\mathcal{I}$ transforms the status quo vector:

\begin{equation}
\boxed{
\mathcal{I}: \vec{S}_0 \mapsto \vec{S}_1 = \vec{S}_0 + \Delta\vec{S}(\mathcal{I})
}
\label{eq:mp-intervention-transform}
\end{equation}

\subsection{The Continuous Intervention Vector}

Interventions are points in continuous 10C-dimensional space:

\begin{equation}
\boxed{
\vec{I} = (I_{WHO}, I_{WHAT}, I_{HOW}, I_{WHEN}, I_{WHERE}, I_{AWARE}, I_{READY}, I_{STAGE}, I_{HIER}) \in [0,1]^9
}
\label{eq:mp-intervention-vector}
\end{equation}

\textbf{Key insight:} Discrete ``types'' (T1-T8) are projections onto single axes. 90\% of real interventions lie \textit{between} types.

\subsection{Emergence Functions (EIT-6, EIT-7, EIT-8)}

Which components should be emphasized?

\begin{center}
\renewcommand{\arraystretch}{1.2}
\begin{tabular}{@{}lll@{}}
\toprule
\textbf{Axiom} & \textbf{Condition} & \textbf{Implication} \\
\midrule
EIT-6 & $A_0 > 0.55$ & $I_{AWARE}$ low (awareness not the problem) \\
EIT-7 & $W_0 < 0.25$ & $I_{READY}$ maximal (willingness is bottleneck) \\
EIT-8 & $\Psi_c$ changeable & $I_{WHEN}$ high (context intervention) \\
\bottomrule
\end{tabular}
\end{center}

\subsection{Temporal Continuity: Trajectories}

Interventions are trajectories, not point actions:

\begin{equation}
I(t): [t_0, t_1] \to [0,1]^9
\label{eq:mp-trajectory}
\end{equation}

\textbf{Co-evolution:} Agent and intervention develop together:
\begin{equation}
\frac{d\vec{S}}{dt} = f(\vec{S}(t), \vec{I}(t), \Psi(t))
\label{eq:mp-coevolution}
\end{equation}

\subsection{Implications for the Optimization Problem}

\begin{enumerate}[nosep]
\item \textbf{Continuous control:} $I_t \in [0,1]^9$ not discrete choice from menu.
\item \textbf{Context-specific:} Intervention emerges from $\vec{S}_0$, not from catalog.
\item \textbf{Trajectory optimization:} Optimize $I(t)$, not single $I$.
\item \textbf{Emergence functions:} EIT-6/7/8 determine emphasis.
\item \textbf{20-field schema:} Complete specification in Appendix TKT (METHOD-TOOLKIT).
\end{enumerate}

% =============================================================================
\section{Emergent Temporal Theory: Phase-Affinity (Chapter 18)}
\label{sec:mp-temporal}
% =============================================================================

Chapter~18 develops: \textbf{How does intervention effectiveness vary with journey position?} Two key emergences connect time to intervention design.

\subsection{Axiom EIT-4: Journey-Position Emergence}

Journey-position emerges from the 10C Status Quo Vector:

\begin{equation}
\boxed{
\varphi = h(\vec{S}_0) = h(A_0, W_0, \text{Scope}) \approx \sigma(A_0 \cdot W_0 \cdot f(\text{Scope}) - \theta)
}
\label{eq:mp-journey-emergence}
\end{equation}

\begin{center}
\renewcommand{\arraystretch}{1.2}
\begin{tabular}{@{}lcc@{}}
\toprule
\textbf{Journey Region} & \textbf{$\varphi$ Range} & \textbf{10C Conditions} \\
\midrule
Pre-Awareness & $[0, 0.2)$ & $A_0 < 0.3$, $W_0 < 0.2$ \\
Triggered & $[0.2, 0.4)$ & $A_0 > 0.5$, $W_0 < 0.4$ \\
Action & $[0.4, 0.6)$ & $A_0 > 0.7$, $W_0 > 0.6$ \\
Maintenance & $[0.6, 0.8)$ & $W_0 > 0.8$ \\
Stable & $[0.8, 1.0]$ & $\gamma_{habit} > 0.7$ \\
\bottomrule
\end{tabular}
\end{center}

\subsection{Axiom EIT-5: Effect-Horizon Emergence}

The temporal effect-horizon emerges from the intervention vector:

\begin{equation}
\boxed{
\tau(\vec{I}) = g(I_{HIER}, I_{WHO}, I_{HOW})
}
\label{eq:mp-horizon-emergence}
\end{equation}

\begin{center}
\renewcommand{\arraystretch}{1.2}
\begin{tabular}{@{}lcc@{}}
\toprule
\textbf{Horizon} & \textbf{$I_{HIER}$ Range} & \textbf{Timeframe} \\
\midrule
Instant & $< 0.2$ & Hours--Days \\
Short-term & $[0.2, 0.4)$ & Days--Weeks \\
Medium-term & $[0.4, 0.7)$ & Weeks--Months \\
Long-term & $> 0.7$ & Months--Years \\
\bottomrule
\end{tabular}
\end{center}

\subsection{Phase-Affinity as Continuous Field}

\begin{equation}
\boxed{
\alpha: [0,1]^9 \times \mathbb{R}^{10C} \to [0,1], \quad \alpha(\vec{I}, \vec{S}_0) \approx \alpha_{base}(\vec{I}) \cdot \alpha_{match}(\vec{I}, \varphi) \cdot \alpha_{context}(\Psi)
}
\label{eq:mp-phase-affinity}
\end{equation}

\textbf{Matching function:}
\begin{equation}
\alpha_{match}(\vec{I}, \varphi) = 1 - \beta \cdot d(\vec{I}, \vec{I}^*(\varphi))
\label{eq:mp-matching}
\end{equation}

where $\vec{I}^*(\varphi)$ is the optimal intervention vector for journey region $\varphi$.

\subsection{Phase-K* Mapping (EMERGE)}

Stage factor $f_S$ modulates optimal intervention count:

\begin{center}
\renewcommand{\arraystretch}{1.2}
\begin{tabular}{@{}lcc@{}}
\toprule
\textbf{Phase} & \textbf{$f_S$} & \textbf{Rationale} \\
\midrule
Pre-Aware ($\varphi < 0.3$) & 0.7 & Focused awareness needed \\
Triggered ($\varphi \in [0.3, 0.5]$) & 1.2 & Critical window, more interventions \\
Action ($\varphi \in [0.5, 0.8]$) & 1.0 & Standard $K^*$ \\
Maintenance ($\varphi > 0.8$) & 0.8 & Focused habit formation \\
\bottomrule
\end{tabular}
\end{center}

\subsection{Implications for the Optimization Problem}

\begin{enumerate}[nosep]
\item \textbf{Phase-dependent effectiveness:} Same $\vec{I}$ has different $\alpha$ in different $\varphi$ regions.
\item \textbf{Mismatch cost:} $\alpha_{match} < 0.3$ signals phase mismatch---intervention poorly timed.
\item \textbf{Optimal sequencing:} Trajectory $\vec{I}(t)$ must adapt as $\varphi(t)$ evolves.
\item \textbf{Interpolation:} Continuous $\alpha$ handles ``between-phase'' positions naturally.
\item \textbf{EIT-4, EIT-5:} Complete specification in Appendix TKT (METHOD-TOOLKIT).
\end{enumerate}

% =============================================================================
\section{Segment-Heterogeneity Theory (Chapter 19)}
\label{sec:mp-segment}
% =============================================================================

Chapter~19 develops: \textbf{How does intervention effectiveness vary across behavioral segments?} The key insight: ATE can be misleading because CATE varies dramatically across segments \citep{bryan2021behavioural, athey2017econometrics, vivalt2020problems}.

\subsection{The Three Heterogeneity Axioms}

Behavioral segments emerge from three fundamental parameters \citep{fehr1999, fischbacher2014exclusion}:

\begin{tcolorbox}[colback=blue!5!white, colframe=blue!75!black,
    title=\textbf{Axiom PAP1-1: Social Preference Heterogeneity}]
\begin{equation}
\boxed{
U_i(x_i, x_{-i}) = u(x_i) + \alpha_i \cdot v(x_{-i}), \quad \alpha_i \in [0, 1]
}
\label{eq:mp-social-heterogeneity}
\end{equation}
Segmentation: $\alpha < 0.2$ (Self-interested), $\alpha \in [0.2, 0.5]$ (Moderate), $\alpha > 0.5$ (Social-Oriented) \citep{fehr1999, charness2002understanding}.
\end{tcolorbox}

\begin{tcolorbox}[colback=blue!5!white, colframe=blue!75!black,
    title=\textbf{Axiom PAP1-2: Temporal Preference Heterogeneity \citep{laibson1997, frederick2002time}}]
\begin{equation}
\boxed{
U_i(c_0, ..., c_T) = u(c_0) + \beta_i \sum_{t=1}^{T} \delta^t u(c_t), \quad \beta_i \in (0, 1]
}
\label{eq:mp-temporal-heterogeneity}
\end{equation}
Segmentation: $\beta > 0.9$ (Time-consistent), $\beta \in [0.7, 0.9]$ (Mild PB), $\beta < 0.7$ (Strong PB) \citep{odonoghue1999doingdoing, gruber2001addiction}.
\end{tcolorbox}

\begin{tcolorbox}[colback=blue!5!white, colframe=blue!75!black,
    title=\textbf{Axiom PAP1-3: Autonomy-Reactance Heterogeneity \citep{deci1999meta, brehm1966theory}}]
\begin{equation}
\boxed{
\sigma_i(\mathcal{I}) = \sigma_{base} - \rho_i \cdot \text{Perceived\_Manipulation}(\mathcal{I}), \quad \rho_i \in [0, 1]
}
\label{eq:mp-reactance-heterogeneity}
\end{equation}
Segmentation: $\rho < 0.3$ (Low), $\rho \in [0.3, 0.6]$ (Moderate), $\rho > 0.6$ (High Reactance).
\end{tcolorbox}

\subsection{Segment-Multiplier $\sigma_s$}

The Segment-Multiplier scales base effectiveness for segment $s$ \citep{bryan2021behavioural, szaszi2018systematic}:

\begin{equation}
\boxed{
E(\mathcal{I} | s) = E_{base}(\mathcal{I}) \cdot \sigma_s(\mathcal{I}), \quad \sigma_s \in [-0.5, 2.0]
}
\label{eq:mp-segment-multiplier}
\end{equation}

\textbf{Estimation:}
\begin{equation}
\hat{\sigma}_s = \frac{\text{CATE}(s)}{\text{ATE}}
\label{eq:mp-sigma-estimation}
\end{equation}

\begin{center}
\renewcommand{\arraystretch}{1.2}
\begin{tabular}{@{}cll@{}}
\toprule
\textbf{$\sigma_s$} & \textbf{Label} & \textbf{Interpretation} \\
\midrule
$2.0$ & Optimal Match & Double effectiveness \\
$1.0$ & Average & Typical response \\
$0.0$ & No Effect & Intervention misses segment \\
$-0.3$ & Reactance & Backfire effect \\
\bottomrule
\end{tabular}
\end{center}

\subsection{Full Effectiveness Function}

The complete effectiveness integrates all dimensions:

\begin{equation}
\boxed{
E(\mathcal{I} | \varphi, s, L, \Psi) = E_{max} \cdot \alpha_{t,\varphi} \cdot \sigma_s \cdot \lambda_L \cdot \gamma(\Psi)
}
\label{eq:mp-full-effectiveness}
\end{equation}

\textbf{Component mapping:}
\begin{itemize}[nosep]
\item $\alpha_{t,\varphi}$: Phase-Affinity from Ch18
\item $\sigma_s$: Segment-Multiplier from Ch19
\item $\lambda_L$: Level factor from AAA (WHO)
\item $\gamma(\Psi)$: Context modulation from Ch9 (WHEN)
\end{itemize}

\subsection{Phase $\times$ Segment Interaction}

Phase and segment effects are not independent:

\begin{equation}
\boxed{
E(\mathcal{I} | \varphi, s, t) = E_{base}(\mathcal{I}) \cdot \alpha_{t,\varphi}(\mathcal{I}) \cdot \sigma_s(\mathcal{I}) \cdot \delta_{t,s}
}
\label{eq:mp-phase-segment}
\end{equation}

where $\delta_{t,s}$ is the Phase-Segment Interaction Coefficient:
\begin{itemize}[nosep]
\item $\delta_{t,s} > 1.0$: Synergy (e.g., Present-Biased $\times$ Action $= 1.4$)
\item $\delta_{t,s} = 1.0$: Additive (no interaction)
\item $\delta_{t,s} < 1.0$: Interference (e.g., Autonomy-Seeking $\times$ Maintenance $= 0.7$)
\end{itemize}

\subsection{Reactance-Mitigation Theorem}

For Autonomy-Oriented segments with $\rho > 0.5$:

\begin{equation}
\boxed{
\sigma_{with\_autonomy\_cue} = \sigma_{base} + 0.4 \cdot \rho
}
\label{eq:mp-reactance-mitigation}
\end{equation}

An explicit autonomy cue can transform negative $\sigma$ into positive.

\subsection{Choice Satisfaction Constraint}

Menu design affects satisfaction (Miller-Iyengar-Schwartz):

\begin{equation}
\boxed{
Z(n) = Z_{max} \times (1 - 0.03 \times \max(0, n-7))
}
\label{eq:mp-choice-satisfaction}
\end{equation}

\textbf{Implication:} Beyond 7 options, each additional option reduces satisfaction by $\approx 3\%$.

\subsection{Implications for the Optimization Problem}

\begin{enumerate}[nosep]
\item \textbf{Segment-dependent returns:} Same $\vec{I}$ has different $\sigma_s$ for different segments.
\item \textbf{Backfire risk:} For Autonomy-Oriented, explicit nudges may have $\sigma < 0$.
\item \textbf{Phase$\times$Segment planning:} Optimal $\vec{I}(t)$ must consider $\delta_{t,s}$ interactions.
\item \textbf{Menu constraint:} Choice architectures should limit options to $n \leq 7$.
\item \textbf{Targeting strategy:} Segment-First, Self-Selection Menu, or Adaptive approaches.
\end{enumerate}

% =============================================================================
\section{Portfolio Optimization Theory (Chapter 20)}
\label{sec:mp-portfolio}
% =============================================================================

Chapter~20 develops: \textbf{How do we combine interventions into coherent portfolios?} The key insight: Portfolio design is more important than individual intervention selection \citep{iyengar2000choice, benartzi2001naive, thaler2008}.

\subsection{The A1-A8 Heuristic Clusters}

The continuous intervention space $\vec{I} \in [0,1]^9$ clusters naturally into 8 heuristics \citep{szaszi2018systematic, johnson2012beyond}:

\begin{center}
\renewcommand{\arraystretch}{1.1}
\small
\begin{tabular}{@{}llll@{}}
\toprule
\textbf{Cluster} & \textbf{Name} & \textbf{Dominant Components} & \textbf{Effect} \\
\midrule
A1 & Awareness Boost & AWARE, WHERE & +5--10\% \\
A2 & Feedback Loop & READY, AWARE & +8--15\% \\
A3 & Choice Architecture & HOW, WHEN, WHERE & +6--12\% \\
A4 & Temporal Push & STAGE, WHEN & +7--14\% \\
A5 & Identity Framing & WHO, AWARE, WHAT & +10--18\% \\
A6 & Social Proof & HIERARCHY, WHERE & +8--16\% \\
A7 & Financial Incentive & WHAT, READY & +12--25\% \\
A8 & Commitment Device & WHEN, STAGE, READY & +9--17\% \\
\bottomrule
\end{tabular}
\end{center}

\subsection{The Seven Constraints C1-C7}

Portfolio optimization is subject to seven constraint classes:

\begin{equation}
\boxed{
\begin{aligned}
\max_{P} \quad & \sum_{s} \sum_{i \in P} w_s \cdot E[\Delta W_s(\vec{I}_i)] + \sum_{i<j} \gamma_{ij} \sqrt{E_i E_j} \\
\text{s.t.} \quad
& \text{C1: } \sum_i \text{cost}(I_i) \leq B_{total} \quad \text{(Budget)} \\
& \text{C2: } \gamma(I_i, I_j) \geq \gamma_{min} \; \forall i,j \quad \text{(Crowding)} \\
& \text{C3: } t(I_i) < t(I_j) \text{ if } I_i \to I_j \quad \text{(Timing)} \\
& \text{C4: } \sum \text{cap\_req} \leq \text{cap\_avail} \quad \text{(Capacity)} \\
& \text{C5: } \text{autonomy}(I_i) \geq \eta_{min} \quad \text{(Ethical)} \\
& \text{C6: } \sigma(I_i, S_j) \geq \sigma_{min} \quad \text{(Segment)} \\
& \text{C7: } |P| \leq K^* \quad \text{(K-Optimal)}
\end{aligned}
}
\label{eq:mp-portfolio-constraints}
\end{equation}

\subsection{K* Formula (EMERGE Algorithm)}

Optimal portfolio size emerges from five factors \citep{iyengar2000choice, schwartz2004paradox, scheibehenne2010can}, using \textbf{EXC-3 Hybrid Formulation}:

\begin{equation}
\boxed{
K^* = K_{base} \cdot f_B + \delta_C + \delta_S + \delta_\gamma + \delta_L
}
\label{eq:mp-k-optimal}
\end{equation}

\begin{itemize}[nosep]
\item $f_B = \sqrt{B/B_{ref}}$ --- Budget factor (\textbf{MULTIPLICATIVE}: $B=0 \Rightarrow K^*=0$)
\item $\delta_C = -0.2 \cdot K_{base} \cdot (c - c_{ref})$ --- Complexity adjustment (\textbf{ADDITIVE})
\item $\delta_S = K_{base} \cdot (\alpha_{BCJ} - 1)$ --- Stage adjustment (\textbf{ADDITIVE})
\item $\delta_\gamma = -K_{base} \cdot |\bar{\gamma}_{neg}|$ --- Crowding adjustment (\textbf{ADDITIVE})
\item $\delta_L = 0.1 \cdot K_{base} \cdot (L - 1)$ --- Scope adjustment (\textbf{ADDITIVE})
\end{itemize}

\subsection{Behavioral Change Prediction}

The complete prediction pipeline from 10C Status Quo to behavioral change \citep{prochaska1992stages, ajzen1991theory, webb2006meta}:

\begin{equation}
\boxed{
P(BC | \vec{S}_0, P) = \sigma\left( \beta_0 + \beta_1 \cdot E(P | \vec{S}_0) \right)
}
\label{eq:mp-bc-prediction}
\end{equation}

where conditional portfolio effectiveness is:

\begin{equation}
\boxed{
E(P | \vec{S}_0) = \sum_{i \in P} E_{max,i} \cdot \alpha_{t_i,\varphi_0} \cdot \sigma_{s(A_0,W_0)} \cdot \lambda_L \cdot \gamma(\Psi) + \sum_{i<j} \gamma_{ij} \sqrt{E_i E_j}
}
\label{eq:mp-conditional-effectiveness}
\end{equation}

\subsection{Six Heuristic Frameworks}

Six equally valid frameworks derive from the continuous model:

\begin{center}
\renewcommand{\arraystretch}{1.2}
\small
\begin{tabular}{@{}llll@{}}
\toprule
\textbf{Framework} & \textbf{Basis} & \textbf{When to Use} & \textbf{Error Risk} \\
\midrule
F1: Component & A1-A8 clusters & Sparse data, quick screening & Aggregation \\
F2: Behavioral & Problem type & Rich qualitative data & Context blindness \\
F3: Phase-Based & Journey stage & Temporal data available & Phase uncertainty \\
F4: Segment-First & Behavioral segment & Rich behavioral data & Classification error \\
F5: Outcome-First & FEPSDE priority & Multi-outcome portfolio & Cross-outcome conflict \\
F6: Context-Adaptive & $\Psi$ monitoring & Dynamic contexts & Misclassification \\
\bottomrule
\end{tabular}
\end{center}

\subsection{Hierarchical Optimization Levels}

Portfolio optimization operates at three levels:

\begin{center}
\renewcommand{\arraystretch}{1.2}
\begin{tabular}{@{}llll@{}}
\toprule
\textbf{Level} & \textbf{Timeframe} & \textbf{K* Multiplier} & \textbf{Typical Size} \\
\midrule
Strategic (L1) & Annual & $\times 0.8$ & 3--5 interventions \\
Tactical (L2) & Quarterly & $\times 1.0$ & 5--8 campaigns \\
Operative (L3) & Weekly & $\times 1.2$ & 8--12 measures \\
\bottomrule
\end{tabular}
\end{center}

\subsection{Implications for the Optimization Problem}

\begin{enumerate}[nosep]
\item \textbf{Complementarity matters:} $\gamma_{ij}$ determines whether combinations help or hurt.
\item \textbf{Crowding-out risks:} Social + Financial ($\gamma = -0.3$), Financial + Intrinsic ($\gamma = -0.3$).
\item \textbf{Timing constraints:} $I_{AWARE} \prec I_{WHAT} \prec I_{WHEN}$ sequencing required.
\item \textbf{K* is emergent:} Portfolio size is not free; emerges from context and constraints.
\item \textbf{Six frameworks:} Choice depends on data availability and population heterogeneity.
\item \textbf{Hierarchical cascade:} Strategic $\to$ Tactical $\to$ Operative with increasing K*.
\end{enumerate}

% =============================================================================
\section{Dynamic Portfolio Evolution (Chapter 21)}
\label{sec:mp-dynamic}
% =============================================================================

Chapter~21 develops: \textbf{As status quo $S_t$ evolves, when should we redesign portfolio $P_t$?} Static optimization (Ch20) becomes dynamic when phase, context, or parameters change \citep{prochaska1992stages, norcross2011stages, milkman2011harnessing}.

\subsection{Three Evolution Mechanisms}

Status quo $S_t$ evolves through three mechanisms \citep{ajzen1991theory, webb2006meta}:

\begin{enumerate}[nosep]
\item \textbf{Phase Progression:} $\varphi_t$ changes as individuals progress through BCJ stages
\item \textbf{Context Shifts:} $\Psi_t$ changes due to policy, economic, or social shocks
\item \textbf{Bayesian Learning:} $\Theta_t$ updates as we observe intervention outcomes
\end{enumerate}

\subsection{Dynamic K*(t) Formula}

Optimal portfolio size evolves over time:

\begin{equation}
\boxed{
K^*(t) = K^*_0 \cdot g(\varphi_t) \cdot h(\Psi_t) \cdot \ell(\Theta_t)
}
\label{eq:mp-dynamic-k}
\end{equation}

\begin{itemize}[nosep]
\item $g(\varphi_t) = 1 + 0.3 \cdot |\varphi_t - \varphi_{t-1}|$ --- Phase progression factor
\item $h(\Psi_t) = 1 + 0.5 \cdot \mathbb{1}[\text{shock}]$ --- Context shift factor
\item $\ell(\Theta_t) = 1 - 0.2 \cdot \text{converged}$ --- Learning convergence factor
\end{itemize}

\subsection{Phase Velocity ODE}

Phase progression rate depends on current awareness and willingness:

\begin{equation}
\boxed{
\frac{d\varphi}{dt} = v_0 + v_A \cdot A_t + v_W \cdot W_t
}
\label{eq:mp-phase-velocity}
\end{equation}

\subsection{Redesign Triggers}

Portfolio redesign is triggered when:

\begin{itemize}[nosep]
\item $|\Delta\varphi| > 0.2$: Phase jump (individual moved to new BCJ stage)
\item $||\Delta\Psi|| > 0.15$: Context shock (policy/economic/social shift)
\item $||\Delta\Theta|| < \epsilon$ for $T$ periods: Convergence (learning stabilized)
\end{itemize}

\subsection{Convergence Criteria}

Learning has converged when parameters stabilize:

\begin{equation}
\boxed{
||\Theta_t - \Theta_{t-1}|| < \epsilon \quad \text{for } T \text{ consecutive periods}
}
\label{eq:mp-convergence}
\end{equation}

Practical values: $\epsilon = 0.05$ (5\% threshold), $T = 2$ periods.

\subsection{K* Evolution Pattern}

All domains show the same pattern:

\begin{center}
\renewcommand{\arraystretch}{1.2}
\begin{tabular}{@{}lll@{}}
\toprule
\textbf{Phase} & \textbf{K* Range} & \textbf{Rationale} \\
\midrule
Baseline & 2--3 & Focused initial portfolio \\
Transition/Shock & 4--5 & More interventions during adaptation \\
Convergence & 3 & Refocused maintenance portfolio \\
\bottomrule
\end{tabular}
\end{center}

\subsection{Implications for the Optimization Problem}

\begin{enumerate}[nosep]
\item \textbf{Time-varying K*:} Portfolio size is not static; evolves with $(\varphi_t, \Psi_t, \Theta_t)$.
\item \textbf{Redesign triggers:} Phase jumps, context shocks, and convergence signal portfolio updates.
\item \textbf{Learning loop:} Design $\to$ Implement $\to$ Measure $\to$ Update $\Theta$ $\to$ Redesign.
\item \textbf{Convergence stopping:} When $\Delta\Theta < \epsilon$, stop redesigning and maintain.
\item \textbf{Crisis response:} Context shocks require qualitative portfolio changes, not just intensity.
\end{enumerate}

% =============================================================================
\section{Policy Applications Theory (Chapter 22)}
\label{sec:mp-policy-applications}
% =============================================================================

Policy portfolio design follows the same principles as individual-level design, but with
\textbf{scale-up modifications} for population-level effects \citep{sunstein2014nudge, benartzi2017should, thaler2008}.

\subsection{Policy Portfolio Definition}

A \textbf{policy portfolio} is a set of coordinated interventions designed to change
population-level behavior in a domain \citep{halpern2015inside, service2014think}, subject to coherence constraints and cross-domain
spillover effects.

\begin{equation}
\boxed{
P_{\text{domain}} = \{I_1, I_2, \ldots, I_{K^*}\} \quad \text{where } |P_{\text{domain}}| \leq K^*_{\text{policy}}
}
\end{equation}

\subsection{Three Scale-Up Modifications}

Portfolio design at the policy level requires three modifications from individual-level design:

\begin{enumerate}[nosep]
\item \textbf{Segment Heterogeneity at Scale:} Populations include diverse segments (Present-Biased, Social-Oriented, Autonomy-Seeking). Effective policies target each segment differently.

\item \textbf{Implementation Dynamics:} Policy changes constrained by administrative capacity, political feasibility, and path-dependency. Sequencing matters.

\item \textbf{Cross-Domain Spillovers:} Policies in one domain (Health) affect outcomes in others (Labor, Pensions). Coherence requires managing spillovers.
\end{enumerate}

\subsection{Policy-Level K* Formula}

The optimal portfolio size at policy level (L0/L1 scope):

\begin{equation}
\boxed{
K^*_{\text{policy}} = K_{\text{base}} \cdot f_B + \delta_C + \delta_S + \delta_\gamma + \delta_L
}
\end{equation}

where $f_L$ is the \textbf{scope-level adjustment factor}:

\begin{center}
\begin{tabular}{@{}llc@{}}
\toprule
\textbf{Scope Level} & \textbf{Description} & \textbf{$f_L$} \\
\midrule
L0 (Systemic) & Few, deep interventions & 0.6 → $K^* \approx 3$--4 \\
L1 (Strategic) & Policy bundles per domain & 0.8 → $K^* \approx 4$--6 \\
\bottomrule
\end{tabular}
\end{center}

\subsection{Policy Coherence Score}

For multi-domain policy, compute the \textbf{Policy Coherence Score}:

\begin{equation}
\boxed{
\text{Coherence} = \frac{\sum_{i < j} \gamma_{ij}}{\binom{n}{2}}
}
\end{equation}

where $n$ is the number of policy domains and $\gamma_{ij}$ is the cross-domain spillover.

\textbf{Interpretation:}
\begin{itemize}[nosep]
\item $\text{Coherence} > 0.3$: Policies well-coordinated (reinforce each other)
\item $\text{Coherence} \in [0.1, 0.3]$: Policies neutral (don't interfere much)
\item $\text{Coherence} < 0.1$: Policies poorly coordinated (may interfere)
\end{itemize}

\subsection{Cross-Domain Spillover Matrix}

The spillover effect between domains follows the cross-domain complementarity:

\begin{equation}
\gamma_{\text{cross}}(D_i, D_j) = f(\text{Congruence}, \text{Capacity}, \text{Behavioral Coherence})
\end{equation}

Example empirical values:
\begin{center}
\begin{tabular}{@{}llc@{}}
\toprule
\textbf{Domain Pair} & \textbf{Effect} & \textbf{$\gamma$} \\
\midrule
Health ↔ Environment & Aligned messages & +0.4 \\
Pension ↔ Labor & Work-life interaction & +0.6 \\
Health ↔ Pension & Conflicting messages & +0.1 \\
Environment ↔ Labor & Green jobs effect & +0.1 \\
\bottomrule
\end{tabular}
\end{center}

\subsection{Crisis Adjustment Factor}

Context shocks require portfolio adjustment via the \textbf{crisis adjustment factor}:

\begin{equation}
\boxed{
E_{\text{crisis}} = E_{\text{base}} \cdot \kappa(\Psi) \quad \text{where } \kappa(\Psi) \in [0.4, 1.0]
}
\end{equation}

During crisis:
\begin{itemize}[nosep]
\item Information campaigns (AWARE): $\kappa \approx 0.4$ (60\% decline)
\item Social comparison (WHAT$_S$): $\kappa \approx 0.4$ (60\% decline)
\item Financial incentives (WHAT$_F$): $\kappa \approx 1.3$ (increase!)
\item Choice architecture (WHEN): $\kappa \approx 0.2$ (80\% decline, backlash risk)
\end{itemize}

\subsection{Domain-Specific K* Patterns}

Empirical $K^*$ emergence across four major policy domains:

\begin{center}
\begin{tabular}{@{}llcc@{}}
\toprule
\textbf{Domain} & \textbf{Context} & \textbf{$K^*$} & \textbf{Primary 10C} \\
\midrule
Health (Prevention) & PreAware population & 4 & AWARE, READY \\
Pensions (Reform) & Trust-eroding context & 3 & WHAT$_S$, STAGE \\
Environmental & Crisis-prone context & 4--5 & WHAT$_F$, WHO \\
Labor Markets & Disruption context & 5 & STAGE, WHAT$_S$ \\
\bottomrule
\end{tabular}
\end{center}

\subsection{Policy Constraints (Extended)}

Building on C1--C7 (Chapter 20), policy-level adds:

\begin{itemize}[nosep]
\item \textbf{C6 (Segment):} $\sigma_s \geq 0$ for all major segments (no backlash at 25\%+ of population)
\item \textbf{C7 (K-Optimal):} $|P_{\text{domain}}| \leq K^*_{\text{policy}}$ (focus over scatter)
\item \textbf{C2 (Crowding) Extended:} $\gamma_{ij} \geq -0.1$ between policy domains (cross-domain coherence)
\end{itemize}

\subsection{Axioms: Policy Applications (MP-A119 to MP-A125)}

\begin{tcolorbox}[colback=red!3!white, colframe=red!50!black,
    title=\textbf{Axioms: Policy Applications Theory (Chapter 22)}]

\textbf{MP-A119 (Policy Portfolio Definition):} A policy portfolio $P_{\text{domain}}$ is a set of coordinated interventions designed to change population-level behavior, subject to coherence constraints and cross-domain spillovers.

\textbf{MP-A120 (Scale-Up Principle):} Policy-level design requires three modifications from individual-level: (a) segment heterogeneity at scale, (b) implementation dynamics, (c) cross-domain spillovers.

\textbf{MP-A121 (Policy-Level K*):} Optimal portfolio size at policy level: $K^*_{\text{policy}} = K_{\text{base}} \cdot f_B + \delta_C + \delta_S + \delta_\gamma + \delta_L$, where $f_L \in [0.6, 0.8]$ for L0/L1 scope.

\textbf{MP-A122 (Scope-Level Factor):} Policy scope determines $f_L$: L0 (Systemic) → $f_L = 0.6$, L1 (Strategic) → $f_L = 0.8$. Higher scope levels require fewer, deeper interventions.

\textbf{MP-A123 (Policy Coherence):} Cross-domain coherence measured by: $\text{Coherence} = \sum_{i<j} \gamma_{ij} / \binom{n}{2}$. Target: Coherence $> 0.25$ (acceptable) or $> 0.35$ (good).

\textbf{MP-A124 (Crisis Adjustment):} Context shocks modulate effectiveness via $\kappa(\Psi) \in [0.4, 1.0]$. Different intervention types respond differently: WHAT$_F$ increases, WHEN decreases.

\textbf{MP-A125 (Implementation Sequence):} Policy implementation requires sequencing constraints: phase-appropriate interventions first, high-spillover domains coordinated, administrative capacity respected.

\end{tcolorbox}

\subsection{Theorems: Policy Applications (MP-T50 to MP-T53)}

\begin{theorem}[MP-T50: Policy Portfolio Effectiveness]
The effectiveness of a policy portfolio across a heterogeneous population is:
\begin{equation}
E(P_{\text{domain}}) = \sum_{i} E_i \cdot \alpha_i \cdot \bar{\sigma} + \sum_{i<j} \gamma_{ij} \sqrt{E_i \cdot E_j}
\end{equation}
where $\bar{\sigma} = \sum_s w_s \cdot \sigma_s$ is the weighted segment multiplier.
\end{theorem}

\begin{theorem}[MP-T51: Cross-Domain Coherence Optimization]
Given $n$ policy domains with portfolios $P_1, \ldots, P_n$, the coherence-optimal configuration maximizes:
\begin{equation}
\max_{P_1, \ldots, P_n} \sum_{d} E(P_d) + \lambda \cdot \text{Coherence}(P_1, \ldots, P_n)
\end{equation}
subject to C1--C7 within each domain and $\gamma_{ij} \geq -0.1$ across domains.
\end{theorem}

\begin{theorem}[MP-T52: Crisis Response Theorem]
When context shock occurs ($||\Delta\Psi|| > 0.15$), optimal policy response is:
\begin{enumerate}[nosep]
\item Expand WHAT$_F$ (financial) interventions: $\kappa(\Psi) > 1$
\item Contract WHEN (choice architecture): $\kappa(\Psi) < 0.3$
\item Maintain AWARE at reduced scope: $\kappa(\Psi) \approx 0.4$
\item Reframe WHO (identity) toward crisis-relevant messaging
\end{enumerate}
Portfolio returns to baseline as $\Psi$ stabilizes.
\end{theorem}

\begin{theorem}[MP-T53: Domain-Specific K* Emergence]
The optimal $K^*$ emerges domain-specifically:
\begin{equation}
K^*_d = K_{\text{base}} \cdot f_B(d) + \delta_C(d) + \delta_S(d) + \delta_\gamma(d) + \delta_L(d) \text{ (EXC-3)}
\end{equation}
Empirical patterns: Health $K^* \approx 4$, Pensions $K^* \approx 3$, Environment $K^* \approx 4$--5, Labor $K^* \approx 5$.
\end{theorem}

\subsection{Implications for the Optimization Problem}

\begin{enumerate}[nosep]
\item \textbf{Multi-domain objective:} Maximize $\sum_d E(P_d)$ subject to cross-domain coherence.
\item \textbf{Scope-level K*:} Policy portfolios have lower K* (3--6) than individual interventions (5--8).
\item \textbf{Crisis dynamics:} Context shocks require rapid portfolio recomposition, not just intensity adjustment.
\item \textbf{Spillover management:} Cross-domain $\gamma$ must be monitored and optimized.
\item \textbf{Implementation sequencing:} Administrative capacity constraints binding at policy level.
\end{enumerate}

% =============================================================================
\section{Multi-Level Implementation Theory (Chapter 23)}
\label{sec:mp-multilevel}
% =============================================================================

The 9C + Portfolio framework is \textbf{scale-invariant}---the same diagnostic and
optimization logic applies at all organizational levels \citep{ostrom1990governing, acemoglu2015state, north1990institutionsinstitutions}.

\subsection{Scale Invariance Principle}

The core insight: Whether designing interventions for a single person or entire nations,
the same structure applies \citep{simon1962architecture, williamson2000new}:

\begin{equation}
\boxed{
\text{Diagnose } S_0 \to \text{Design } \vec{I} \to \text{Modulate } (\alpha, \sigma) \to \text{Optimize } P \to \text{Execute \& Learn}
}
\end{equation}

What changes: (1) Unit of analysis, (2) Parameters, (3) Constraints, (4) Spillovers.
The structure stays invariant; the specifics vary.

\subsection{Five Organizational Levels}

The optimization problem applies at five organizational levels:

\begin{equation}
L_0 \text{ (Individual)} \to L_1 \text{ (Household)} \to L_2 \text{ (Organization)} \to L_3 \text{ (National)} \to L_4 \text{ (Global)}
\end{equation}

\begin{center}
\begin{tabular}{@{}clcc@{}}
\toprule
\textbf{Level} & \textbf{Unit} & \textbf{$K^*$} & \textbf{Key Challenge} \\
\midrule
L0 & Individual & 2--4 & Within-person heterogeneity \\
L1 & Household & 3--5 & Phase mismatch between partners \\
L2 & Organization & 5--8 & Hierarchical coordination \\
L3 & National & 4--6 & Regional Ψ heterogeneity \\
L4 & Global & 3--5 & No enforcement mechanism \\
\bottomrule
\end{tabular}
\end{center}

\subsection{K* by Organizational Level}

Optimal portfolio size varies systematically by level:

\begin{equation}
\boxed{
K^*(L) = K_{\text{base}} \cdot f_{\text{level}}(L)
}
\end{equation}

where $f_{\text{level}}$:

\begin{center}
\begin{tabular}{@{}lcc@{}}
\toprule
\textbf{Level} & \textbf{$f_{\text{level}}$} & \textbf{Rationale} \\
\midrule
L0 (Individual) & 0.5--0.7 & Focused personal portfolio \\
L1 (Household) & 0.6--0.8 & Coordination between partners \\
L2 (Organization) & 1.0--1.3 & Hierarchical abstimmung \\
L3 (National) & 0.8--1.0 & Policy-level (fewer, deeper) \\
L4 (Global) & 0.6--0.8 & International coordination \\
\bottomrule
\end{tabular}
\end{center}

\subsection{Cross-Level Constraint}

Interventions at different levels must not undermine each other:

\begin{equation}
\boxed{
\gamma(L_i, L_j) \geq 0 \quad \text{for all } i \neq j
}
\end{equation}

\textbf{Violation example:} CEO portfolio (``We empower you'') + Plant Manager portfolio
(``Cut costs or lose job'') yields $\gamma = -0.6$ (strong negative). Result: Implementation collapse.

\subsection{Cross-Level Spillovers}

\textbf{Top-Down Causation:} Higher-level policy affects lower-level behavior.

\begin{equation}
E_{L_i}(P) = E_{\text{base}}(P) \cdot \prod_{j > i} \kappa_{j \to i}(\Psi_j)
\end{equation}

where $\kappa_{j \to i}$ is the spillover factor from level $j$ to level $i$.

\textbf{Bottom-Up Pressure:} Lower-level non-compliance constrains higher-level policy.

\begin{equation}
\text{Compliance}_{L_j} = f\left(\sum_{i < j} \text{Adoption}_{L_i}\right)
\end{equation}

If individual adoption (L0) is low, national policy (L3) becomes politically infeasible.

\subsection{Cascade Dynamics}

Policy signals transform as they pass through organizational hierarchy:

\begin{equation}
\boxed{
\text{Message}_{L_i} = \text{Message}_{L_{i+1}} \cdot \tau_i + \epsilon_i
}
\end{equation}

where $\tau_i$ is the transmission fidelity and $\epsilon_i$ is local noise/interpretation.

\textbf{Cascade failure:} When $\tau_i < 0.5$, message can reverse direction at each level.

\subsection{Hybrid Optimality Theorem}

\begin{equation}
\boxed{
E_{\text{hybrid}} > \max(E_{\text{top-down}}, E_{\text{bottom-up}})
}
\end{equation}

The optimal implementation strategy combines:
\begin{itemize}[nosep]
\item \textbf{Top-down:} Goals, budgets, minimum standards
\item \textbf{Bottom-up:} Design and execute implementation
\end{itemize}

\subsection{Level Aggregation}

Micro behavior aggregates to macro outcomes via:

\begin{equation}
E_{\text{macro}}(L_j) = \sum_{i < j} w_i \cdot E_i \cdot \kappa_{i \to j}
\end{equation}

where $w_i$ is the weight of level $i$ in the aggregation.

\subsection{Axioms: Multi-Level Implementation (MP-A126 to MP-A132)}

\begin{tcolorbox}[colback=red!3!white, colframe=red!50!black,
    title=\textbf{Axioms: Multi-Level Implementation Theory (Chapter 23)}]

\textbf{MP-A126 (Scale Invariance):} The 9C diagnostic and portfolio optimization logic applies identically at all organizational levels (L0--L4). Structure invariant; parameters vary.

\textbf{MP-A127 (Five Levels):} Five organizational levels: L0 (Individual), L1 (Household), L2 (Organization), L3 (National), L4 (Global). Each has distinct coordination challenges.

\textbf{MP-A128 (K* by Level):} Optimal portfolio size varies by level: $K^*(L) = K_{\text{base}} \cdot f_{\text{level}}(L)$. L2 (Organization) has highest K*; L0/L4 have lowest.

\textbf{MP-A129 (Cross-Level Constraint):} $\gamma(L_i, L_j) \geq 0$ for all $i \neq j$. Interventions at different levels must not undermine each other.

\textbf{MP-A130 (Cross-Level Spillovers):} Top-down causation: higher levels affect lower. Bottom-up pressure: lower-level non-compliance constrains higher-level policy.

\textbf{MP-A131 (Cascade Dynamics):} Policy signals transform through hierarchy with transmission fidelity $\tau_i$. Cascade failure occurs when $\tau_i < 0.5$.

\textbf{MP-A132 (Hybrid Optimality):} Hybrid implementation (top-down goals + bottom-up execution) outperforms pure top-down or pure bottom-up strategies.

\end{tcolorbox}

\subsection{Theorems: Multi-Level Implementation (MP-T54 to MP-T57)}

\begin{theorem}[MP-T54: Scale Invariance]
For any organizational level $L \in \{L_0, L_1, L_2, L_3, L_4\}$, the optimization problem has the same structure:
\begin{equation}
\max_{P \in \mathcal{P}(L)} E(P|S_0, \Psi) \quad \text{s.t. } C_1\text{--}C_7
\end{equation}
Only the parameters $(S_0, \Psi, \Theta)$ and constraint bounds vary by level.
\end{theorem}

\begin{theorem}[MP-T55: Cross-Level Coherence]
A multi-level implementation succeeds iff:
\begin{equation}
\gamma(L_i, L_j) \geq 0 \quad \forall i, j \in \{0, 1, 2, 3, 4\}
\end{equation}
Violation at any level pair causes cascade failure with probability $p > 0.6$.
\end{theorem}

\begin{theorem}[MP-T56: Hybrid Dominance]
For implementation problems with heterogeneous local contexts:
\begin{equation}
E_{\text{hybrid}} \geq 1.2 \cdot \max(E_{\text{top-down}}, E_{\text{bottom-up}})
\end{equation}
The hybrid premium is highest when local Ψ variance is high.
\end{theorem}

\begin{theorem}[MP-T57: Cascade Fidelity]
For an $n$-level cascade, end-to-end message fidelity is:
\begin{equation}
\tau_{\text{total}} = \prod_{i=1}^{n-1} \tau_i
\end{equation}
If each $\tau_i = 0.8$, a 4-level cascade has $\tau_{\text{total}} = 0.41$ (message reversed).
\end{theorem}

\subsection{Implications for the Optimization Problem}

\begin{enumerate}[nosep]
\item \textbf{Level-specific K*:} Portfolio size varies from 2--4 (L0) to 5--8 (L2) to 3--5 (L4).
\item \textbf{Cross-level coherence:} Must verify $\gamma(L_i, L_j) \geq 0$ before multi-level launch.
\item \textbf{Hybrid implementation:} Combine top-down goals with bottom-up execution for +20\% effectiveness.
\item \textbf{Cascade monitoring:} Track message fidelity $\tau_i$ at each level; intervene if $\tau_i < 0.6$.
\item \textbf{Aggregation:} Micro behavior determines macro feasibility; model bottom-up constraints explicitly.
\end{enumerate}

% =============================================================================
\section{Emergent Life Journeys Theory (Chapter 24)}
\label{sec:mp-life-journeys}
% =============================================================================

The seven life journeys are NOT pre-defined categories but \textbf{EMERGENT outcomes}
of the 9C + Portfolio system when applied to different behavioral domains over decades \citep{elder1998children, bengtson2001families, hofstede2001cultures}.

\subsection{Parameter-Emergence Principle}

Each domain has distinct parameters $(\alpha, \beta, \lambda, \rho)$ that lead to
characteristic habituation curves, frequency evolutions, and segment transitions \citep{heckman2006skill, cunha2007technology}.

\begin{equation}
\boxed{
\frac{dA}{dt} = \alpha_{\text{domain}} \cdot T_1(t) \cdot (1 - A(t)) - \beta_{\text{domain}} \cdot A(t)
}
\end{equation}

\begin{equation}
\boxed{
\frac{d\phi}{dt} = g_\phi(W(t), A(t), \Psi(t), \{T_i(t)\})
}
\end{equation}

\subsection{Seven Emergent Journeys}

Domain-specific parameters generate seven characteristic pathways:

\begin{center}
\small
\begin{tabular}{@{}lccccl@{}}
\toprule
\textbf{Journey} & $\boldsymbol{\alpha}$ & $\boldsymbol{\beta}$ & $\boldsymbol{\lambda}$ & \textbf{Time to Eq.} & \textbf{Type} \\
\midrule
Health & 0.05 & 0.02 & 0.3 & 5--10y & Smooth \\
Money & 0.005 & 0.01 & 0.1 & 30y+ & Slow \\
Career & Episodic & 0.001 & 0.4 & 5--10y/role & Jumps \\
Relationships & 0.02 & 0.03 & 0.5 & 10--20y & Bifurcating \\
Self-Governance & 0.10 & 0.02 & 0.6 & 2--5y & Fast \\
Living Space & 0.01 & 0.02 & 0.2 & 15--30y & Discontinuous \\
Meaning/Legacy & 0.002 & 0.001 & 0.8 & 30y+ & Convex \\
\bottomrule
\end{tabular}
\end{center}

\subsection{Domain-Specific Convergence}

\textbf{Health (Fast):} $\phi^*_H \approx 0.85$ within 5--10 years. Maintenance-focused thereafter.

\textbf{Money (Slow):} $\phi^*_M \approx 0.60$--0.70 over 30+ years. Never fully automatic.

\textbf{Career (Episodic):} Jump discontinuities at transitions (age 22, 28, 35, etc.).

\subsection{Frequency Evolution Formula}

Optimal intervention frequency evolves with phase:

\begin{equation}
\boxed{
f_{\text{optimal}}(\phi) = \begin{cases}
1/\text{year} & \text{if } \phi < 0.3 \text{ (PreAware)} \\
1/\text{month} & \text{if } 0.3 < \phi < 0.6 \text{ (Triggered)} \\
1/\text{week} & \text{if } 0.6 < \phi < 0.85 \text{ (Action)} \\
1/\text{quarter} & \text{if } \phi > 0.85 \text{ (Maintenance)}
\end{cases}
}
\end{equation}

\subsection{Intergenerational Spillover Matrix}

Life journeys transmit across generations via $\Gamma^{inter}$ \citep{doepke2017parenting, chowdhury2022preferences, sacerdote2007nature}:

\begin{equation}
\boxed{
\phi^{child}_{d,0} = \phi_{\text{base}} + \sum_{d'=1}^{7} \gamma^{inter}_{d',d} \cdot \left(\phi^{parent}_{d',T} - \bar{\phi}_{d'}\right)
}
\end{equation}

\textbf{Key spillover coefficients:}
\begin{itemize}[nosep]
\item $\gamma^{inter}_{SG,SG} = 0.60$ (Self-Governance highest transmission)
\item $\gamma^{inter}_{C,C} = 0.55$ (Career strongly inherited)
\item $\gamma^{inter}_{M,M} = 0.50$ (Money moderately inherited)
\item $\gamma^{inter}_{H,H} = 0.45$ (Health transmitted via habits)
\end{itemize}

\subsection{Critical Transmission Windows (CTWs)}

Parental influence is strongest during domain-specific age windows \citep{heckman2008schools, cunha2010estimating, turkheimer2003socioeconomic}:

\begin{center}
\small
\begin{tabular}{@{}llc@{}}
\toprule
\textbf{Domain} & \textbf{CTW (Child Age)} & \textbf{Mechanism} \\
\midrule
Health & 0--6 years & Food habits, activity patterns \\
Self-Governance & 3--12 years & Executive function modeling \\
Money & 12--22 years & Financial socialization \\
Career & 14--22 years & Aspiration formation \\
Meaning & 15--25 years & Value transmission \\
\bottomrule
\end{tabular}
\end{center}

\subsection{Intergenerational Multiplier}

Self-Governance is the ``master domain'' for transmission:

\begin{equation}
\boxed{
E\left[\sum_{d} \Delta\phi^{child}_{d,0}\right] \approx 1.70 \cdot \Delta\phi^{parent}_{SG}
}
\end{equation}

Every unit increase in parental Self-Governance produces 1.70 units of aggregate child benefit.

\subsection{Family-Level ROI}

Family interventions during CTWs have multiplier effects:

\begin{equation}
\boxed{
ROI_{\text{family}} \approx 1.7 \times ROI_{\text{individual}}
}
\end{equation}

\subsection{Axioms: Emergent Life Journeys (MP-A133 to MP-A140)}

\begin{tcolorbox}[colback=red!3!white, colframe=red!50!black,
    title=\textbf{Axioms: Emergent Life Journeys Theory (Chapter 24)}]

\textbf{MP-A133 (Parameter-Emergence):} Life journeys emerge from domain-specific parameters $(\alpha, \beta, \lambda, \rho)$ applied to 9C ODEs over decades. NOT pre-defined categories.

\textbf{MP-A134 (Seven Journeys):} Seven characteristic pathways emerge: Health, Money, Career, Relationships, Self-Governance, Living Space, Meaning/Legacy.

\textbf{MP-A135 (Domain ODE):} Each domain follows: $dA/dt = \alpha_d \cdot T_1 \cdot (1-A) - \beta_d \cdot A$ with domain-specific parameters.

\textbf{MP-A136 (Convergence Rates):} Convergence time varies: Health (5--10y), Money (30y+), Career (episodic 5--10y/role).

\textbf{MP-A137 (Frequency Evolution):} Optimal intervention frequency $f_{\text{optimal}}(\phi)$ decreases as $\phi \to 1$ (maintenance).

\textbf{MP-A138 (Intergenerational Spillover):} $\phi^{child}_{d,0} = \phi_{\text{base}} + \sum \gamma^{inter}_{d',d} \cdot (\phi^{parent}_{d',T} - \bar{\phi}_{d'})$. Children inherit behavioral capital.

\textbf{MP-A139 (Critical Transmission Windows):} Each domain has CTW where parental influence is strongest (Health: 0--6y, Self-Gov: 3--12y, Money: 12--22y).

\textbf{MP-A140 (Self-Governance Master Domain):} $E[\sum \Delta\phi^{child}] \approx 1.70 \cdot \Delta\phi^{parent}_{SG}$. Self-Governance has highest cross-domain transmission.

\end{tcolorbox}

\subsection{Theorems: Emergent Life Journeys (MP-T58 to MP-T62)}

\begin{theorem}[MP-T58: Journey Emergence]
Given domain-specific parameters $(\alpha_d, \beta_d, \lambda_d, \rho_d)$ and 9C ODEs, exactly seven characteristic trajectory types emerge corresponding to the seven life domains.
\end{theorem}

\begin{theorem}[MP-T59: Domain Equilibrium]
For domain $d$ with parameters $(\alpha_d, \beta_d)$, the equilibrium awareness is:
\begin{equation}
A^*_d = \frac{\alpha_d \cdot T_1^* + \lambda_d \cdot E[\text{observed effect}]}{\alpha_d + \beta_d + \lambda_d}
\end{equation}
Health: $A^*_H \approx 0.75$. Money: $A^*_M \approx 0.35$.
\end{theorem}

\begin{theorem}[MP-T60: Intergenerational ROI]
For family interventions during CTW:
\begin{equation}
ROI_{\text{family}} = ROI_{\text{individual}} \times (1 + \gamma^{inter} \cdot \frac{T_{\text{child}}}{T_{\text{parent}}}) \approx 1.7 \times ROI_{\text{individual}}
\end{equation}
\end{theorem}

\begin{theorem}[MP-T61: Self-Governance Multiplier]
Self-Governance investments have cross-domain effects:
\begin{equation}
\frac{\partial \phi^{child}_d}{\partial \phi^{parent}_{SG}} > 0 \quad \forall d \in \{\text{7 domains}\}
\end{equation}
with aggregate multiplier $\approx 1.70$.
\end{theorem}

\begin{theorem}[MP-T62: Lifespan Optimization]
The optimal 60-year (3-generation) intervention strategy prioritizes:
\begin{enumerate}[nosep]
\item Self-Governance during CTW (ages 3--12)
\item Health during CTW (ages 0--6)
\item Money/Career during CTW (ages 12--22)
\end{enumerate}
to maximize $\sum_{\text{generations}} \sum_d \phi_d$.
\end{theorem}

\subsection{Implications for the Optimization Problem}

\begin{enumerate}[nosep]
\item \textbf{Domain-specific timing:} Intervention timing must match domain convergence rates.
\item \textbf{Lifespan perspective:} 30--60 year planning horizon, not single-period optimization.
\item \textbf{Intergenerational effects:} Family-level interventions during CTWs have 1.7× ROI.
\item \textbf{Self-Governance priority:} Highest cross-domain, cross-generational multiplier.
\item \textbf{Emergence, not taxonomy:} Seven journeys arise from parameters, not pre-definition.
\end{enumerate}

% =============================================================================
\section{Chapter Contributions Summary}
\label{sec:mp-chapters}
% =============================================================================

\begin{center}
\renewcommand{\arraystretch}{1.2}
\small
\begin{longtable}{clll}
\toprule
\textbf{Ch} & \textbf{Title} & \textbf{Contribution} & \textbf{Formula/Concept} \\
\midrule
\endhead
1 & Introduction & Motivation for framework & $C_{ij} \neq 0$ necessary \\
2 & Rationality & Stability gap, K vs Q & $C_{ij} = 0$ relaxed, QRE, ESS \\
3 & Limits & 6 anomaly domains, 4 interdependencies & $\beta$-$\delta$, $\lambda$, $\alpha$-$\beta$ \\
4 & Empirical & 5 facts, diagnostic, calibration vs. simulation & $\alpha$, $\beta$, scaling, functional knowledge \\
5 & Complementarity & HOW, supermodularity, 4 types & $\Gamma$ (15) + $\Lambda$ (28) + $\delta$ (120) = 163 \\
6 & Reference $C^*$ & Fixed point, 3 derivations, $C^*(\Psi)$ & $C^* = \Phi(C^*, \Psi)$, $k^n$ convergence \\
7 & Non-Concavity & Peaks/valleys, Roberts fit, K-Q trap & Valley crossing, calibration imperative \\
8 & Mathematical & Two-axis model, 144 components, $C$ matrix & $\Pi(s,\sigma;\Psi)$, block structure, tractability \\
\midrule
9 & Context (WHEN) & 8Ψ dimensions, endogenous, multi-speed & $\Psi_{t+1} = f(\Psi_t, a_t)$, tipping points, universal modulation \\
10 & FEPSDE (WHAT+WHO) & K vs Q, 3 categories, 6 dims & $Q(C^*) = \sum \omega_i \delta_d^t [G - \lambda P]$, 144 components \\
11 & Awareness (AWARE) & Transformation, $K^*$, EMERGE & $U_{eff} = A(t^*) \times U_{pot}$; 104 axioms, 59 biases \\
12 & Willingness (READY) & $\tau$, WAX, $\theta$, $\Delta$ gap & $WAX = \tau \min + (1-\tau) \sum$; 99 axioms \\
13 & BCJ (STAGE) & 6 stages, $f_S$, transitions, $T_{lock}$ & $\theta^{eff} = \theta_0 \cdot m(S)$; 85 axioms, 152 theorems \\
14 & BCS (Segments) & $\mathcal{H}$, $k^*$, $\sigma_s$, BCJ$\times$BCS, leverage & $\sigma_s = CATE/ATE$; 13 axioms \\
15 & WEC (Hierarchy) & L0-L3, $N_{L2}$, W11-W13, CF1-5 & $N_{L2} = \alpha \gamma n(1-m)/\log n$; 12 axioms, 15 theorems \\
16 & $P_{eff}$ (OBJECTIVE) & Master eq., sigmoid, NTM, UNTCM, X5 & $P_{eff} = \sigma(WEC \cdot \alpha_{BCJ} \cdot \beta_{BCS})$; 10 axioms \\
17 & EIT (Instruments) & A0, $\vec{S}_0$, $\vec{I}$, EIT-6/7/8, trajectories & $\vec{I} \in [0,1]^9$; 13 axioms \\
18 & Temporal (Phase) & EIT-4/5, $\varphi$, $\tau$, $\alpha(\vec{I},\vec{S}_0)$, $f_S$ & $\varphi = h(\vec{S}_0)$; $\alpha_{match}$ \\
\midrule
19 & Segment-Spec. & PAP1-1/2/3, $\sigma_s$, $\delta_{t,s}$, $Z(n)$, reactance & $\sigma_s = CATE/ATE$; $\delta_{Phase,Seg}$ \\
20 & Portfolios & A1-A8 clusters, C1-C7, $K^*$, 6 Frameworks, hierarchical & $K^* = K_b \cdot f_B + \delta_C + \delta_S + \delta_\gamma + \delta_L$; $P(BC)$ \\
21 & Dynamic & $K^*(t)$, phase velocity, redesign triggers, convergence & $K^*(t) = K^*_0 \cdot g \cdot h \cdot \ell$; $d\varphi/dt$ \\
\midrule
22 & Policy Applications & $K^*_{\text{policy}}$, Coherence, spillovers, $\kappa(\Psi)$ & $K^*_{\text{policy}}$ (EXC-3 hybrid); Coherence $= \sum \gamma / \binom{n}{2}$ \\
23 & Multi-Level & Scale invariance, L0--L4, $K^*(L)$, spillovers, cascade & $\gamma(L_i, L_j) \geq 0$; $\tau_{\text{total}} = \prod \tau_i$ \\
24 & Life Journeys & 7 emergent journeys, $\Gamma^{inter}$, CTWs, ROI$_{\text{family}}$ & $dA/dt = \alpha T_1(1-A) - \beta A$; $\gamma^{inter}_{SG} = 1.70$ \\
\bottomrule
\end{longtable}
\end{center}

% =============================================================================
\section{The 4-Stage Pipeline (Chapter 1b)}
\label{sec:mp-pipeline}
% =============================================================================

The transformation from potential utility to behavioral change follows 4 stages:

\begin{equation}
\boxed{
U_{pot} \xrightarrow{\text{AWARE}} U_{eff} \xrightarrow{\text{READY}} WAX \geq \theta \xrightarrow{\text{STAGE}} S(t)
}
\end{equation}

\begin{center}
\renewcommand{\arraystretch}{1.4}
\begin{tabular}{clll}
\toprule
\textbf{Stage} & \textbf{Name} & \textbf{Transformation} & \textbf{10C Involved} \\
\midrule
1 & Utility Definition & $U_{pot} = f(L, d, \gamma, \Psi, \Theta)$ & WHO, WHAT, HOW, WHEN, WHERE \\
2 & Awareness Filter & $U_{eff} = A(\cdot) \times U_{pot}$ & AWARE \\
3 & Action Threshold & Action $\Leftrightarrow WAX \geq \theta$ & READY \\
4 & Behavioral Change & $S(t+1) = f(S(t), A, WAX, \theta, \Psi)$ & STAGE \\
\bottomrule
\end{tabular}
\end{center}

\textbf{The Central Equation:}
\begin{equation}
\text{Action} \Leftrightarrow WAX\Big(A\big(U_{pot}(L, d, \gamma, \Psi, \Theta)\big), \tau\Big) \geq \theta(\Psi)
\end{equation}

% =============================================================================
\section{Bellman Formulation}
\label{sec:mp-bellman}
% =============================================================================

The dynamic optimization problem can be written in Bellman form:

\begin{equation}
\boxed{
V(S_t, \Psi_t) = \max_{I_t \in \mathcal{S}(\Psi_t, S_t)} \left\{ P_{eff}(S_t, I_t, \Psi_t) + \delta \cdot V(S_{t+1}, \Psi_{t+1}) \right\}
}
\label{eq:mp-bellman}
\end{equation}

subject to:
\begin{align}
S_{t+1} &= S_t + \Delta S(I_t) \\
\Psi_{t+1} &= f(\Psi_t, I_t, a_t)
\end{align}

\textbf{Complications:}
\begin{itemize}
\item Non-concavity: Multiple local optima (Chapter 7)
\item Endogenous state space: $\mathcal{S}(\Psi)$ changes with $\Psi$
\item High dimensionality: $S \in \mathbb{R}^{9+}$, $\Psi \in [0,1]^8$
\item Heterogeneity: Different $V$ for different segments (Chapter 14)
\end{itemize}

% =============================================================================
\section{Goal-Directed Optimization (Zielgerichtete Optimierung)}
\label{sec:mp-goal-directed}
% =============================================================================

The general formulation in Section \ref{sec:mp-bellman} maximizes cumulative $P_{eff}$.
In practice, behavioral interventions often have a \textbf{specific target}---a desired
behavioral state $S^*$ to be reached. This section formalizes goal-directed optimization.

\subsection{The Target State $S^*$}

Define the \textbf{target state} as a specific configuration in 10C space:
\begin{equation}
S^* = (L^*, d^*, \gamma^*, \Psi^*, \Theta^*, A^*, W^*, \varphi^*, N^*_{L2})
\label{eq:mp-target}
\end{equation}

\begin{example}[Retirement Savings Target]
\begin{itemize}
\item $A^* = 0.8$ (high awareness of retirement needs)
\item $W^* = 0.7$ (strong willingness to save)
\item $\varphi^* = S4$ (Action/Maintenance phase)
\item Behavioral target: ``Employee saves 10\% of salary automatically''
\end{itemize}
\end{example}

\subsection{Distance Metric}

Define a distance metric from current state to target:
\begin{equation}
d(S_t, S^*) = \sqrt{\sum_{k=1}^{9} w_k \cdot (S_{t,k} - S^*_k)^2}
\label{eq:mp-distance}
\end{equation}

where $w_k$ are dimension-specific weights reflecting importance.

\textbf{Alternative metrics:}
\begin{itemize}
\item \textbf{Weighted Euclidean}: $d_E = \sqrt{\sum w_k (S_k - S^*_k)^2}$
\item \textbf{Manhattan}: $d_M = \sum w_k |S_k - S^*_k|$ (dimension-wise progress)
\item \textbf{Mahalanobis}: $d_{Mah} = \sqrt{(S-S^*)^T \Sigma^{-1} (S-S^*)}$ (accounts for correlations)
\item \textbf{Behavioral}: $d_B = 1 - P(a = a^* | S)$ (probability of target behavior)
\end{itemize}

\subsection{Alternative Objective Formulations}

\subsubsection{Formulation 1: Target-Reaching (First-Hitting-Time)}

Minimize time to reach target:
\begin{equation}
\boxed{
\min_{\{I_t\}} T^* = \min\{t : S_t \in \mathcal{B}_\epsilon(S^*)\}
}
\label{eq:mp-first-hitting}
\end{equation}

subject to $I_t \in \mathcal{S}(\Psi_t, S_t)$, where $\mathcal{B}_\epsilon(S^*)$ is an $\epsilon$-ball around the target.

\subsubsection{Formulation 2: Distance Minimization}

Minimize terminal distance to target:
\begin{equation}
\boxed{
\min_{\{I_t\}_{t=0}^{T}} d(S_T, S^*)
}
\label{eq:mp-terminal}
\end{equation}

subject to $I_t \in \mathcal{S}(\Psi_t, S_t)$ for all $t$.

\subsubsection{Formulation 3: Path-Integrated Distance}

Minimize cumulative distance over trajectory:
\begin{equation}
\boxed{
\min_{\{I_t\}} \sum_{t=0}^{T} \delta^t \cdot d(S_t, S^*) + \lambda \cdot C(I_t)
}
\label{eq:mp-path-integrated}
\end{equation}

where $C(I_t)$ is intervention cost.

\subsubsection{Formulation 4: Probability of Target Achievement}

Maximize probability of reaching target within time horizon:
\begin{equation}
\boxed{
\max_{\{I_t\}} P\left(S_T \in \mathcal{B}_\epsilon(S^*)\right)
}
\label{eq:mp-prob-target}
\end{equation}

This connects to the original $P_{eff}$ formulation when $S^*$ corresponds to the target behavior.

\subsection{Reachability Analysis}

\begin{definition}[Reachable Set]
The \textbf{reachable set} from $S_0$ in $T$ steps is:
\begin{equation}
\mathcal{R}(S_0, T) = \{S_T : \exists \{I_t\}_{t=0}^{T-1} \text{ with } I_t \in \mathcal{S}(\Psi_t, S_t)\}
\end{equation}
\end{definition}

\begin{definition}[Target Reachability]
Target $S^*$ is \textbf{reachable} from $S_0$ if:
\begin{equation}
\exists T < \infty : S^* \in \mathcal{R}(S_0, T)
\end{equation}
\end{definition}

\textbf{Reachability conditions:}
\begin{enumerate}
\item \textbf{Direct reachability}: $S^* \in \mathcal{R}(S_0, T)$ for finite $T$
\item \textbf{Context-enabled reachability}: $S^* \notin \mathcal{R}(S_0, T)$ but $\exists$ context change sequence enabling access (Valley Problem, see Appendix UNMAPPED_SS)
\item \textbf{Unreachable}: No feasible path exists (hard constraints block all trajectories)
\end{enumerate}

\subsection{Optimal Control to Target}

The goal-directed Bellman equation becomes:
\begin{equation}
\boxed{
V^*(S_t, \Psi_t) = \min_{I_t \in \mathcal{S}} \left\{ d(S_t, S^*) + \delta \cdot V^*(S_{t+1}, \Psi_{t+1}) \right\}
}
\label{eq:mp-bellman-target}
\end{equation}

with terminal condition $V^*(S^*, \Psi) = 0$ for all $\Psi$.

\textbf{Key insight}: This is a \textbf{shortest path problem} in the 10C state space, but with:
\begin{itemize}
\item \textbf{Dynamic graph}: Edges (feasible transitions) depend on $\Psi$
\item \textbf{Non-convex costs}: Complementarity creates local minima
\item \textbf{Stochastic transitions}: Behavior is probabilistic
\end{itemize}

\subsection{Multi-Target and Prioritization}

When multiple targets exist, define:
\begin{equation}
S^*_{composite} = \arg\min_{S^* \in \{S^*_1, \ldots, S^*_n\}} \sum_{i} \pi_i \cdot d(S_T, S^*_i)
\label{eq:mp-multi-target}
\end{equation}

where $\pi_i$ are target priorities.

\textbf{Pareto optimization}: When targets conflict, find Pareto-optimal trajectories:
\begin{equation}
\{I^*_t\} \in \text{Pareto}(\{d(S_T, S^*_1), \ldots, d(S_T, S^*_n)\})
\end{equation}

\subsection{Connection to $P_{eff}$ Formulation}

The general $P_{eff}$ maximization and goal-directed optimization are connected:
\begin{equation}
\max P_{eff} \approx \max P(a = a^*) = \max P(S_T \in \mathcal{B}_\epsilon(S^*))
\end{equation}

when the target behavior $a^*$ corresponds to the target state $S^*$.

\textbf{Equivalence condition}: If $S^*$ is defined such that $P_{eff}(S^*) = 1$, then:
\begin{equation}
\max \sum_t \delta^t P_{eff}(S_t) \equiv \min \sum_t \delta^t d(S_t, S^*)
\end{equation}

under appropriate metric choice.

\subsection{Summary: Goal-Directed Extensions}

\begin{center}
\begin{tabular}{|l|c|c|}
\hline
\textbf{Aspect} & \textbf{General (MP)} & \textbf{Goal-Directed} \\
\hline
Objective & $\max P_{eff}$ & $\min d(S, S^*)$ \\
Target & Implicit & Explicit $S^*$ \\
Horizon & Open-ended & $T^*$ = first-hitting \\
Success & Higher $P_{eff}$ & Reach $\mathcal{B}_\epsilon(S^*)$ \\
Metric & Probability & Distance in 10C \\
\hline
\end{tabular}
\end{center}

% =============================================================================
\section{Worked Example}
\label{sec:mp-example}
% =============================================================================

\textbf{Anna's Heat Pump Decision}

\textit{Status Quo ($t=0$):}
\begin{itemize}
\item $S_0$: L=0 (individual), high F-utility from savings, low A=0.55, W=-0.50, $\varphi$=S2 (Contemplation)
\item $\Psi_0$: High economic uncertainty, moderate social norm
\item $P_{eff,0} = \sigma(-0.50) = 38\%$
\end{itemize}

\textit{Intervention ($I_0$):}
\begin{itemize}
\item $I_{AWARE} = 0.7$ (information campaign)
\item $I_{WHAT,S} = 0.5$ (social norm messaging)
\item $I_{READY} = 0.6$ (commitment device)
\end{itemize}

\textit{After Intervention ($t=1$):}
\begin{itemize}
\item $S_1$: A=0.72, W=+0.14, $\varphi$=S3 (Preparation)
\item $\Psi_1$: Higher social norm salience
\item $P_{eff,1} = \sigma(+0.14) = 54\%$
\end{itemize}

\textit{Effect:} $\Delta P_{eff} = +16$ percentage points.

\textit{Context Evolution:}
Anna's action (or inaction) changes $\Psi_2$---if she installs, social norm strengthens for neighbors (spillover).

% =============================================================================
\section{Axioms}
\label{sec:mp-axioms}
\label{ax:mp-core}
% =============================================================================

\begin{tcolorbox}[colback=gray!5!white, colframe=gray!75!black,
    title=\textbf{MP Axioms}]

\textbf{Axiom MP-A1 (Objective):} The goal is to maximize $P_{eff}$, the probability of behavioral change.

\textbf{MP-A2 (State):} The state $S$ is fully characterized by the 10C dimensions.

\textbf{MP-A3 (Instruments):} Interventions $I \in [0,1]^9$ are continuous, not discrete types.

\textbf{MP-A4 (Feasibility):} The solution space $\mathcal{S}(\Psi, S)$ depends on context and state.

\textbf{MP-A5 (Endogeneity):} Context evolves: $\Psi_{t+1} = f(\Psi_t, I_t, a_t)$.

\textbf{MP-A6 (Non-Concavity):} The landscape has multiple local maxima due to complementarity.

\textbf{MP-A7 (K-Q Distinction):} Coherence ($K \approx 1$) does not imply optimality ($\max Q$).

\textbf{MP-A8 (Path Dependence):} The optimal path depends on starting point and history.

\textbf{MP-A9 (Multi-Level):} The problem applies at all levels $L_0$ to $L_4$.

\textbf{MP-A10 (Emergence):} Intervention types, life journeys, and equilibria are emergent.

\vspace{0.5em}
\textit{Goal-Directed Extensions:}

\textbf{MP-A11 (Target State):} A target state $S^*$ is a specific point in 10C space representing the desired behavioral configuration.

\textbf{MP-A12 (Distance Metric):} A metric $d(S, S^*)$ quantifies proximity to target, with $d(S^*, S^*) = 0$.

\textbf{MP-A13 (Reachability):} Not all targets are reachable from all starting points; reachability depends on $\mathcal{S}(\Psi, S)$.

\textbf{MP-A14 (Equivalence):} Goal-directed and $P_{eff}$ formulations are equivalent when $P_{eff}(S^*) = 1$ and $d$ is appropriately chosen.

\vspace{0.5em}
\textit{Rationality Foundations (Chapter 2):}

\textbf{MP-A15 (Non-Separability):} Classical separability $C_{ij} = 0$ is relaxed. Utility interdependence $C_{ij} \neq 0$ is the norm, not the exception.

\textbf{MP-A16 (Stability Gap):} Multiple equilibria exist; which one is reached depends on selection dynamics, not optimization alone. Coherence ($K$) differs from optimality ($Q$).

\textbf{MP-A17 (Evolutionary Selection):} Behavioral configurations $C^*$ emerge through selection, not calculation. Replicator dynamics $\dot{x}_i = x_i[\pi_i - \bar{\pi}]$ govern strategy spread.

\textbf{MP-A18 (Bounded Rationality):} Agents exhibit QRE behavior with rationality parameter $\lambda$: $P(s_i) = \exp(\lambda\pi_i)/\sum_j\exp(\lambda\pi_j)$. The parameter $\lambda$ is determined by cognitive context $\Psi_C$.

\vspace{0.5em}
\textit{Anomaly Evidence (Chapter 3):}

\textbf{MP-A19 (Four Interdependencies):} Utility exhibits four structural interdependencies: cross-agent ($\partial U_i/\partial x_j \neq 0$), temporal ($\partial U_t/\partial a_{t-1} \neq 0$), intertemporal inconsistency, and context dependence ($\partial U/\partial \Psi \neq 0$).

\textbf{MP-A20 (Present Bias):} Agents discount quasi-hyperbolically: $U_t = u(c_t) + \beta\sum_\tau \delta^\tau u(c_{t+\tau})$ with $\beta \approx 0.7 < 1$. Time-inconsistent preferences create demand for commitment devices.

\textbf{MP-A21 (Reference Dependence):} Utility is reference-dependent: $U(x) \to U(x - r(\Psi))$ with loss aversion $\lambda \approx 2$. Reference points are context-determined.

\textbf{MP-A22 (Type Heterogeneity):} Population contains heterogeneous types (selfish 20-60\%, inequality-averse 10-20\%, altruistic 25-55\%). BCS segmentation captures this distribution.

\vspace{0.5em}
\textit{Empirical Foundations (Chapter 4):}

\textbf{MP-A23 (Calibrated Parameters):} Framework parameters $\Theta$ derive from 1,900+ experimental papers. Key estimates: $\alpha \approx 0.85$ (envy), $\beta \approx 0.315$ (guilt), $\lambda \approx 2$ (loss aversion), $\beta_{time} \approx 0.7$ (present bias).

\textbf{MP-A24 (Cross-Cultural Variation):} Complementarity structure is universal ($C_{ij} \neq 0$ everywhere); parameters are culturally specific: $C_{ij} = C_{ij}(\Psi_{culture})$. Ultimatum offers range 26-58\% across societies.

\textbf{MP-A25 (Calibration Imperative):} Prediction requires calibration on experimental variation, not training on text. LLMs know ``that'' effects exist but not ``when'' or ``how strongly.''

\textbf{MP-A26 (Self-Improvement):} Deviations $\Delta = C - C^*(\Psi)$ are information, not noise. Five deviation types map to corrective actions. The framework is systematically improvable.

\vspace{0.5em}
\textit{Calibration vs. Simulation (Chapter 4x):}

\textbf{MP-A27 (Propositional vs. Functional):} LLMs have propositional knowledge (``X exists''); prediction requires functional knowledge ($Y = f(X, \Psi)$ with estimated parameters). Text-training $\neq$ behavioral calibration.

\textbf{MP-A28 (Scaling Law):} Predictive accuracy scales with calibration base: 50 experiments → $R^2 \approx 0.4$; 1,500 experiments → $R^2 \approx 0.8$. Information is in experimental variation, not text.

\textbf{MP-A29 (Heterogeneity Signal):} Cross-experimental variation is signal, not noise. The variation identifies $\partial C / \partial \Psi$ --- the functional relationship between context and behavior.

\vspace{0.5em}
\textit{Complementarity Structure (Chapter 5):}

\textbf{MP-A30 (Supermodularity):} Complementarity is defined by $\gamma_{ij} = \partial^2 U / \partial x_i \partial x_j$. The lattice formulation $f(x \vee y) + f(x \wedge y) \geq f(x) + f(y)$ implies coordination dominates mixing.

\textbf{MP-A31 (Four Complementarity Types):} Complementarity manifests in four forms: (i) Strategic ($\partial^2 \Pi_i/\partial a_i \partial a_j > 0$) → multiple equilibria; (ii) Temporal ($\partial^2 U/\partial a_t \partial a_{t+1} > 0$) → path dependence; (iii) Cross-Domain ($\partial^2 U/\partial a^A \partial a^B > 0$) → spillovers; (iv) Cross-Agent ($\partial^2 U_i/\partial x_i \partial x_j > 0$) → coordination problems.

\textbf{MP-A32 (163 Parameters):} The complete complementarity structure requires $\Gamma$ (15 FEPSDE interactions) + $\Lambda$ (28 context interactions) + $\delta$ (120 modulation tensor) = 163 parameters calibrated from experimental evidence (Appendix UNMAPPED_BBB).

\textbf{MP-A33 (Context-Dependent $\gamma$):} Complementarity varies with context: $\gamma_{dd'}(\Psi) = \bar{\gamma}_{dd'} + \sum_k \delta_{dd',k} \cdot (\Psi_k - 0.5)$. The same behavior pair exhibits different interaction strength depending on $\Psi$.

\textbf{MP-A34 (Bundle Dominance):} Strong complementarity ($\gamma > \gamma_{crit}$) implies coherent bundles dominate piecemeal adoption. Ichniowski et al. (1997): HRM bundles raise productivity 7\%; individual practices have negligible effect.

\textbf{MP-A35 (Exclusion Principle):} Classical economics ($\gamma_{ij} = 0$) is the null hypothesis. Complementarity claims ($\gamma \neq 0$) require empirical justification. EBF identifies \textit{where} the null fails, not a priori rejection.

\vspace{0.5em}
\textit{Reference Structure (Chapter 6):}

\textbf{MP-A36 (Fixed-Point Reference):} The reference structure $C^*$ is the unique fixed point $C^* = \Phi(C^*, \Psi)$ where beliefs reproduce themselves through experience. Under contraction ($k < 1$), existence and uniqueness follow from Banach theorem.

\textbf{MP-A37 (Triple Convergence):} Three independent derivations---axiomatic (self-consistency), evolutionary (ESS), empirical (ergodic limit)---yield identical $C^*$. This triple convergence validates the reference structure as fundamental.

\textbf{MP-A38 (Context-Dependent $C^*$):} The reference structure depends on context: $C^*(\Psi_1) \neq C^*(\Psi_2)$ for different $\Psi$. The coherence target is a moving target as context evolves.

\textbf{MP-A39 (Exponential Convergence):} Approach to $C^*$ is exponentially fast: $\|C_n - C^*\| \leq k^n \|C_0 - C^*\|$ with typical $k \approx 0.7$. Systems converge to reference within $\sim$8 iterations.

\vspace{0.5em}
\textit{Non-Concavity and Fit (Chapter 7):}

\textbf{MP-A40 (Non-Concavity):} When $\gamma_{ij} > 0$ for all $i \neq j$ (supermodularity), the payoff function $\Pi$ is not globally concave. The Hessian has positive off-diagonal elements, ensuring multiple local maxima.

\textbf{MP-A41 (Valley Problem):} Transitioning from Peak A to Peak B requires crossing a valley where $K \downarrow$ temporarily. Sequential change fails; all complementary elements must change simultaneously.

\textbf{MP-A42 (Roberts Fit):} Fit (Roberts 2004) = Coherence ($K \approx 1$). Organizational configurations are locally stable when all elements reinforce each other. Toyota, Southwest, Goldman are distinct coherent configurations.

\textbf{MP-A43 (Coherent Obsolescence):} High $K$ does not guarantee high $Q$. A system can be perfectly coherent ($K \approx 1$) in an obsolete configuration ($Q \downarrow$ as $\Psi$ shifts). Kodak 2000: high film-$K$, collapsing digital-$Q$.

\textbf{MP-A44 (Calibration for Navigation):} Strategic navigation requires quantitative parameters: valley depth, peak height, $P(success)$. Text describes \textit{that} landscapes exist; calibration determines \textit{which} path to take.

\vspace{0.5em}
\textit{Mathematical Foundations (Chapter 8):}

\textbf{MP-A45 (Two-Axis Reduction):} Configuration space reduces to $(s, \sigma) \in [0,1]^2$: strategy (stability$\leftrightarrow$dynamism) and structure (hierarchy$\leftrightarrow$flexibility). This Roberts reduction applies at all levels.

\textbf{MP-A46 (Formal Payoff):} $\Pi(s,\sigma;\Psi) = -\gamma(s-\sigma)^2 + \alpha(\Psi)s\sigma + \beta(\Psi)(1-s)(1-\sigma)$. Misfit penalty + innovation synergy + stability synergy. Equilibria at $(0,0)$, $(1,1)$; saddle at $(0.5,0.5)$.

\textbf{MP-A47 (144 Components):} Full welfare aggregates $3 \times 2 \times 6 \times 4 = 144$ utility components: category (INU/KNU/IDN) $\times$ valence (G/P) $\times$ FEPSDE dimension $\times$ time horizon.

\textbf{MP-A48 (Block Structure):} The complementarity matrix $C \in \mathbb{R}^{144 \times 144}$ has hierarchical block structure. Within-category blocks are dense; cross-category blocks are sparse. Effective dimensionality: $\sim$500 parameters.

\textbf{MP-A49 (Computational Tractability):} Low-rank approximation $C \approx U\Sigma V^T$ with $k \approx 10$--$20$ factors captures 90\%+ variance. Hierarchical Bayesian estimation, GPU acceleration, and probabilistic programming make the 144-structure feasible (2026 technology).

\vspace{0.5em}
\textit{Context as Endogenous (Chapter 9 --- CORE-WHEN):}

\textbf{MP-A50 (8Ψ Dimensions):} Context is an 8-dimensional vector $\boldsymbol{\Psi} = (\Psi_I, \Psi_S, \Psi_C, \Psi_K, \Psi_E, \Psi_T, \Psi_M, \Psi_F) \in [0,1]^8$: Institutional, Social, Cognitive, Cultural, Economic, Temporal, Material, Physical.

\textbf{MP-A51 (Endogenous Context):} Context evolves through aggregate behavior: $\Psi_{t+1} = f(\Psi_t, a_t) = \Psi_t + \eta(a_t - a^*(\Psi_t))$. Context is not exogenous background---it is shaped by the behavior it influences.

\textbf{MP-A52 (Multi-Speed Dynamics):} Context components have different adjustment speeds: $\eta_{tech} \gg \eta_{norm} > \eta_{inst} \gg \eta_{culture}$. Technology changes in years; culture changes over generations. This creates cultural lag.

\textbf{MP-A53 (Universal Modulation):} Every EBF parameter is context-dependent: $\gamma(\Psi)$, $\lambda(\Psi)$, $\theta(\Psi)$, $\delta(\Psi)$, $C^*(\Psi)$. This explains heterogeneous treatment effects and why ``one size fits all'' policies fail.

\textbf{MP-A54 (Tipping Points):} S-curve dynamics $\Psi_{t+1} = \Psi_t + \eta\Psi_t(1-\Psi_t)\Delta a_t$ create tipping points. Small interventions near critical thresholds trigger regime change (e.g., smartphone adoption 2007-2015, smoking decline 1965-2020).

\medskip
\textit{Welfare and FEPSDE (Chapter 10 --- CORE-WHAT + CORE-WHO):}

\textbf{MP-A55 (K-Q Separation):} Coherence index $K$ measures equilibrium existence; quality index $Q$ measures equilibrium welfare. High $K$ does not imply high $Q$. A perfectly coherent extractive society ($K = 1$) may have low welfare ($Q \ll Q_{max}$). Both $K$ and $Q$ must be optimized.

\textbf{MP-A56 (FEPSDE Dimensions):} Utility spans six empirically distinct dimensions $d \in \{F, E, P, S, D, Eco\}$: Financial, Emotional, Physical, Social, Digital, Ecological. Each is universally relevant, measurable, and policy-actionable. Single-dimensional optimization (e.g., GDP) is incomplete.

\textbf{MP-A57 (Utility Categories):} Utility is evaluated in three categories: INU (individual---``what I want for myself''), KNU$^L$ (collective---``what I want for us'' at level $L$), IDN (identity---``who I want to be''). People sacrifice INU for KNU (parents for children) and both for IDN (refusing to lie when profitable). Categories are irreducible.

\textbf{MP-A58 (Loss Aversion):} For each dimension $d$, utility is $U_d = G_d - \lambda_d(\Psi) \cdot P_d$ with $\lambda_d > 1$. Empirically: $\lambda_F \approx 2$ (Kahneman-Tversky), $\lambda_P \approx 2$--$3$, $\lambda_S \approx 2$--$4$. Harm prevention creates more welfare than equivalent benefit provision.

\textbf{MP-A59 (144-Component Structure):} Complete welfare aggregates $3 \times 2 \times 6 \times 4 = 144$ utility components: category (INU/KNU/IDN) $\times$ valence (G/P) $\times$ FEPSDE dimension $\times$ time horizon $t \in \{0,1,2,3\}$. The aggregate quality function is $Q(C^*) = \sum_i \sum_d \sum_t \omega_i \delta_d^t [G_{i,d,t} - \lambda_d P_{i,d,t}]$.

\medskip
\textit{Awareness: The Transformation Condition (Chapter 11 --- CORE-AWARE):}

\textbf{MP-A60 (Transformation Principle):} Effective utility is potential utility filtered by awareness: $U_{eff,i,d,t} = A_{i,d,t}(t^*) \times U_{pot,i,d,t}$ where $A \in [0,1]$. Without awareness ($A = 0$), even high $U_{pot}$ has no behavioral effect. This explains why objectively beneficial options are ignored.

\textbf{MP-A61 (Moment of Truth):} Awareness is evaluated at decision point $t^*$. $A(t < t^*)$ and $A(t > t^*)$ may differ dramatically. Post-decision regret (``I knew this'') indicates $A(t^*) \ll A(t > t^*)$. The timing of awareness measurement is critical.

\textbf{MP-A62 (Three-System Architecture):} Awareness operates at three levels: SYS 0 (instinct, instantaneous, hardwired), SYS 1 (intuition, fast, automatic), SYS 2 (deliberation, slow, effortful). Instant Exception Rule: $A(U^P_{t=0}) = 1$ always---immediate pain overrides all other processing.

\textbf{MP-A63 (K-Optimal Portfolio):} Optimal intervention count emerges from context: $K^*(\Psi) = K_{base} + \delta_{budget} + \delta_{complexity} + \delta_{stage} + \delta_{crowding}$ with $K_{base} = 4$ and bounds $K^* \in [2, 9]$ from cognitive capacity limits (Miller's 7±2).

\textbf{MP-A64 (EMERGE Algorithm):} Interventions emerge from optimization, not menu selection: $\vec{I}^* = \text{EMERGE}(\Psi, \Gamma(\Psi), B, K^*) = \arg\max E(\vec{I} | \Psi)$. The 5-step pipeline: Context $\to$ $\gamma(\Psi)$ $\to$ $K^*$ $\to$ EMERGE $\to$ $\vec{I}^*$.

\medskip
\textit{Willingness to Act (Chapter 12 --- CORE-READY):}

\textbf{MP-A65 (Action Condition):} Action occurs if and only if $WAX \geq \theta$. Awareness (AWARE) transforms potential to effective utility; Willingness (READY) determines whether effective utility triggers action. The knowing-doing gap is explained by $WAX < \theta$.

\textbf{MP-A66 (Thinking Style):} $\tau \in [0,1]$ determines how dimensions integrate: $\tau \to 1$ (connected/complementary---only act if everything is right), $\tau \to 0$ (separated/additive---total benefit matters). $\tau = \tau_0 + \Delta\tau(\Psi)$ is context-dependent.

\textbf{MP-A67 (WAX Master Equation):} Willingness to act is $WAX = \tau \cdot \min_i\{U^*_i\} + (1-\tau) \cdot \sum_i U^*_i$ where $U^*_i = A_i \cdot U_i$ is effective utility. High $\tau$ makes the weakest dimension dominate; low $\tau$ allows compensation across dimensions.

\textbf{MP-A68 (Threshold Decomposition):} The action threshold is $\theta = \theta_{base}(d) + \sum_j \delta_j \cdot \Psi_j$. Base threshold varies by decision type (0.2 for habit continuation to 1.5+ for life-changing). Context adjusts: uncertainty $\to \theta \uparrow$, subsidies/social proof $\to \theta \downarrow$.

\textbf{MP-A69 (Gap and Regimes):} The gap $\Delta = WAX - \theta$ determines regime: $\Delta < 0$ (not ready, P < 50\%), $\Delta \approx 0$ (tipping point, P ≈ 50\%), $\Delta > 0$ (ready, P > 50\%). Two paths to close gap: increase $WAX$ (motivation) or decrease $\theta$ (barriers).

\medskip
\textit{Behavioral Change Journey (Chapter 13 --- CORE-STAGE):}

\textbf{MP-A70 (BCJ Six Stages):} The behavioral change journey has six stages $S \in \{S_1, ..., S_6\}$: Unaware ($AWX \approx 0$), Aware (low $AWX$), Considering ($WAX < \theta$), Intending ($WAX \approx \theta$), Acting ($WAX > \theta$), Maintaining ($\theta \to 0$). Each stage has characteristic parameter profiles.

\textbf{MP-A71 (Stage-Dependent Effectiveness):} Same intervention has different effects at different stages. Awareness interventions work at $S_1$-$S_2$; threshold-lowering works at $S_4$. One-size-fits-all interventions fail because stage determines receptivity.

\textbf{MP-A72 (Stage Modulation):} Effective threshold is modulated by stage: $\theta^{eff}(S) = \theta_0 \cdot m(S)$ with $m(S_1) \to \infty$ and $m(S_6) \to 0$. Stage factor $f_S \in [0.6, 1.2]$ affects $K^*$ with maximum at $S_4$ (intention-behavior gap).

\textbf{MP-A73 (Transition Dynamics):} Forward transitions: $P(S_i \to S_{i+1}) = f(AWX, WAX, \theta, \Psi, I)$. Backward transitions (relapse): $P(S_i \to S_{i-1}) = g(\Delta\Psi, t, \text{alternatives})$. Both directions possible; behavioral change is not unidirectional.

\textbf{MP-A74 (Habit Lock-In):} At $S_6$, when $\theta \to 0$, behavior becomes ``locked in'' (Habit Lock-In Time $T_{lock}$). Minimal psychological cost creates behavioral inertia. This is a novel EBF prediction explaining persistence and resistance to change.

\medskip
\textit{Behavioral Change Segments (Chapter 14):}

\textbf{MP-A75 (Heterogeneity Space):} Population heterogeneity forms a space $\mathcal{H} = \{(\omega_i, A_i, W_i)\}$ over utility weights, awareness, and willingness. Segments emerge as natural clusters in this space, not as a priori categories.

\textbf{MP-A76 (Context-Dependent Clustering):} Optimal cluster count $k^* = \arg\min_k\{BIC(k)\} = f(\mathcal{H}, \Psi)$ varies by context. Rogers 5-segment structure emerges only under specific conditions (novel, voluntary, social influence). Mandatory behaviors show 2-3 clusters.

\textbf{MP-A77 (Segment Multiplier):} Treatment effects are heterogeneous: $\sigma_s = CATE(s)/ATE$. $\sigma_s > 1.3$ = segment responds above average; $\sigma_s < 0$ = intervention backfires for this segment. Portfolio optimization requires segment-specific $\sigma_s$ estimates.

\textbf{MP-A78 (BCJ $\times$ BCS Interaction):} Optimal intervention is jointly determined by journey stage (BCJ) and segment type (BCS): $I^* = f(S_{BCJ}, C_{BCS})$. BCJ determines \textit{which} lever; BCS determines \textit{how} to design it. One-size-fits-all fails on both dimensions.

\textbf{MP-A79 (Leverage Profile):} Each segment $k$ has leverage profile $\ell_k^d = \text{Var}(\omega_i^d | i \in C_k) / \sum_{d'}\text{Var}(\omega_i^{d'})$. High $\ell_k^d$ = targeting dimension $d$ creates maximum impact. Interventions must match segment leverage to be effective.

\medskip
\textit{Decision Hierarchy and WEC Coherence (Chapter 15):}

\textbf{MP-A80 (WEC Coherence Requirement):} Effective interventions require simultaneous satisfaction of W (Willingness: can/will person act?), E (Expectation: will intervention work?), and C (Compatibility: do layers reinforce?). Missing any component causes 50-70\% efficiency loss.

\textbf{MP-A81 (Four-Level Hierarchy):} Decision hierarchies have four universal levels: L0 (meta---what counts?), L1 (strategic roots---few, high interaction), L2 (derived strategic---relational to L1), L3 (operative---low interaction). A decision is strategic if removing it breaks coordination, not because it's important.

\textbf{MP-A82 (Universal L2 Formula):} The number of L2 decisions emerges from context: $N_{L2} = \alpha \cdot \gamma_{avg} \times n \times (1-m) / \log(n)$ with $\alpha \approx 0.8$ (bounded rationality constant). Context determines scale (Individual: 3.1, Tribal: 9.2), not structure.

\textbf{MP-A83 (Dynamic Willingness W11-W13):} Willingness is dynamically modulated: $WAX(t) = WAX^0 \times f(A) \times g(KON) \times h(JNY)$. Boundary condition: below $WAX_{blockade}$, no action regardless of other factors. Action probability: $P(Action) = P(WAX \geq blockade) \times P(Opportunity) \times P(Capability)$.

\textbf{MP-A84 (Incoherence Penalty):} Missing predicted L2s causes 50-70\% efficiency loss (HI Theorem T8). Incoherence is costlier than complexity. Better to manage many coherent L2s than few incoherent ones. Cannot assume fixed L2 count---must use formula.

\medskip
\textit{Effective Probability (Chapter 16 --- OBJECTIVE FUNCTION):}

\textbf{MP-A85 ($P_{eff}$ Master Equation):} The objective function is $P_{eff} = \sigma(WEC \cdot \alpha_{BCJ} \cdot \beta_{BCS})$ where $\sigma(x) = 1/(1+e^{-kx})$ is the sigmoid transformation. $P_{eff}$ integrates all previous CORE components into a single probability output.

\textbf{MP-A86 (Sigmoid Threshold):} At $WEC = 0$, $P_{eff} = 50\%$. $WEC < 0 \Rightarrow P_{eff} < 50\%$ (unlikely to act); $WEC > 0 \Rightarrow P_{eff} > 50\%$ (likely to act). Small changes near threshold have maximal $\Delta P_{eff}$ (sigmoid gradient is highest at inflection).

\textbf{MP-A87 (Natural Trajectory Decay):} Without intervention, $P_{eff}$ decays: $A(t) = A_0 \cdot e^{-\lambda t}$ and $W(t) = W_0 \cdot e^{-\mu t}$. The Natural Trajectory Model (NTM) predicts $P_{eff} \to 0$ over time. Intervention is required to counteract decay.

\textbf{MP-A88 (UNTCM Three Regimes):} The Unified NTM-Context Model identifies three regimes: I (decay-dominant, $\kappa < \kappa^*$), II (feedback-dominant, $\kappa > \kappa^*$), III (oscillatory). Tipping point $\bar{B}^* \approx 10$--$25\%$ determines regime transition. Small interventions near $\bar{B}^*$ trigger regime change.

\textbf{MP-A89 (Loss Aversion Lock-in X5):} Status quo strengthens over time: $\theta_{eff}(t) = \theta_0 + \lambda_{LA} \cdot |S_{target} - r(t)|$ with $\lambda_{LA} \approx 2.25$. Late intervention faces $\sim$225\% higher effective threshold due to reference point crystallization. Early intervention is cheaper.

\medskip
\textit{Emergent Intervention Theory (Chapter 17):}

\textbf{MP-A90 (Axiom A0 Foundation):} All of EBF emerges from one axiom: $A0 = \exists a,G,U,K$ (agent $a$ pursues goals $G$ under uncertainty $U$ in context $K$). From A0, the 10C questions emerge necessarily. The intervention space is not arbitrary but emerges from this fundamental axiom.

\textbf{MP-A91 (10C Status Quo Vector):} Before intervention, diagnose current state: $\vec{S}_0 = (L, d, \gamma_{current}, \Psi, \Theta, A_0, W_0, \varphi_0, Scope)$. No intervention without 10C Status Quo Assessment. This vector is the input for intervention emergence.

\textbf{MP-A92 (Intervention as Transformation):} An intervention $\mathcal{I}$ transforms status quo: $\mathcal{I}: \vec{S}_0 \mapsto \vec{S}_1 = \vec{S}_0 + \Delta\vec{S}(\mathcal{I})$. The question is not ``which type?'' but ``which transformation $\Delta\vec{S}$ is optimal for this specific $\vec{S}_0$?''

\textbf{MP-A93 (Continuous Intervention Space):} Interventions are points in continuous 9-dimensional space: $\vec{I} = (I_{WHO}, I_{WHAT}, I_{HOW}, I_{WHEN}, I_{WHERE}, I_{AWARE}, I_{READY}, I_{STAGE}, I_{HIER}) \in [0,1]^9$. \textbf{Note: 9D within 10C}---The vector targets COREs 1--9; CORE 10 (EIT) is the \textit{methodology} that generates interventions, not a target itself. Discrete ``types'' (T1-T8) are projections onto single axes. 90\% of real interventions lie between types. The color spectrum analogy: labels are useful but not ontological.

\textbf{MP-A94 (Emergence Functions EIT-6/7/8):} Intervention emphasis emerges from status quo: EIT-6 ($A_0 > 0.55 \Rightarrow I_{AWARE}$ low), EIT-7 ($W_0 < 0.25 \Rightarrow I_{READY}$ maximal), EIT-8 ($\Psi_c$ changeable $\Rightarrow I_{WHEN}$ high). Context determines which dimensions to emphasize.

\medskip
\textit{Emergent Temporal Theory (Chapter 18):}

\textbf{MP-A95 (EIT-4: Journey-Position Emergence):} Journey-position $\varphi$ emerges from 10C Status Quo: $\varphi = h(A_0, W_0, Scope)$. Journey phases are \textit{regions} in 10C space, not ontological categories. Boundaries are fuzzy; between-phase positions are valid states.

\textbf{MP-A96 (EIT-5: Effect-Horizon Emergence):} Temporal effect-horizon $\tau$ emerges from intervention vector: $\tau(\vec{I}) = g(I_{HIER}, I_{WHO}, I_{HOW})$. Interventions changing $I_{HOW}$ (complementarities) automatically have longer horizons than those affecting only $I_{AWARE}$.

\textbf{MP-A97 (Phase-Affinity Field):} Phase-affinity is continuous: $\alpha: [0,1]^9 \times \mathbb{R}^{10C} \to [0,1]$. Factorization: $\alpha(\vec{I}, \vec{S}_0) \approx \alpha_{base}(\vec{I}) \cdot \alpha_{match}(\vec{I}, \varphi) \cdot \alpha_{context}(\Psi)$. No universally ``best'' intervention---optimality depends on journey region.

\textbf{MP-A98 (Phase-K* Mapping):} Stage factor $f_S$ modulates $K^*$: Pre-Aware ($f_S = 0.7$), Triggered ($f_S = 1.2$), Action ($f_S = 1.0$), Maintenance ($f_S = 0.8$). Triggered phase has highest $f_S$ (critical window). Mismatch detection: $\alpha_{match} < 0.3$ signals poor timing.

\textbf{--- Chapter 19: Segment-Heterogeneity ---}

\textbf{MP-A99 (PAP1-1: Social Preference Heterogeneity):} Individuals differ in social preferences: $U_i(x_i, x_{-i}) = u(x_i) + \alpha_i \cdot v(x_{-i})$ with $\alpha_i \in [0,1]$. Segmentation: $\alpha < 0.2$ (Self-interested), $\alpha > 0.5$ (Social-Oriented). Social-Oriented respond to norms ($\sigma = 1.8$), Self-interested to incentives.

\textbf{MP-A100 (PAP1-2: Temporal Preference Heterogeneity):} Individuals differ in present bias: $U_i = u(c_0) + \beta_i \sum \delta^t u(c_t)$ with $\beta_i \in (0,1]$. Segmentation: $\beta > 0.9$ (Time-consistent), $\beta < 0.7$ (Strong Present-Biased). Present-Biased need commitment ($\sigma = 1.9$), not information.

\textbf{MP-A101 (PAP1-3: Autonomy-Reactance Heterogeneity):} Individuals differ in reactance sensitivity: $\sigma_i(\mathcal{I}) = \sigma_{base} - \rho_i \cdot \text{Perceived\_Manipulation}$. High $\rho_i$ with perceived manipulation yields $\sigma_i < 0$ (backfire). Autonomy-Seeking require invisible interventions.

\textbf{MP-A102 (Segment-Multiplier):} Segment-multiplier $\sigma_s$ scales base effectiveness: $E(\mathcal{I}|s) = E_{base} \cdot \sigma_s$ with $\sigma_s \in [-0.5, 2.0]$. Estimation: $\hat{\sigma}_s = CATE(s)/ATE$. Same intervention has dramatically different effects across segments.

\textbf{MP-A103 (Phase$\times$Segment Interaction):} Phase and segment effects interact multiplicatively with interaction coefficient $\delta_{t,s}$: $E = E_{base} \cdot \alpha_{t,\varphi} \cdot \sigma_s \cdot \delta_{t,s}$. Key synergies: Present-Biased$\times$Action ($\delta = 1.4$), Social$\times$Triggered ($\delta = 1.5$). Key interference: Autonomy$\times$Maintenance ($\delta = 0.7$).

\textbf{MP-A104 (Reactance-Mitigation):} For Autonomy-Oriented ($\rho > 0.5$), explicit autonomy cues improve effectiveness: $\sigma_{with\_cue} = \sigma_{base} + 0.4 \cdot \rho$. Can transform negative $\sigma$ to positive by respecting freedom.

\textbf{MP-A105 (Choice Satisfaction Constraint):} Menu satisfaction follows Miller-Iyengar-Schwartz: $Z(n) = Z_{max} \times (1 - 0.03 \times \max(0, n-7))$. Beyond 7 options, each additional option reduces satisfaction by $\approx 3\%$. Optimal menu size: 5--7 options.

\textbf{--- Chapter 20: Portfolio Optimization ---}

\textbf{MP-A106 (A1-A8 Heuristic Clusters):} Continuous $\vec{I} \in [0,1]^9$ clusters into 8 heuristics via PCA+k-means: A1 (Awareness), A2 (Feedback), A3 (Choice Architecture), A4 (Temporal), A5 (Identity), A6 (Social), A7 (Financial), A8 (Commitment). Clusters are practical shortcuts; continuous model is rigorous.

\textbf{MP-A107 (Portfolio Complementarity):} Portfolio effectiveness includes complementarity: $E(P) = \sum_i E_i + \sum_{i<j} \gamma_{ij} \sqrt{E_i E_j}$. Known conflicts: Social+Financial ($\gamma = -0.3$), Financial+Intrinsic ($\gamma = -0.3$). Coherent portfolios maximize synergies, avoid conflicts.

\textbf{MP-A108 (Seven Constraints C1-C7):} Portfolio optimization subject to: C1 (Budget), C2 (Crowding: $\gamma \geq \gamma_{min}$), C3 (Timing: sequencing), C4 (Capacity), C5 (Ethical: autonomy $\geq \eta_{min}$), C6 (Segment: $\sigma \geq \sigma_{min}$), C7 (K-Optimal: $|P| \leq K^*$). Constraints are hard; violations invalidate portfolio.

\textbf{MP-A109 (K* Emergence---EXC-3 Hybrid):} Optimal portfolio size emerges: $K^* = K_{base} \cdot f_B + \delta_C + \delta_S + \delta_\gamma + \delta_L$ where $f_B$ is \textit{multiplicative} (veto logic: $B=0 \Rightarrow K^*=0$) and $\delta_C, \delta_S, \delta_\gamma, \delta_L$ are \textit{additive} (EXC-1 default). K* is not arbitrary---it emerges from context per Exclusion Principle (FRM, EXC-3).

\textbf{MP-A110 (BC Prediction Pipeline):} Behavioral change probability: $P(BC | \vec{S}_0, P) = \sigma(\beta_0 + \beta_1 \cdot E(P | \vec{S}_0))$. Conditional effectiveness: $E(P|\vec{S}_0) = \sum_i E_{max,i} \cdot \alpha_{t,\varphi_0} \cdot \sigma_s \cdot \lambda_L \cdot \gamma(\Psi) + \sum_{i<j} \gamma_{ij}\sqrt{E_i E_j}$.

\textbf{MP-A111 (Six Frameworks):} Six equally valid heuristic frameworks: F1 (Component/A1-A8), F2 (Behavioral/Problem-type), F3 (Phase-Based), F4 (Segment-First), F5 (Outcome-First/FEPSDE), F6 (Context-Adaptive/$\Psi$). Choice depends on data availability, heterogeneity, context stability.

\textbf{MP-A112 (Hierarchical Optimization):} Portfolio optimization at three levels: Strategic (annual, $K^* \times 0.8$), Tactical (quarterly, $K^* \times 1.0$), Operative (weekly, $K^* \times 1.2$). Cascade: Strategic $\to$ Tactical $\to$ Operative with increasing granularity.

\textbf{--- Chapter 21: Dynamic Portfolio Evolution ---}

\textbf{MP-A113 (Three Evolution Mechanisms):} Status quo $S_t$ evolves through: (1) Phase progression $\varphi_t$ (BCJ stages), (2) Context shifts $\Psi_t$ (policy/economic/social), (3) Bayesian learning $\Theta_t$ (parameter updates). All three trigger portfolio redesign.

\textbf{MP-A114 (Dynamic K* Formula):} Optimal portfolio size evolves: $K^*(t) = K^*_0 \cdot g(\varphi_t) \cdot h(\Psi_t) \cdot \ell(\Theta_t)$. Factors: $g = 1 + 0.3|\Delta\varphi|$ (phase), $h = 1 + 0.5 \cdot \mathbb{1}[\text{shock}]$ (context), $\ell = 1 - 0.2 \cdot \text{converged}$ (learning).

\textbf{MP-A115 (Phase Velocity ODE):} Phase progression rate: $d\varphi/dt = v_0 + v_A \cdot A_t + v_W \cdot W_t$. Higher awareness and willingness accelerate phase progression. Calibrated from empirical data (typically 3--12 months per phase).

\textbf{MP-A116 (Redesign Triggers):} Portfolio redesign triggered by: $|\Delta\varphi| > 0.2$ (phase jump), $||\Delta\Psi|| > 0.15$ (context shock), or significant $\Delta\Theta$ (parameter update). Triggers are OR-connected; any one suffices.

\textbf{MP-A117 (Convergence Criteria):} Learning converges when $||\Theta_t - \Theta_{t-1}|| < \epsilon$ for $T$ consecutive periods. Practical: $\epsilon = 0.05$, $T = 2$. Convergence signals: stop redesigning, maintain current portfolio.

\textbf{MP-A118 (K* Evolution Pattern):} Universal pattern across domains: Baseline (K* = 2--3) $\to$ Transition/Shock (K* = 4--5) $\to$ Convergence (K* = 3). Crisis requires expanded portfolio; stability allows refocusing.

\end{tcolorbox}

% =============================================================================
\section{Theorems}
\label{sec:mp-theorems}
% =============================================================================

\begin{theorem}[MP-T1: Existence]
Under regularity conditions (bounded $\Psi$, continuous $P_{eff}$, compact $\mathcal{S}$),
an optimal policy $\{I^*_t\}$ exists.
\end{theorem}

\begin{theorem}[MP-T2: Non-Uniqueness]
Due to non-concavity (MP-A6), multiple local optima may exist. The global optimum
depends on context $\Psi$ (which peak is highest).
\end{theorem}

\begin{theorem}[MP-T3: Path Dependence]
The optimal trajectory $\{S^*_t, \Psi^*_t\}$ depends on initial conditions $(S_0, \Psi_0)$
and the sequence of interventions. Different paths may lead to different equilibria.
\end{theorem}

\begin{theorem}[MP-T4: Solution Space Emergence]
The solution space $\mathcal{S}_{t+1}$ at time $t+1$ emerges from context $\Psi_{t+1}$,
which itself depends on interventions $I_t$. Thus:
\[
\mathcal{S}_{t+1} = \mathcal{S}(f(\Psi_t, I_t, a_t), S_t + \Delta S(I_t))
\]
The feasible set evolves endogenously.
\end{theorem}

\begin{theorem}[MP-T5: Target Reachability]
Target $S^*$ is reachable from $(S_0, \Psi_0)$ if and only if there exists a finite
sequence of interventions $\{I_t\}_{t=0}^{T-1}$ such that:
\begin{enumerate}
\item $I_t \in \mathcal{S}(\Psi_t, S_t)$ for all $t$ (feasibility)
\item $S_T \in \mathcal{B}_\epsilon(S^*)$ (terminal condition)
\end{enumerate}
If $S^* \notin \mathcal{R}(S_0, T)$ for any finite $T$, the target is unreachable
and requires context change or constraint relaxation.
\end{theorem}

\begin{theorem}[MP-T6: Optimal Path to Target]
If target $S^*$ is reachable, the optimal path $\{I^*_t\}$ minimizing $\sum_t d(S_t, S^*)$
satisfies the goal-directed Bellman equation:
\[
V^*(S, \Psi) = \min_{I \in \mathcal{S}} \left\{ d(S, S^*) + \delta \cdot V^*(S', \Psi') \right\}
\]
The value function $V^*(S^*, \Psi) = 0$ for all $\Psi$.
\end{theorem}

\begin{theorem}[MP-T7: Evolutionary Stability]
A behavioral configuration $C^*$ is \textbf{evolutionarily stable} if and only if:
\begin{enumerate}
\item $\pi(C^*, C^*) \geq \pi(C', C^*)$ for all $C' \neq C^*$ (Nash condition)
\item If $\pi(C^*, C^*) = \pi(C', C^*)$, then $\pi(C^*, C') > \pi(C', C')$ (stability condition)
\end{enumerate}
Every ESS is a Nash equilibrium, but not vice versa. Stability requires robustness to invasion.
\end{theorem}

\begin{theorem}[MP-T8: K-Q Independence]
Coherence $K$ and welfare $Q$ are independent: There exist configurations with:
\begin{itemize}
\item High $K$, high $Q$ (optimal peak)
\item High $K$, low $Q$ (stable trap)
\item Low $K$, transitional $Q$ (valley crossing)
\end{itemize}
Maximizing $K$ does not guarantee maximizing $Q$. Transition between peaks requires
temporary $K$ decrease (incoherence) during valley crossing.
\end{theorem}

\begin{theorem}[MP-T9: Present Bias and Commitment]
For agents with $\beta < 1$ (present bias), commitment devices $\kappa$ that
restrict future choice sets can be welfare-improving:
\[
\exists \kappa : U^{sophisticated}(\text{with } \kappa) > U^{sophisticated}(\text{without } \kappa)
\]
The optimal intervention portfolio $P^*$ includes commitment mechanisms when
the target population has $\bar{\beta} < 0.85$.
\end{theorem}

\begin{theorem}[MP-T10: Reference Point Dynamics]
Reference points $r_t$ evolve according to:
\[
r_{t+1} = (1-\rho) \cdot r_t + \rho \cdot x_t
\]
where $\rho \in (0,1)$ is the adaptation rate. This creates path dependence in
$P_{eff}$: the same intervention $I$ produces different effects depending on
the reference point trajectory $\{r_\tau\}_{\tau \leq t}$.
\end{theorem}

\begin{theorem}[MP-T11: Supermodular Equilibria]
If the objective function $P_{eff}(S, I, \Psi)$ is supermodular in $(S, I)$, then:
\begin{enumerate}
\item Pure-strategy Nash equilibria exist (no convexity required)
\item Equilibria form a complete lattice with smallest ($\underline{e}$) and largest ($\bar{e}$) equilibria
\item Comparative statics are monotone: higher parameters induce higher equilibria
\item Best responses are isotone: $BR_i(a_{-i}) \uparrow$ when $a_{-i} \uparrow$
\end{enumerate}
The optimization problem thus has well-behaved solutions despite non-convexity.
\end{theorem}

\begin{theorem}[MP-T12: Context Modulation]
For context-dependent complementarity $\gamma_{dd'}(\Psi) = \bar{\gamma}_{dd'} + \sum_k \delta_{dd',k}(\Psi_k - 0.5)$:
\begin{enumerate}
\item The landscape topology changes with $\Psi$: same intervention, different effects
\item Critical thresholds exist: $\exists \Psi^* : \gamma(\Psi^*) = \gamma_{crit}$ triggers phase transition
\item Context can unlock/lock equilibria: $e^* \in \mathcal{E}(\Psi)$ but $e^* \notin \mathcal{E}(\Psi')$
\end{enumerate}
Implication: Context interventions ($\Delta\Psi$) can be more effective than direct interventions ($\Delta I$).
\end{theorem}

\begin{theorem}[MP-T13: Reference Structure Convergence]
Under regularity (D-REG), reciprocity (D-REC), and contraction (D-CON) axioms:
\begin{enumerate}
\item \textbf{Existence:} $\exists$ unique $C^*$ satisfying $C^* = \Phi(C^*, \Psi)$
\item \textbf{Triple Convergence:} $C^*_{axiomatic} = C^*_{evolutionary} = C^*_{empirical}$
\item \textbf{Rate:} $\|C_n - C^*\| \leq k^n\|C_0 - C^*\|$ with $k < 1$
\item \textbf{Context Dependence:} $C^*(\Psi)$ is continuous in $\Psi$
\end{enumerate}
The coherence index $K = 1 - \|C - C^*\|/\|C^*\|$ is therefore well-defined for any $(C, \Psi)$.
\end{theorem}

\begin{theorem}[MP-T14: Non-Concavity and Multiple Equilibria]
Under supermodularity ($\gamma_{ij} > 0$ for $i \neq j$):
\begin{enumerate}
\item \textbf{Non-Concavity:} $\Pi(C)$ has positive Hessian off-diagonals $\Rightarrow$ not globally concave
\item \textbf{Multiple Peaks:} $\exists$ at least two local maxima $C^*_A, C^*_B$ with $K(C^*_A) \approx K(C^*_B) \approx 1$
\item \textbf{Valley Separation:} $\exists$ valley $V$ between peaks where $K < K_{crit}$
\item \textbf{Transition Cost:} $\min_{\text{path}} \int K(t) dt < K(C^*_A) \cdot T$ (must lose coherence to gain it)
\end{enumerate}
Partial change (moving in one dimension while others fixed) increases valley depth.
\end{theorem}

\begin{theorem}[MP-T15: K-Q Independence]
Coherence $K$ and quality $Q$ are independent dimensions:
\begin{enumerate}
\item $\exists$ configurations with high $K$, high $Q$ (optimal coherent)
\item $\exists$ configurations with high $K$, low $Q$ (coherent obsolescence)
\item $\exists$ configurations with low $K$, high potential $Q$ (transition state)
\end{enumerate}
Context shift $\Delta\Psi$ can change $Q(C^*)$ without changing $K$, creating the coherent obsolescence trap.
\end{theorem}

\begin{theorem}[MP-T16: Two-Axis Equilibrium Structure]
For the payoff function $\Pi(s,\sigma;\Psi) = -\gamma(s-\sigma)^2 + \alpha s\sigma + \beta(1-s)(1-\sigma)$:
\begin{enumerate}
\item Two local maxima exist: $A = (0,0)$ with $\Pi_A = \beta$, $B = (1,1)$ with $\Pi_B = \alpha$
\item One saddle point at $M = (0.5, 0.5)$ with $\Pi_M = (\alpha + \beta)/4$
\item Two local minima (misfit): $(0,1)$ and $(1,0)$ with $\Pi = -\gamma$
\item Peak height determined by context: $\alpha(\Psi) > \beta(\Psi) \Rightarrow B$ dominates
\end{enumerate}
\end{theorem}

\begin{theorem}[MP-T17: 144-Component Tractability]
The $144 \times 144$ complementarity matrix $C$ with 10,440 unique entries is computationally tractable:
\begin{enumerate}
\item Block structure: 9 blocks with within-block density $\gg$ cross-block density
\item Low-rank: $C \approx U\Sigma V^T$ with rank $k \approx 15$ captures $>90\%$ variance
\item Effective parameters: $\sim$500 non-negligible entries
\item Estimation: Hierarchical Bayes pools information across similar $(i,j,\Psi)$ contexts
\end{enumerate}
The structure-exploitation makes the full EBF optimization problem solvable.
\end{theorem}

\begin{theorem}[MP-T18: Context Feedback Loop]
The system $(\Psi_t, C^*_t, a_t)$ evolves according to:
\begin{enumerate}
\item $\Psi_t \to C^*(\Psi_t)$: Context determines reference structure
\item $C^*_t \to a^*_t$: Reference structure determines optimal behavior
\item $a_t \to \Psi_{t+1}$: Aggregate behavior updates context
\end{enumerate}
If $a_t = a^*(\Psi_t)$, context is stationary. If $a_t \neq a^*$, context drifts.
Multiple stable equilibria exist when self-reinforcing feedback dominates.
\end{theorem}

\begin{theorem}[MP-T19: Tipping Point Dynamics]
For S-curve context dynamics $\Psi_{t+1} = \Psi_t + \eta\Psi_t(1-\Psi_t)\Delta a_t$:
\begin{enumerate}
\item Stable equilibria exist at $\Psi = 0$ and $\Psi = 1$
\item Unstable equilibrium at $\Psi = 0.5$ (tipping point)
\item Near $\Psi = 0.5$: small $\Delta a$ causes large $\Delta\Psi$ (regime change)
\item Transition speed maximized at $\Psi = 0.5$: $\partial^2\Psi/\partial t^2 = 0$
\end{enumerate}
Policy implication: Interventions near tipping points have disproportionate effect.
\end{theorem}

\begin{theorem}[MP-T20: K-Q Independence]
Coherence ($K$) and quality ($Q$) are independent indices. Specifically:
\begin{enumerate}
\item $K = 1$ (perfect coherence) is compatible with any $Q \in [Q_{min}, Q_{max}]$
\item $\partial K/\partial Q = 0$ at equilibrium
\item Optimizing $K$ alone may decrease $Q$ (coordination on bad equilibrium)
\item Optimizing $Q$ alone may decrease $K$ (welfare improvement without coordination)
\end{enumerate}
The complete objective must optimize $f(K, Q)$, not $K$ or $Q$ alone.
\end{theorem}

\begin{theorem}[MP-T21: Welfare Aggregation]
The aggregate quality function $Q(C^*)$ aggregates 144 utility components via:
\[
Q = \sum_{i \in \{INU, KNU, IDN\}} \sum_{d \in FEPSDE} \sum_{t=0}^{3}
    \omega_i(\Psi) \cdot \delta_d(\Psi)^t \left[ G_{i,d,t} - \lambda_d(\Psi) \cdot P_{i,d,t} \right]
\]
Properties:
\begin{enumerate}
\item Loss aversion ($\lambda_d > 1$) implies harm prevention dominates benefit provision
\item Dimension-specific discounting allows health ($\delta_P > \delta_F$) to be privileged
\item Category weights $\omega_i$ vary with context (collectivist vs. individualist cultures)
\item Level-indexed KNU$^L$ enables multi-level welfare accounting
\end{enumerate}
\end{theorem}

\begin{theorem}[MP-T22: Awareness Transformation]
The transformation from potential to effective utility satisfies:
\begin{enumerate}
\item $U_{eff} = A(t^*) \cdot U_{pot}$ filters utility through awareness at Moment of Truth
\item $\sum_i \sum_d \sum_t U_{eff,i,d,t} \leq \sum_i \sum_d \sum_t U_{pot,i,d,t}$ (awareness reduces, never amplifies)
\item Instant Exception: $A(U^P_{t=0}) = 1$ always (pain overrides)
\item Behavioral influence requires $A > 0$ AND perception at $t^*$
\end{enumerate}
The optimization problem's objective uses $U_{eff}$, not $U_{pot}$.
\end{theorem}

\begin{theorem}[MP-T23: K-Optimal Bounds]
The optimal intervention count $K^*$ satisfies:
\begin{enumerate}
\item $K^* \in [2, 9]$ bounded by cognitive capacity (Miller's 7±2)
\item $K^* = K_{base} + \sum_j \delta_j(\Psi)$ is context-dependent
\item More resources ($\delta_{budget} > 0$) allows larger $K^*$
\item Earlier BCJ stages ($\delta_{stage} < 0$) require smaller $K^*$
\item Crowding risk ($\delta_{crowding} < 0$) reduces effective $K^*$
\end{enumerate}
EMERGE algorithm: $\vec{I}^* = \arg\max_{\|\vec{I}\|_0 \leq K^*} E(\vec{I} | \Psi)$.
\end{theorem}

\begin{theorem}[MP-T24: Action Condition]
Behavioral action occurs if and only if $WAX \geq \theta$:
\begin{enumerate}
\item $WAX = \tau \cdot \min_i\{U^*_i\} + (1-\tau) \cdot \sum_i U^*_i$ (willingness formula)
\item $\theta = \theta_{base}(d) + \sum_j \delta_j \cdot \Psi_j$ (threshold formula)
\item Gap $\Delta = WAX - \theta$ determines P(Action) via $\sigma(\Delta)$
\item Three regimes: Below ($\Delta < 0$), Critical ($\Delta \approx 0$), Above ($\Delta > 0$)
\end{enumerate}
The knowing-doing gap is explained by $WAX < \theta$ even when $A(\cdot) > 0$.
\end{theorem}

\begin{theorem}[MP-T25: Dual Path Intervention]
To change behavior, interventions can target either $WAX$ or $\theta$ (or both):
\begin{enumerate}
\item Path 1 ($\uparrow WAX$): Increase $\tau$, increase $A$, increase $U$
\item Path 2 ($\downarrow \theta$): Reduce friction, change defaults, provide subsidies
\item Combined effect: $\Delta_{new} = (WAX + \Delta WAX) - (\theta - \Delta\theta)$
\item Optimal strategy: Address both paths when $|\Delta|$ is large
\end{enumerate}
High $\tau$ individuals require all dimensions addressed; low $\tau$ allows focused interventions.
\end{theorem}

\begin{theorem}[MP-T26: Stage-Dependent Optimization]
Optimal intervention strategy depends on BCJ stage:
\begin{enumerate}
\item Stages $S_1$-$S_2$: Maximize $\Delta AWX$ (awareness creation)
\item Stage $S_3$: Balance $\uparrow WAX$ and $\downarrow \theta$ (consideration support)
\item Stage $S_4$: Maximize $\theta$ reduction (close intention-behavior gap)
\item Stage $S_5$: Minimize friction (implementation support)
\item Stage $S_6$: Reinforce and prevent relapse (habit maintenance)
\end{enumerate}
Stage factor $\delta_S$ captures this in EMERGE: $K^* = K_{base} \cdot f_B + ... + \delta_S + ...$ (hybrid per EXC-3)
\end{theorem}

\begin{theorem}[MP-T27: Habit Lock-In]
Once $S_6$ is reached with $\theta \to 0$:
\begin{enumerate}
\item Behavior becomes automatic (no conscious decision required)
\item Psychological cost $\to 0$ (minimal effort to maintain)
\item Behavioral inertia resists change (switching costs high)
\item New identity forms: ``I am someone who does this''
\end{enumerate}
Habit Lock-In Time $T_{lock}$ is the duration after which $\theta < \theta_{crit}$.
Relapse probability decreases exponentially: $P(relapse) \propto e^{-t/\tau}$ for $t > T_{lock}$.
\end{theorem}

\begin{theorem}[MP-T28: Heterogeneous Treatment Effects]
Treatment effects vary systematically by segment:
\begin{enumerate}
\item Segment multiplier: $\sigma_s = CATE(s)/ATE = E[Y(1)-Y(0)|S=s]/E[Y(1)-Y(0)]$
\item Portfolio effectiveness: $P_{eff}^{portfolio} = \sum_s w_s \cdot \sigma_s \cdot P_{eff}^{avg}$
\item Backfire risk: If $\sigma_s < 0$ for segment $s$, intervention harms that segment
\item Optimal allocation: Weight interventions toward high-$\sigma_s$ segments
\end{enumerate}
Ignoring heterogeneity underestimates effectiveness for responsive segments and misses backfire risks.
\end{theorem}

\begin{theorem}[MP-T29: BCJ-BCS Joint Optimization]
Optimal intervention design requires joint consideration of stage and segment:
\begin{enumerate}
\item BCJ determines receptivity: Stage $S$ determines which intervention types can work
\item BCS determines framing: Segment $C$ determines how to design the intervention
\item Interaction: $I^* = f(S_{BCJ}, C_{BCS})$ with stage-segment affinity matrix $\alpha_{S,C}$
\item Portfolio structure: Different segments may require different stage progressions
\end{enumerate}
A segment-aware, stage-appropriate intervention outperforms generic approaches by factor 3-10x.
\end{theorem}

\begin{theorem}[MP-T30: Universal L2 Formula]
The number of derived strategic decisions $N_{L2}$ is determined by context:
\begin{enumerate}
\item Formula: $N_{L2} = \alpha \cdot \gamma_{avg} \times n \times (1-m) / \log(n)$
\item Bounded rationality: $\alpha \approx 0.8$ is invariant across scales (humans manage $\sim$4 L2s coherently)
\item Scale independence: Structure (L0-L1-L2-L3) is universal; only $N_{L2}$ varies
\item Empirical fit: $< 5\%$ error across individual, household, organization, nation, tribal contexts
\end{enumerate}
Attempting to ``simplify'' by reducing L2 count below formula prediction destroys coherence.
\end{theorem}

\begin{theorem}[MP-T31: WEC Coherence Constraint]
Interventions must satisfy WEC coherence across all predicted L2 decisions:
\begin{enumerate}
\item For each L2: W (willingness supported?), E (expectation realistic?), C (compatible with other L2s?)
\item Red flag: If any WEC component is NO for any L2, intervention will fail
\item Efficiency: Coherent design yields 80\%+ success; incoherent design yields 30-40\%
\item Diagnosis: Incoherence manifests as unexpected resistance, low adoption, early dropout
\end{enumerate}
The WEC coherence check is the critical validation step before implementation.
\end{theorem}

\begin{theorem}[MP-T32: $P_{eff}$ as Objective Function]
The optimization problem's objective is the effective probability of behavioral change:
\begin{enumerate}
\item Master equation: $P_{eff} = \sigma(WEC \cdot \alpha_{BCJ} \cdot \beta_{BCS})$
\item Sigmoid nonlinearity: Maximum sensitivity at $WEC = 0$ (threshold)
\item All CORE components (WHO, WHAT, HOW, WHEN, WHERE, AWARE, READY, STAGE) integrate into $P_{eff}$
\item Intervention design optimizes $\max \sum_t \delta^t P_{eff}(S_t, I_t, \Psi_t)$
\end{enumerate}
$P_{eff}$ is the single output that practitioners measure and optimize.
\end{theorem}

\begin{theorem}[MP-T33: Natural Trajectory and Decay]
Without intervention, $P_{eff}$ decays to zero (Natural Trajectory Model):
\begin{enumerate}
\item Awareness decay: $A(t) = A_0 \cdot e^{-\lambda t}$ with dimension-specific $\lambda_d$
\item Willingness decay: $W(t) = W_0 \cdot e^{-\mu t}$
\item Behavioral decay: $B(t) \to 0$ in absence of reinforcement
\item Loss aversion lock-in: $\theta_{eff}(t) = \theta_0 + \lambda_{LA} \cdot |S_{target} - r(t)|$ with $\lambda_{LA} \approx 2.25$
\end{enumerate}
Implication: Intervention cost increases superlinearly with delay ($\sim$3.25$\times$ at 12 months).
\end{theorem}

\begin{theorem}[MP-T34: UNTCM Regime Transitions]
The Unified NTM-Context Model identifies regime transitions at tipping points:
\begin{enumerate}
\item Regime I (Decay): $\kappa < \kappa^*$ --- intervention required to prevent fade
\item Regime II (Feedback): $\kappa > \kappa^*$ --- self-sustaining adoption after crossing threshold
\item Regime III (Oscillatory): Complex dynamics, seasonal patterns
\item Tipping point: $\bar{B}^* \approx 10$--$25\%$ adoption triggers transition I $\to$ II
\end{enumerate}
Strategic implication: Target early adopters to reach $\bar{B}^*$ and trigger positive feedback.
\end{theorem}

\begin{theorem}[MP-T35: Continuous Intervention Space]
The intervention space is continuous, not discrete:
\begin{enumerate}
\item Vector representation: $\vec{I} = \sum_{k \in 10C} I_k \cdot \hat{e}_k \in [0,1]^9$
\item No sharp boundaries between ``types'' --- 90\% of interventions lie between
\item Combinability: Vectors can be added, distances and angles defined
\item Optimizability: Gradients computable for continuous optimization
\end{enumerate}
Implication: Design interventions as vectors, not as labels from a catalog.
\end{theorem}

\begin{theorem}[MP-T36: Intervention Emergence from Status Quo]
Optimal intervention emphasis emerges from 10C status quo assessment:
\begin{enumerate}
\item $A_0 \leq 0.35$: $I_{AWARE}$ high (awareness campaign needed)
\item $A_0 > 0.55$: $I_{AWARE}$ low (awareness not the problem)
\item $W_0 < 0.25$ AND $A_0 > 0.5$: $I_{READY}$ maximal (willingness is bottleneck)
\item $\Psi_c$ easily changeable: $I_{WHEN}$ high (context intervention)
\end{enumerate}
The intervention profile is not chosen from menu but emerges from diagnosis.
\end{theorem}

\begin{theorem}[MP-T37: Intervention Trajectory Optimization]
Interventions are trajectories, not point actions:
\begin{enumerate}
\item Definition: $I(t): [t_0, t_1] \to [0,1]^9$ continuous path through 10C space
\item Co-evolution: $d\vec{S}/dt = f(\vec{S}(t), \vec{I}(t), \Psi(t))$
\item Point interventions fail because they ignore temporal dynamics
\item Optimal control: Solve for trajectory $I^*(t)$ that maximizes $\int P_{eff}(t) dt$
\end{enumerate}
Adaptive interventions that ``flow'' through 10C space outperform static ones.
\end{theorem}

\begin{theorem}[MP-T38: Phase-Affinity Matching]
Intervention effectiveness depends on journey-phase matching:
\begin{enumerate}
\item Matching function: $\alpha_{match}(\vec{I}, \varphi) = 1 - \beta \cdot d(\vec{I}, \vec{I}^*(\varphi))$
\item Optimal vector $\vec{I}^*(\varphi)$ varies by region (Pre-Aware: $I_{AWARE}$ high; Stable: $I_{WHO}$ high)
\item Same intervention in wrong phase can have $\alpha_{match} < 0.3$ (poor timing)
\item Interpolation: Between-phase positions handled naturally via continuous $\alpha$
\end{enumerate}
The right intervention at the wrong time is the wrong intervention.
\end{theorem}

\begin{theorem}[MP-T39: Phase-K* Stage Factor]
Optimal intervention count $K^*$ is modulated by journey stage:
\begin{enumerate}
\item $K^* = K_{base} \cdot f_B + \delta_C + \delta_S + \delta_\gamma + \delta_L$ (EXC-3 hybrid)
\item Stage factor $f_S$: Pre-Aware (0.7), Triggered (1.2), Action (1.0), Maintenance (0.8)
\item Triggered phase requires more interventions (critical window)
\item Maintenance phase requires focused portfolio (habit formation)
\end{enumerate}
Phase determines not just \textit{which} interventions but \textit{how many}.
\end{theorem}

\begin{theorem}[MP-T40: Segment-Heterogeneity Dominance]
Treatment Effect Heterogeneity dominates Average Treatment Effects:
\begin{enumerate}
\item $\text{CATE}(s) = E[Y_1 - Y_0 | S = s] \neq \text{ATE}$
\item When $\text{CATE}(s_1) > 0$ and $\text{CATE}(s_2) < 0$, ATE $\approx 0$ is misleading
\item Optimal policy targets segments with $\sigma_s > 1$, avoids $\sigma_s < 0$
\item Simpson's Paradox: Effective intervention appears ineffective in aggregate
\end{enumerate}
One-size-fits-all wastes potential and harms some segments.
\end{theorem}

\begin{theorem}[MP-T41: Phase-Segment Synergy]
Phase and segment effects interact non-additively:
\begin{enumerate}
\item Full model: $E = E_{base} \cdot \alpha_{t,\varphi} \cdot \sigma_s \cdot \delta_{t,s}$
\item Key synergies ($\delta > 1$): Present-Biased$\times$Action (1.4), Social$\times$Triggered (1.5)
\item Key interference ($\delta < 1$): Autonomy$\times$Maintenance (0.7)
\item Optimal intervention timing varies by segment
\end{enumerate}
Intervention planning requires the 2D Phase$\times$Segment surface, not 1D phase alone.
\end{theorem}

\begin{theorem}[MP-T42: Reactance Boundary]
For Autonomy-Oriented segments, intervention visibility determines effectiveness:
\begin{enumerate}
\item Explicit nudges: $\sigma < 0$ (backfire) when $\rho > 0.5$
\item Invisible defaults: $\sigma = 1.5$ (effective)
\item Autonomy cue: $\sigma_{with\_cue} = \sigma_{base} + 0.4\rho$ (mitigation)
\item Self-designed commitment: $\sigma = 1.4$ (respects autonomy)
\end{enumerate}
For high-reactance segments, visibility is the key design dimension.
\end{theorem}

\begin{theorem}[MP-T43: Portfolio Superadditivity]
Coherent portfolios exhibit superadditivity when complementarities are positive:
\begin{enumerate}
\item If $\gamma_{ij} > 0$ for all $i,j \in P$: $E(P) > \sum_i E(\mathcal{I}_i)$
\item Superadditivity magnitude: $\Delta E_{synergy} = \sum_{i<j} \gamma_{ij} \sqrt{E_i E_j}$
\item Crowding-out ($\gamma_{ij} < 0$) causes subadditivity: portfolio worse than sum
\end{enumerate}
Portfolio design is more important than individual intervention selection.
\end{theorem}

\begin{theorem}[MP-T44: K* Optimality]
Portfolio size $K^*$ is optimal in the sense of constrained efficiency:
\begin{enumerate}
\item $|P| < K^*$: Underutilization of budget and opportunities
\item $|P| = K^*$: Optimal balance of diversity and coherence
\item $|P| > K^*$: Diminishing returns, increased coordination costs, potential conflicts
\item K* emergence: $K^* = K_{base} \cdot f_B + \delta_C + \delta_S + \delta_\gamma + \delta_L$ (EXC-3 hybrid)
\end{enumerate}
More interventions is not always better; K* defines the sweet spot.
\end{theorem}

\begin{theorem}[MP-T45: Constraint Hierarchy]
The seven constraints have implicit priority ordering:
\begin{enumerate}
\item C5 (Ethical) and C6 (Segment: no backfire) are hard constraints---violations disqualify portfolio
\item C2 (Crowding) and C7 (K-Optimal) are semi-hard---violations degrade performance significantly
\item C1 (Budget), C3 (Timing), C4 (Capacity) are soft---violations increase cost but may be acceptable
\end{enumerate}
Constraint priority guides trade-off decisions when perfect satisfaction impossible.
\end{theorem}

\begin{theorem}[MP-T46: Hierarchical Cascade]
Optimal portfolio design follows hierarchical cascade:
\begin{enumerate}
\item Strategic level sets $\vec{I}_{weighted}$, segment priorities, budget allocation
\item Tactical level operationalizes into campaigns with K* $\times$ 1.0
\item Operative level implements touchpoints with K* $\times$ 1.2
\item Feedback loop: Operative $\to$ Tactical $\to$ Strategic review
\end{enumerate}
Cross-level coherence requires $\gamma(L_i, L_j) \geq 0$ for all levels.
\end{theorem}

\begin{theorem}[MP-T47: Dynamic K* Evolution]
Portfolio size K*(t) evolves predictably with status quo changes:
\begin{enumerate}
\item Phase progression: $g(\varphi_t) = 1 + 0.3|\Delta\varphi|$ increases K* during transitions
\item Context shock: $h(\Psi_t) = 1 + 0.5 \cdot \mathbb{1}[\text{shock}]$ expands portfolio temporarily
\item Learning convergence: $\ell(\Theta_t) = 1 - 0.2 \cdot \text{converged}$ refocuses stable portfolio
\item Combined: $K^*(t) = K^*_0 \cdot g \cdot h \cdot \ell$
\end{enumerate}
K* is not static---it responds to all three evolution mechanisms.
\end{theorem}

\begin{theorem}[MP-T48: Redesign Trigger Conditions]
Portfolio redesign is triggered when any of three conditions hold:
\begin{enumerate}
\item Phase jump: $|\varphi_t - \varphi_{t-1}| > 0.2$ (individual moved BCJ stages)
\item Context shock: $||\Psi_t - \Psi_{t-1}|| > 0.15$ (policy/economic/social shift)
\item Parameter update: $||\Theta_t - \Theta_{t-1}|| > \delta$ (significant learning)
\end{enumerate}
Triggers are OR-connected: any single trigger suffices for redesign decision.
\end{theorem}

\begin{theorem}[MP-T49: Learning Convergence]
Bayesian parameter learning converges when:
\begin{enumerate}
\item Stability condition: $||\Theta_t - \Theta_{t-1}|| < \epsilon$ for $T$ consecutive periods
\item Practical values: $\epsilon = 0.05$ (5\% threshold), $T = 2$ periods
\item Convergence implies: Stop redesigning, maintain current portfolio $P^*$
\item Typical timeline: 6--12 months (2--4 redesign cycles)
\end{enumerate}
Convergence detection prevents over-optimization and reduces coordination costs.
\end{theorem}

% =============================================================================
\section{Cross-References}
\label{sec:mp-crossref}
% =============================================================================

\textbf{CORE Appendices:}
\begin{itemize}[nosep]
\item AAA (CORE-WHO): Level hierarchy $L$
\item C (CORE-WHAT): FEPSDE dimensions $d$
\item B (CORE-HOW): Complementarity $\gamma$
\item V (CORE-WHEN): Context $\Psi$
\item BBB (CORE-WHERE): Parameters $\Theta$
\item AU (CORE-AWARE): Awareness $A(\cdot)$
\item AV (CORE-READY): Willingness $WAX, \theta$
\item AW (CORE-STAGE): Journey $\varphi$
\item HI (CORE-HIERARCHY): Decision levels $N_{L2}$
\item IE (CORE-EIT): Intervention theory, EMERGE
\end{itemize}

\textbf{FORMAL Appendices:}
\begin{itemize}[nosep]
\item NC (FORMAL-NC): Non-concavity proofs
\item PBB (FORMAL-PROBABILITY): $P_{eff}$ axioms P1-P10
\item UN (FORMAL-UNTCM): Dynamic ODE system
\item D (FORMAL-CSTAR): Reference structure $C^*$
\end{itemize}

\textbf{METHOD Appendices:}
\begin{itemize}[nosep]
\item HHH (METHOD-TOOLKIT): Operational tools, A1-A8, 20-field schema
\item KKK (METHOD-DYNAMIC): Portfolio dynamics
\item SS (CORE-SOLSPACE): Solution space formalization [NEW]
\end{itemize}

% =============================================================================
\section{Summary}
\label{sec:mp-summary}
% =============================================================================

\begin{tcolorbox}[colback=green!5!white, colframe=green!75!black,
    title=\textbf{The Complete EBF Optimization Problem}]

\textbf{Objective:}
\[
\max_{\{I_t\}} \sum_{t=0}^{T} \delta^t \cdot P_{eff}(S_t, I_t, \Psi_t)
\]

\textbf{Where:}
\begin{itemize}[nosep]
\item $P_{eff} = \sigma(WEC \cdot \alpha_{BCJ} \cdot \beta_{BCS})$ is the probability of behavioral change (Ch16)
\item $S = (L, d, \gamma, \Psi, \Theta, A, W, \varphi, N_{L2})$ is the 10C state (Ch9-15)
\item $I \in [0,1]^9$ is the intervention vector (Ch17, IE)
\item $\mathcal{S}(\Psi, S)$ is the solution space (SS)
\end{itemize}

\textbf{Subject to:}
\begin{itemize}[nosep]
\item State dynamics: $S_{t+1} = S_t + \Delta S(I_t)$
\item Context evolution: $\Psi_{t+1} = f(\Psi_t, I_t, a_t)$
\item Hard constraints: C1-C7 (Ch20)
\item Structural constraints: Non-concavity, K-Q distinction (Ch7)
\item Multi-level coherence: $\gamma(L_i, L_j) \geq 0$ (Ch23)
\end{itemize}

\textbf{Key Insights:}
\begin{enumerate}
\item The landscape has multiple local optima (not convex optimization)
\item The solution space emerges from context (endogenous)
\item Intervention types are emergent clusters, not primitives
\item The system is path-dependent
\item All 24 chapters contribute essential components
\end{enumerate}

\end{tcolorbox}

% =============================================================================
\section{Glossary}
\label{sec:mp-glossary}
% =============================================================================

This glossary provides comprehensive definitions of all key terms, symbols, and concepts
used in the MP CORE-MASTER appendix, organized by category.

% -----------------------------------------------------------------------------
\subsection{Symbol Disambiguation: Greek Letter Overloading}
\label{sec:mp-symbol-disambiguation}
% -----------------------------------------------------------------------------

\begin{tcolorbox}[colback=yellow!5!white, colframe=yellow!75!black,
    title=\textbf{IMPORTANT: Symbol Disambiguation}]

Several Greek letters are used with \textbf{different meanings} in different contexts.
This table disambiguates overloaded symbols:

\begin{center}
\renewcommand{\arraystretch}{1.3}
\begin{tabular}{@{}llll@{}}
\toprule
\textbf{Symbol} & \textbf{Meaning} & \textbf{Context} & \textbf{Reference} \\
\midrule
\multicolumn{4}{l}{\textit{$\alpha$ (alpha) --- 6 different meanings:}} \\
$\alpha_i$ & Fehr-Schmidt envy parameter & Social preferences & Ch5, $\approx 0.85$ \\
$\alpha_{BCJ}$ & Journey stage factor & $P_{eff}$ equation & Ch13, Ch16 \\
$\alpha^L(\Psi)$ & Level salience weight & WHO aggregation & Ch2, $\sum_L = 1$ \\
$\alpha_{t,\varphi}$ & Phase-type affinity & Intervention effectiveness & Ch18 \\
$\alpha_d$ & Domain absorption rate & ODE parameters & Ch24 \\
$\alpha$ (in $N_{L2}$) & Bounded rationality constant & Hierarchy formula & Ch15, $\approx 0.8$ \\
\midrule
\multicolumn{4}{l}{\textit{$\beta$ (beta) --- 4 different meanings:}} \\
$\beta_i$ & Fehr-Schmidt guilt parameter & Social preferences & Ch5, $\approx 0.315$ \\
$\beta_{BCS}$ & Segment factor & $P_{eff}$ equation & Ch14, Ch16 \\
$\beta$ (time) & Present bias parameter & Quasi-hyperbolic & Ch3, $\approx 0.7$ \\
$\beta_d$ & Domain decay rate & ODE parameters & Ch24 \\
\midrule
\multicolumn{4}{l}{\textit{$\varphi$ vs $\tau$ --- resolved:}} \\
$\varphi$ & BCJ Phase (journey stage) & STAGE dimension & Ch13 \\
$\tau$ & Thinking Style parameter & READY dimension & Ch12 \\
\bottomrule
\end{tabular}
\end{center}

\textbf{Convention:} When $\alpha$ appears without subscript in a formula, check the equation label
or surrounding context to determine which meaning applies.

\vspace{0.5em}
\hrule
\vspace{0.5em}

\textbf{$L$ (Level) --- 3 different hierarchies:}

\begin{center}
\renewcommand{\arraystretch}{1.2}
\begin{tabular}{@{}llll@{}}
\toprule
\textbf{Hierarchy} & \textbf{Notation} & \textbf{Values} & \textbf{Chapter} \\
\midrule
\textbf{WHO (Welfare)} & $L$ & Individual(1), Dyad(2), Group(3), Society(4) & Ch2 \\
\textbf{Decision} & $L_0$--$L_3$ & Meta, Strategic Roots, Derived, Operative & Ch15 \\
\textbf{Scope (Org)} & $L_0$--$L_4$ & Individual, Household, Org, National, Global & Ch23 \\
\bottomrule
\end{tabular}
\end{center}

\textbf{In State Vector:} $S = (L, d, \gamma, \Psi, \Theta, A, W, \varphi, N_{L2})$
\begin{itemize}[nosep]
    \item $L$ = WHO Level (welfare aggregation: 1--4)
    \item $N_{L2}$ = Number of L2 (derived strategic) decisions from HIERARCHY
\end{itemize}

\textbf{Convention:} Use subscript to disambiguate: $L_{WHO}$ (welfare), $L_{dec}$ (decision), $L_{scope}$ (organizational).

\vspace{0.5em}
\hrule
\vspace{0.5em}

\textbf{$\vec{I}$ (Intervention Vector) --- 9D within 10C:}

\begin{center}
\renewcommand{\arraystretch}{1.2}
\begin{tabular}{@{}lll@{}}
\toprule
\textbf{Framework} & \textbf{Dimension} & \textbf{Explanation} \\
\midrule
\textbf{10C} & 10 & The theoretical framework (10 CORE questions) \\
\textbf{$\vec{I}$} & 9 & The intervention vector $\vec{I} \in [0,1]^9$ \\
\bottomrule
\end{tabular}
\end{center}

\textbf{Why 9D in a 10C Framework?}
\begin{itemize}[nosep]
    \item The vector $\vec{I} = (I_{WHO}, I_{WHAT}, I_{HOW}, I_{WHEN}, I_{WHERE}, I_{AWARE}, I_{READY}, I_{STAGE}, I_{HIER})$ targets \textbf{COREs 1--9}
    \item \textbf{CORE 10 (EIT)} is the \textit{methodology} for generating interventions---not a target dimension itself
    \item Analogy: EIT is the ``how to design'' (process), COREs 1--9 are the ``what to change'' (targets)
\end{itemize}

\textbf{Reference:} Appendix UNMAPPED_IE (CORE-EIT), core-framework-definition.yaml (stage\_6)

\end{tcolorbox}

% -----------------------------------------------------------------------------
\subsection{Core Variables and Functions}
% -----------------------------------------------------------------------------

\begin{longtable}{@{}p{2.8cm}p{10cm}@{}}
\toprule
\textbf{Symbol} & \textbf{Definition} \\
\midrule
\endhead
$P_{eff}$ & \textbf{Effective probability of behavioral change}---the objective function of EBF. Computed as $P_{eff} = \sigma(WEC \cdot \alpha_{BCJ} \cdot \beta_{BCS})$ where $\sigma$ is the sigmoid function. \\[6pt]
$S$ & \textbf{10C Status Quo Vector}---the complete state vector comprising all 10C dimensions: $S = (L, d, \gamma, \Psi, \Theta, A, W, \varphi, N_{L2})$. \\[6pt]
$\vec{I}$ & \textbf{Intervention Vector}---a 9-dimensional continuous vector $\vec{I} \in [0,1]^9$ specifying intervention intensity on COREs 1--9. (CORE 10/EIT is the methodology, not a target---see Section~\ref{sec:mp-symbol-disambiguation}.) \\[6pt]
$\Psi$ & \textbf{Context Vector}---an 8-dimensional vector $\Psi \in [0,1]^8$ capturing the environmental conditions that modulate all behavioral parameters. \\[6pt]
$\mathcal{S}(\Psi, S)$ & \textbf{Solution Space}---the set of feasible interventions given context $\Psi$ and state $S$: $\mathcal{S}(\Psi, S) = \{I \in [0,1]^9 : \text{feasible}(I, \Psi, S)\}$. \\[6pt]
$\Theta$ & \textbf{Parameter Set}---all calibrated parameters from 1,900+ empirical studies (Appendix UNMAPPED_BBB). \\[6pt]
$V(S, \Psi)$ & \textbf{Value Function}---the Bellman value function representing optimal expected future welfare from state $(S, \Psi)$. \\[6pt]
$\delta$ & \textbf{Discount Factor}---time preference parameter; may be dimension-specific $\delta_d(\Psi)$. \\[6pt]
$\sigma(\cdot)$ & \textbf{Sigmoid Function}---$\sigma(x) = 1/(1+e^{-kx})$, transforms WEC to probability. \\
\bottomrule
\end{longtable}

% -----------------------------------------------------------------------------
\subsection{The 10C Framework Dimensions}
% -----------------------------------------------------------------------------

\begin{longtable}{@{}p{2.8cm}p{10cm}@{}}
\toprule
\textbf{10C Dimension} & \textbf{Definition} \\
\midrule
\endhead
\textbf{WHO} (AAA) & \textbf{Who has utility?}---The welfare hierarchy from $L_0$ (self) to $L_\infty$ (all sentient beings). Determines aggregation level. \\[6pt]
\textbf{WHAT} (C) & \textbf{What is utility?}---The six FEPSDE dimensions (Financial, Emotional, Physical, Social, Development, Ecological) plus three categories (INU, KNU, IDN). \\[6pt]
\textbf{HOW} (B) & \textbf{How do utility components interact?}---The complementarity structure $\gamma_{ij}$ determining whether components are complements, substitutes, or independent. \\[6pt]
\textbf{WHEN} (V) & \textbf{When does context matter?}---The 8 context dimensions $\Psi$ that modulate all behavioral parameters. \\[6pt]
\textbf{WHERE} (BBB) & \textbf{Where do parameters come from?}---The empirical calibration from 1,900+ papers providing functional knowledge. \\[6pt]
\textbf{AWARE} (AU) & \textbf{How aware am I of utility at the Moment of Truth?}---The awareness function $A(\cdot) \in [0,1]$ that filters potential into effective utility. \\[6pt]
\textbf{READY} (AV) & \textbf{How ready am I to act?}---The willingness-to-act index $WAX$ and action threshold $\theta$. \\[6pt]
\textbf{STAGE} (AW) & \textbf{Where am I in the behavioral journey?}---The BCJ phase $\varphi \in \{S_1, ..., S_6\}$ (PreAware to Maintenance). \\[6pt]
\textbf{HIERARCHY} (HI) & \textbf{What decision level?}---The hierarchical level $N_{L2} \in \{L_0, L_1, L_2, L_3\}$ affecting cognitive processing. \\[6pt]
\textbf{EIT} (IE) & \textbf{How do interventions emerge?}---The \textit{methodology} for deriving intervention vectors $\vec{I} \in [0,1]^9$. \textbf{Note:} EIT is not a target dimension---it generates interventions that target COREs 1--9. \\
\bottomrule
\end{longtable}

% -----------------------------------------------------------------------------
\subsection{Utility and Welfare Structure}
% -----------------------------------------------------------------------------

\begin{longtable}{@{}p{2.8cm}p{10cm}@{}}
\toprule
\textbf{Symbol} & \textbf{Definition} \\
\midrule
\endhead
\textbf{FEPSDE} & \textbf{Six utility dimensions}: Financial ($U^F$), Emotional ($U^E$), Physical ($U^P$), Social ($U^S$), Development ($U^D$), Ecological ($U^{Eco}$). \\[6pt]
\textbf{INU} & \textbf{Individual Utility}---``What do I want for myself?'' Material self-interest. \\[6pt]
\textbf{KNU} & \textbf{Kollective Utility}---``What do I want for us?'' Welfare of reference groups, level-indexed $KNU^L$. \\[6pt]
\textbf{IDN} & \textbf{Identity Utility}---``Who am I / want to be?'' Self-concept and values. \\[6pt]
$U_{pot}$ & \textbf{Potential Utility}---the 144-component utility structure before awareness filtering. \\[6pt]
$U_{eff}$ & \textbf{Effective Utility}---$U_{eff} = A(t^*) \times U_{pot}$; utility that actually influences behavior. \\[6pt]
$K$ & \textbf{Coherence Index}---measures whether a system is at an equilibrium: $K = 1 - \|C - C^*\|/\|C^*\|$. High $K$ = ``at a peak''. \\[6pt]
$Q$ & \textbf{Quality Index}---measures the welfare value of an equilibrium: $Q = \Pi(C^*; \Psi)$. High $Q$ = ``high peak''. \\[6pt]
$K$ vs $Q$ & \textbf{Central distinction}: $K \approx 1$ means coherent (stable), but does NOT imply high welfare. A society can be stably trapped at a low-$Q$ equilibrium. \\[6pt]
$\lambda_d$ & \textbf{Loss Aversion}---dimension-specific coefficient where losses loom larger than gains: $U_d = G_d - \lambda_d \cdot P_d$ with $\lambda \approx 2$. \\[6pt]
$G$, $P$ & \textbf{Gain and Pain}---the two valences of utility. 144 components = $3 \times 2 \times 6 \times 4$. \\[6pt]
$\omega_i$ & \textbf{Category Weight}---relative importance of INU, KNU, IDN in the aggregate quality function. \\
\bottomrule
\end{longtable}

% -----------------------------------------------------------------------------
\subsection{Complementarity Structure}
% -----------------------------------------------------------------------------

\begin{longtable}{@{}p{2.8cm}p{10cm}@{}}
\toprule
\textbf{Symbol} & \textbf{Definition} \\
\midrule
\endhead
$\gamma_{ij}$ & \textbf{Complementarity Coefficient}---cross-partial derivative $\partial^2 U / \partial x_i \partial x_j$. Positive = complements, negative = substitutes, zero = independent. \\[6pt]
$\Gamma$ & \textbf{FEPSDE Complementarity Matrix}---15 parameters capturing interactions between utility dimensions. \\[6pt]
$\Lambda$ & \textbf{Context Interaction Matrix}---28 parameters capturing $\Psi_k \times \Psi_{k'}$ interactions. \\[6pt]
$\delta_{dd',k}$ & \textbf{Modulation Tensor}---120 parameters capturing how context $k$ modulates complementarity between dimensions $d$ and $d'$. \\[6pt]
$C$ & \textbf{Complementarity Matrix}---full $144 \times 144$ matrix of utility component interactions. \\[6pt]
$C^*$ & \textbf{Reference Structure}---the self-consistent fixed point $C^* = \Phi(C^*, \Psi)$ against which coherence is measured. \\[6pt]
$C_{ij} \neq 0$ & \textbf{Classical Violation}---EBF relaxes Arrow-Debreu's separability assumption to allow utility interdependence. \\[6pt]
Supermodularity & A function is \textbf{supermodular} if $\partial^2 f / \partial x_i \partial x_j \geq 0$ for all $i \neq j$; creates multiple equilibria. \\
\bottomrule
\end{longtable}

% -----------------------------------------------------------------------------
\subsection{Context Dimensions ($\Psi$)}
% -----------------------------------------------------------------------------

\begin{longtable}{@{}p{2.8cm}p{10cm}@{}}
\toprule
\textbf{Symbol} & \textbf{Definition} \\
\midrule
\endhead
$\Psi_I$ & \textbf{Institutional Context}---formal rules, property rights, rule of law (very slow change). \\[6pt]
$\Psi_S$ & \textbf{Social Context}---trust, norms, social capital (medium change speed). \\[6pt]
$\Psi_C$ & \textbf{Cognitive Context}---bounded rationality, attention capacity, financial literacy (fast change). \\[6pt]
$\Psi_K$ & \textbf{Cultural Context}---values, worldviews, belief systems (very slow change). \\[6pt]
$\Psi_E$ & \textbf{Economic Context}---markets, prices, income, development level (fast change). \\[6pt]
$\Psi_T$ & \textbf{Temporal Context}---time pressure, stability, planning horizons (variable speed). \\[6pt]
$\Psi_M$ & \textbf{Material Context}---technology, infrastructure (medium change speed). \\[6pt]
$\Psi_F$ & \textbf{Physical Context}---geography, climate, factor flexibility (very slow change). \\[6pt]
$\eta$ & \textbf{Context Adjustment Speed}---rate at which context dimension responds to behavior. \\[6pt]
Endogenous $\Psi$ & Context evolves through behavior: $\Psi_{t+1} = f(\Psi_t, I_t, a_t)$; creates feedback loops. \\
\bottomrule
\end{longtable}

% -----------------------------------------------------------------------------
\subsection{Awareness and Willingness (AWARE, READY)}
% -----------------------------------------------------------------------------

\begin{longtable}{@{}p{2.8cm}p{10cm}@{}}
\toprule
\textbf{Symbol} & \textbf{Definition} \\
\midrule
\endhead
$A(\cdot)$ & \textbf{Awareness Function}---$A_{i,d,t}(t^*) \in [0,1]$ measuring salience of utility component at Moment of Truth. \\[6pt]
$t^*$ & \textbf{Moment of Truth}---the decision point where awareness configuration determines behavior. \\[6pt]
$AWX$ & \textbf{Awareness-Weighted Utility}---utility weighted by awareness: $U^* = A \times U$. \\[6pt]
SYS 0/1/2 & \textbf{Three Processing Levels}: SYS 0 (Instinct, hardwired), SYS 1 (Intuition, fast), SYS 2 (Deliberation, slow). \\[6pt]
$WAX$ & \textbf{Willingness-to-Act Index}---master equation: $WAX = \tau \cdot \min\{U^*\} + (1-\tau) \cdot \sum U^*$. \\[6pt]
$\theta$ & \textbf{Action Threshold}---context-dependent barrier: $\theta = \theta_{base}(d) + \sum_j \delta_j \Psi_j$. Action occurs iff $WAX \geq \theta$. \\[6pt]
$\tau$ (tau) & \textbf{Thinking Style Parameter}---$\tau \in [0,1]$. High $\tau$ = connected/complementary thinking (``only act if everything is right''); low $\tau$ = separated/additive thinking. \textit{Note: $\tau$ (Thinking Style) $\neq$ $\varphi$ (BCJ Phase); see $\varphi$ in STAGE glossary below.} \\[6pt]
$WEC$ & \textbf{Willingness-Exceeds-Cost}---$WEC = WAX - \theta$; the gap that determines action probability. \\[6pt]
$\Delta$-Gap & \textbf{Knowing-Doing Gap}---difference between potential ($U_{pot}$) and effective ($U_{eff}$) utility due to awareness filtering. \\
\bottomrule
\end{longtable}

% -----------------------------------------------------------------------------
\subsection{Behavioral Journey (STAGE) and Segments (BCS)}
% -----------------------------------------------------------------------------

\begin{longtable}{@{}p{2.8cm}p{10cm}@{}}
\toprule
\textbf{Symbol} & \textbf{Definition} \\
\midrule
\endhead
\textbf{BCJ} & \textbf{Behavioral Change Journey}---the 6-phase progression from PreAware to Maintenance. \\[6pt]
$\varphi$ (phase) & \textbf{BCJ Phase}---current stage in journey: $S_1$ (PreAware), $S_2$ (Triggered), $S_3$ (Seeking), $S_4$ (Acting), $S_5$ (Habituating), $S_6$ (Maintenance). \textit{Note: $\varphi$ (BCJ Phase) $\neq$ $\tau$ (Thinking Style); see $\tau$ in READY glossary above.} \\[6pt]
$\alpha_{BCJ}$ & \textbf{Journey Stage Factor}---phase-dependent multiplier in $P_{eff}$ equation. Range $(0, 1]$. \textit{Note: This $\alpha$ is distinct from $\alpha_i$ (Fehr-Schmidt) and $\alpha^L$ (level weights).} \\[6pt]
$f_S$ & \textbf{Stage Factor}---adjustment to K* based on BCJ phase. \\[6pt]
$T_{lock}$ & \textbf{Lock-in Time}---minimum time required in each phase before transition. \\[6pt]
\textbf{BCS} & \textbf{Behavioral Change Segments}---population heterogeneity clusters based on behavioral parameters. \\[6pt]
$\beta_{BCS}$ & \textbf{Segment Factor}---segment-dependent multiplier in $P_{eff}$ equation. \\[6pt]
$\sigma_s$ & \textbf{Segment Multiplier}---scales base effectiveness for segment $s$: $\sigma_s = CATE_s / ATE$. Range $[-0.5, 2.0]$. \\[6pt]
$\alpha_i$ (Fehr-Schmidt) & \textbf{Envy Parameter}---sensitivity to disadvantageous inequity in social preferences, mean $\approx 0.85$, range $[0, 4+]$. \textit{Literature-standard symbol; distinct from $\alpha_{BCJ}$ and $\alpha^L$.} \\[6pt]
$\beta_i$ (social) & \textbf{Social Preference (Fehr-Schmidt)}---sensitivity to advantageous inequity (guilt), mean $\approx 0.315$. \\[6pt]
$\rho$ & \textbf{Reactance Parameter}---sensitivity to perceived manipulation, $\rho \in [0,1]$. \\[6pt]
$\delta_{t,s}$ & \textbf{Phase-Segment Interaction}---coefficient capturing synergy/interference between phase and segment. \\
\bottomrule
\end{longtable}

% -----------------------------------------------------------------------------
\subsection{Intervention Types and Phase-Affinity}
% -----------------------------------------------------------------------------

\begin{longtable}{@{}p{2.8cm}p{10cm}@{}}
\toprule
\textbf{Symbol} & \textbf{Definition} \\
\midrule
\endhead
$\alpha_{t,\varphi}$ & \textbf{Phase-Type Affinity}---effectiveness multiplier for intervention type $t$ in phase $\varphi$. Values: $\alpha > 0.7$ (recommended), $\alpha < 0.3$ (not recommended). \\[6pt]
\textbf{A1--A8} & \textbf{Heuristic Intervention Clusters}: A1 (Awareness Boost), A2 (Feedback Loop), A3 (Choice Architecture), A4 (Temporal Push), A5 (Identity Framing), A6 (Social Proof), A7 (Financial Incentive), A8 (Commitment Device). \\[6pt]
EIT & \textbf{Emergent Intervention Theory} (CORE 10)---the \textit{methodology} for designing interventions. EIT generates vectors $\vec{I} \in [0,1]^9$ that target COREs 1--9. Note: EIT is a process, not a target dimension. \\[6pt]
EMERGE & \textbf{Algorithm}---computes optimal intervention profile: $\vec{I}^* = \text{EMERGE}(\Psi, \Gamma, B, K^*)$. \\
\bottomrule
\end{longtable}

% -----------------------------------------------------------------------------
\subsection{Portfolio Optimization}
% -----------------------------------------------------------------------------

\begin{longtable}{@{}p{2.8cm}p{10cm}@{}}
\toprule
\textbf{Symbol} & \textbf{Definition} \\
\midrule
\endhead
$K^*$ & \textbf{Optimal Portfolio Size}---number of simultaneous interventions. Computed as $K^* = K_{base} \cdot f_B + \delta_C + \delta_S + \delta_\gamma + \delta_L$ (EXC-3 hybrid). \\[6pt]
$f_B$ & \textbf{Budget Factor}---$f_B = \sqrt{B/B_{ref}}$, adjusts K* for available resources. \\[6pt]
$f_C$ & \textbf{Complexity Factor}---$f_C = 1/\sqrt{\text{complexity}}$, reduces K* for complex populations. \\[6pt]
$f_\gamma$ & \textbf{Complementarity Factor}---$f_\gamma = 1 - |\bar{\gamma}_{neg}|$, reduces K* when crowding-out risk is high. \\[6pt]
$f_L$ & \textbf{Scope Level Factor}---adjusts K* by organizational level: L0 ($\times 0.6$), L4 ($\times 1.5$). \\[6pt]
\textbf{C1--C7} & \textbf{Seven Constraint Classes}: C1 (Budget), C2 (Crowding-Out), C3 (Timing), C4 (Capacity), C5 (Ethical/Autonomy), C6 (Segment Safety), C7 (K-Optimal). \\[6pt]
Crowding-Out & Negative complementarity between intervention types: Social + Financial ($\gamma = -0.3$), Financial + Intrinsic ($\gamma = -0.3$). \\[6pt]
$Z(n)$ & \textbf{Choice Satisfaction}---$Z(n) = Z_{max} \times (1 - 0.03 \times \max(0, n-7))$; satisfaction declines beyond 7 options. \\
\bottomrule
\end{longtable}

% -----------------------------------------------------------------------------
\subsection{Dynamic Optimization}
% -----------------------------------------------------------------------------

\begin{longtable}{@{}p{2.8cm}p{10cm}@{}}
\toprule
\textbf{Symbol} & \textbf{Definition} \\
\midrule
\endhead
$K^*(t)$ & \textbf{Dynamic Portfolio Size}---$K^*(t) = K^*_0 \cdot g(\varphi_t) \cdot h(\Psi_t) \cdot \ell(\Theta_t)$; evolves with phase, context, and learning. \\[6pt]
$d\varphi/dt$ & \textbf{Phase Velocity}---rate of BCJ progression: $d\varphi/dt = v_0 + v_A \cdot A_t + v_W \cdot W_t$. \\[6pt]
Redesign Trigger & Conditions requiring portfolio update: $|\Delta\varphi| > 0.2$ (phase jump), $||\Delta\Psi|| > 0.15$ (context shock), $||\Delta\Theta|| < \epsilon$ (convergence). \\[6pt]
Convergence & Learning has stabilized when $||\Theta_t - \Theta_{t-1}|| < \epsilon$ for $T$ consecutive periods. \\[6pt]
$\epsilon$ & \textbf{Convergence Threshold}---typically 0.05 (5\% parameter stability). \\
\bottomrule
\end{longtable}

% -----------------------------------------------------------------------------
\subsection{Policy and Multi-Level Implementation}
% -----------------------------------------------------------------------------

\begin{longtable}{@{}p{2.8cm}p{10cm}@{}}
\toprule
\textbf{Symbol} & \textbf{Definition} \\
\midrule
\endhead
\textbf{L0--L4} & \textbf{Five Organizational Levels}: L0 (Individual), L1 (Household), L2 (Organization), L3 (National), L4 (Global). \\[6pt]
Scale Invariance & The 9C diagnostic and portfolio optimization applies identically at all levels; only parameters vary. \\[6pt]
$K^*(L)$ & \textbf{Level-Specific K*}---varies by level: L0 (2--4), L2 (5--8), L4 (3--5). \\[6pt]
$\gamma(L_i, L_j)$ & \textbf{Cross-Level Complementarity}---must be $\geq 0$; negative values cause implementation failure. \\[6pt]
$\kappa_{j \to i}$ & \textbf{Top-Down Spillover}---factor by which level $j$ policies affect level $i$ behavior. \\[6pt]
$\tau_i$ & \textbf{Cascade Fidelity}---transmission accuracy at level $i$. If $\tau_i < 0.5$, message reverses. \\[6pt]
$\tau_{total}$ & \textbf{End-to-End Fidelity}---$\tau_{total} = \prod_{i=1}^{n-1} \tau_i$; 4-level cascade with $\tau_i = 0.8$ yields $\tau_{total} = 0.41$. \\[6pt]
$\kappa(\Psi)$ & \textbf{Crisis Adjustment Factor}---$\kappa \in [0.4, 1.0]$ during context shocks. WHAT$_F$ increases, WHEN decreases. \\[6pt]
Coherence Score & \textbf{Policy Coherence}---$\text{Coherence} = \sum_{i<j} \gamma_{ij} / \binom{n}{2}$; target $> 0.25$. \\[6pt]
Hybrid Optimality & Top-down goals + bottom-up execution outperforms pure strategies by $\approx 20\%$. \\
\bottomrule
\end{longtable}

% -----------------------------------------------------------------------------
\subsection{Emergent Life Journeys}
% -----------------------------------------------------------------------------

\begin{longtable}{@{}p{2.8cm}p{10cm}@{}}
\toprule
\textbf{Symbol} & \textbf{Definition} \\
\midrule
\endhead
Seven Journeys & \textbf{Emergent Life Domains}: Health, Money, Career, Relationships, Self-Governance, Living Space, Meaning/Legacy. \\[6pt]
Domain ODE & $dA/dt = \alpha_d \cdot T_1(t) \cdot (1-A) - \beta_d \cdot A$; domain-specific dynamics. \\[6pt]
$\alpha_d$ & \textbf{Acquisition Rate}---speed of awareness gain in domain $d$. \\[6pt]
$\beta_d$ & \textbf{Decay Rate}---speed of awareness loss in domain $d$. \\[6pt]
$\lambda_d$ (journey) & \textbf{Learning Rate}---speed of incorporating observed effects. \\[6pt]
$\Gamma^{inter}$ & \textbf{Intergenerational Spillover Matrix}---coefficients for parent-to-child behavioral transmission. \\[6pt]
$\gamma^{inter}_{d',d}$ & \textbf{Spillover Coefficient}---effect of parent domain $d'$ on child domain $d$. Self-Governance: $\gamma^{inter}_{SG,SG} = 0.60$. \\[6pt]
\textbf{CTW} & \textbf{Critical Transmission Window}---age period when parental influence is strongest: Health (0--6y), Self-Gov (3--12y), Money (12--22y). \\[6pt]
1.70× Multiplier & \textbf{Self-Governance Effect}---every unit increase in parental Self-Governance produces 1.70 units aggregate child benefit. \\[6pt]
$ROI_{family}$ & \textbf{Family-Level ROI}---$ROI_{family} \approx 1.7 \times ROI_{individual}$ during CTWs. \\[6pt]
$f_{optimal}(\phi)$ & \textbf{Optimal Intervention Frequency}---decreases as $\phi \to 1$: PreAware (1/year), Action (1/week), Maintenance (1/quarter). \\
\bottomrule
\end{longtable}

% -----------------------------------------------------------------------------
\subsection{Mathematical and Structural Concepts}
% -----------------------------------------------------------------------------

\begin{longtable}{@{}p{2.8cm}p{10cm}@{}}
\toprule
\textbf{Term} & \textbf{Definition} \\
\midrule
\endhead
Non-Concavity & The optimization landscape has multiple local maxima due to complementarity; requires global search. \\[6pt]
Valley & Incoherent configuration between peaks; transition requires crossing through reduced performance. \\[6pt]
Peak & Coherent equilibrium with $K \approx 1$; may have high or low $Q$. \\[6pt]
Path Dependence & Which equilibrium is reached depends on initial conditions and trajectory. \\[6pt]
Fixed Point & $C^* = \Phi(C^*, \Psi)$; self-consistent reference structure. \\[6pt]
Bellman Equation & $V(S_t, \Psi_t) = \max_{I_t} \{P_{eff}(S_t, I_t, \Psi_t) + \delta \cdot V(S_{t+1}, \Psi_{t+1})\}$. \\[6pt]
$\beta$-$\delta$ Model & \textbf{Quasi-Hyperbolic Discounting}---$U_t = u(c_t) + \beta \sum_{\tau=1}^T \delta^\tau u(c_{t+\tau})$ with $\beta < 1$ (present bias). \\[6pt]
QRE & \textbf{Quantal Response Equilibrium}---$P(s_i) = \exp(\lambda \pi_i) / \sum_j \exp(\lambda \pi_j)$; $\lambda$ parameterizes rationality. \\[6pt]
ESS & \textbf{Evolutionarily Stable Strategy}---strategy that cannot be invaded by mutants. \\[6pt]
Replicator Dynamics & $\dot{x}_i = x_i[\pi_i(x) - \bar{\pi}(x)]$; evolutionary selection mechanism. \\[6pt]
$d(S_t, S^*)$ & \textbf{Distance to Target}---metric measuring progress toward goal state. \\[6pt]
$\mathcal{R}(S_0, T)$ & \textbf{Reachable Set}---all states achievable from $S_0$ in $T$ steps. \\
\bottomrule
\end{longtable}

% -----------------------------------------------------------------------------
\subsection{Key Acronyms and Abbreviations}
% -----------------------------------------------------------------------------

\begin{longtable}{@{}p{2.8cm}p{10cm}@{}}
\toprule
\textbf{Acronym} & \textbf{Full Form} \\
\midrule
\endhead
EBF & Evidence-Based Framework for Economic and Social Behavior \\
10C & Ten Core Questions (WHO, WHAT, HOW, WHEN, WHERE, AWARE, READY, STAGE, HIERARCHY + combined output) \\
9C & Nine-dimensional state space (pre-HIERARCHY) \\
BCJ & Behavioral Change Journey \\
BCS & Behavioral Change Segments \\
FEPSDE & Financial, Emotional, Physical, Social, Development, Ecological \\
WAX & Willingness-to-Act Index \\
WEC & Willingness Exceeds Cost \\
EIT & Emergent Intervention Theory (CORE 10, methodology) \\
CTW & Critical Transmission Window \\
ODE & Ordinary Differential Equation \\
ROI & Return on Investment \\
ATE & Average Treatment Effect \\
CATE & Conditional Average Treatment Effect \\
WEIRD & Western, Educated, Industrialized, Rich, Democratic (sampling bias) \\
\bottomrule
\end{longtable}

\vspace{1em}
\textit{For complete symbol definitions across the entire EBF framework, see Appendix G (REF-GLOSSARY).}

% =============================================================================
\section{Theory: The Optimization Framework}
\label{sec:mp-theory}
% =============================================================================

The EBF optimization framework rests on three theoretical pillars:

\subsection{Pillar 1: Behavioral Probability as Objective}

Unlike classical welfare economics (maximize utility $U$), EBF optimizes the
\textbf{probability of behavioral change} $P_{eff}$. This reflects:
\begin{itemize}
\item Behavioral change is uncertain, not deterministic
\item The sigmoid transformation maps WEC to probabilities
\item Population-level effects aggregate individual probabilities
\end{itemize}

\subsection{Pillar 2: Endogenous Context and Solution Space}

Context $\Psi$ is not exogenous but evolves through behavior (Chapter 9):
\[
\Psi_{t+1} = f(\Psi_t, I_t, a_t)
\]
This creates a feedback loop where interventions change what interventions
become feasible---the solution space $\mathcal{S}(\Psi)$ emerges from context.

\subsection{Pillar 3: Non-Convex Landscape with Multiple Equilibria}

Complementarity ($\gamma > 0$) generates non-concavity (Chapter 7, Appendix UNMAPPED_NC):
\begin{itemize}
\item Multiple local optima (peaks) exist
\item Coherence ($K \approx 1$) $\neq$ optimality ($\max Q$)
\item Transitions require coordinated, simultaneous change
\end{itemize}

% =============================================================================
\section{Results: Key Theoretical Contributions}
\label{sec:mp-results}
% =============================================================================

\subsection{Result 1: Complete Problem Specification}

This appendix provides the first \textbf{complete formal specification} of the
EBF optimization problem, integrating all 24 chapters into equations
(\ref{eq:mp-objective})--(\ref{eq:mp-bellman}).

\subsection{Result 2: Chapter-Component Mapping}

Table in Section \ref{sec:mp-chapters} establishes the precise contribution
of each chapter to the optimization problem, enabling:
\begin{itemize}
\item Systematic verification of framework completeness
\item Clear dependencies between components
\item Identification of gaps requiring new chapters/appendices
\end{itemize}

\subsection{Result 3: Dynamic Bellman Formulation}

Equation (\ref{eq:mp-bellman}) casts EBF as a dynamic programming problem,
connecting to established optimal control theory while highlighting
EBF-specific complications (non-concavity, endogenous state space).

% =============================================================================
\section{Implications}
\label{sec:mp-implications}
% =============================================================================

\subsection{For Practitioners}

\begin{itemize}
\item \textbf{Diagnosis first}: Assess all 10C dimensions before intervention design
\item \textbf{Portfolio thinking}: Single interventions rarely optimal; use $K^*$ formula
\item \textbf{Context matters}: Same intervention works differently in different $\Psi$
\item \textbf{Watch for valleys}: Partial change may make things worse (Chapter 7)
\end{itemize}

\subsection{For Researchers}

\begin{itemize}
\item \textbf{Testable predictions}: Each component generates falsifiable hypotheses
\item \textbf{Parameter estimation}: Need calibration for $\gamma$, $\alpha$, $\beta$, etc.
\item \textbf{Solution algorithms}: Non-convex optimization requires specialized methods
\item \textbf{Empirical validation}: Compare predicted vs. observed $P_{eff}$
\end{itemize}

\subsection{For Policy Makers}

\begin{itemize}
\item \textbf{Multi-level coherence}: Policies at L0-L4 must not undermine each other
\item \textbf{Long-term dynamics}: Consider 30-year trajectories (Chapter 24)
\item \textbf{Domain spillovers}: Interventions in one domain affect others
\end{itemize}

% =============================================================================
\section{Critical Foundations}
\label{sec:mp-foundations}
% =============================================================================

This appendix builds on the following foundational works:

\begin{enumerate}
\item \textbf{Classical Foundations}: \citet{arrowdeb} on general equilibrium;
      \citet{debreu1959} on axiomatic foundations; \citet{sonnenschein1972} on
      aggregate demand restrictions

\item \textbf{Evolutionary Game Theory}: \citet{maynardsmith1982} on ESS;
      \citet{weibull1995} on evolutionary dynamics; \citet{gintis2009} on bounded
      rationality as evolutionary adaptation

\item \textbf{Bounded Rationality}: \citet{mckelvey1995} on QRE;
      \citet{smith1776, smith1759} on dual nature of economic agents

\item \textbf{Complementarity Economics}: \citet{milgrom1990} on supermodularity;
      \citet{roberts2004} on organizational fit

\item \textbf{Behavioral Economics}: \citet{kahneman1979prospect} on prospect theory;
      \citet{thaler2008} on nudging; \citet{fehr1999} on social preferences;
      \citet{fehrschmidt1999} on inequality aversion

\item \textbf{Intertemporal Choice}: \citet{laibson1997} on $\beta$-$\delta$ discounting;
      \citet{odonoghue1999doing} on sophistication vs. naïveté;
      \citet{frederick2002time} on time preference estimates

\item \textbf{Reference Dependence}: \citet{koszegirab2006} on expectation-based
      reference points; \citet{knetschetal1990} on endowment effect

\item \textbf{Experimental Economics}: \citet{guth1982} on ultimatum games;
      \citet{bergsend1995} on trust; \citet{fehrgaechter2000} on public goods;
      \citet{camerer2018evaluating} on replication

\item \textbf{Cross-Cultural}: \citet{henrich2000,henrich2010} on WEIRD problem;
      \citet{herrmannthoeni2008} on antisocial punishment

\item \textbf{Dynamic Optimization}: Bellman (1957) on dynamic programming;
      Stokey \& Lucas (1989) on recursive methods

\item \textbf{Behavioral Change}: Prochaska \& DiClemente (1983) on stages of change;
      \citet{ajzen1991} on planned behavior

\item \textbf{Context Effects}: \citet{thaler1985} on mental accounting;
      Cialdini (2001) on social influence
\end{enumerate}

% =============================================================================
\section{Open Issues}
\label{sec:mp-open-issues}
% =============================================================================

\begin{enumerate}
\item \textbf{Computational Tractability}: The non-convex, high-dimensional
      problem may be computationally intractable for exact solution.
      Approximation methods needed.

\item \textbf{Parameter Uncertainty}: Many parameters ($\gamma$, $\alpha$, $\beta$)
      have wide confidence intervals. Robust optimization under uncertainty?

\item \textbf{Model Validation}: How to empirically test the complete problem
      specification? What would falsify it?

\item \textbf{Solution Space Formalization}: Appendix UNMAPPED_SS (CORE-SOLSPACE)
      needed to fully specify $\mathcal{S}(\Psi, S)$.

\item \textbf{Multi-Agent Extensions}: Current formulation assumes single
      decision-maker. How to extend to strategic interactions?

\item \textbf{Learning Dynamics}: How do agents learn about $\Psi$, $\gamma$
      over time? Bayesian updating? Reinforcement learning?
\end{enumerate}

% =============================================================================
\section{References}
\label{sec:mp-references}
% =============================================================================

This appendix draws on the complete EBF bibliography.

\nocite{bcm_master}

\textbf{Key References by Chapter Contribution:}

\begin{itemize}[nosep]
\item Ch1-4 (Foundations): See LIT-META, LIT-FEHR, LIT-THALER, LIT-KAHNEMAN
\item Ch5-7 (Complementarity): See Appendix UNMAPPED_B, NC, MIL (Milgrom-Roberts)
\item Ch9-15 (10C CORE): See individual CORE appendices (AAA, B, C, V, BBB, AU, AV, AW, HI)
\item Ch16 ($P_{eff}$): See Appendix UNMAPPED_PBB (FORMAL-PROBABILITY), UN (UNTCM)
\item Ch17-21 (Interventions): See Appendix UNMAPPED_IE, HHH, KKK
\item Ch22-24 (Scaling): See Appendix UNMAPPED_LLL, NNN, RRR, OOO, PPP
\end{itemize}

\textit{For complete bibliography, see the Master Bibliography (bcm\_master.bib)
with 1,922+ papers integrated into the EBF framework.}

% =============================================================================
% END OF APPENDIX MP
% =============================================================================
